% Livro sobre analise e engenharia de dados.
% Compile com: python compilar.py   (ver README.md)
\documentclass[11pt,twoside,openright,showtrims]{memoir}
\usepackage{estilo/livro}
% Dados do livro, num lugar so. A capa interna, a pagina de creditos, os
% cabecalhos e as propriedades do PDF leem daqui.

\newcommand{\titulolivro}{Do Cubo ao Copiloto}% titulo de trabalho, ver notas/02-titulos.md
\newcommand{\subtitulolivro}{Análise e engenharia de dados antes e depois da inteligência artificial}
\newcommand{\autorlivro}{Felipe Ramires Terrazas}
\newcommand{\edicaolivro}{1ª edição}
\newcommand{\anolivro}{2027}
\newcommand{\cidadelivro}{\pendente{cidade}}
\newcommand{\editoralivro}{Edição do autor}
\newcommand{\isbnlivro}{\pendente{ISBN}}

\addbibresource{referencias.bib}

\hypersetup{pdftitle={\titulolivro}, pdfauthor={\autorlivro}, pdfsubject={\subtitulolivro}}

\begin{document}

% ---------------------------------------------------------------- abertura
\frontmatter
% Paginas iniciais, na ordem tradicional de um livro impresso:
%   i   falsa folha de rosto (so o titulo)
%   ii  verso em branco
%   iii folha de rosto (titulo, subtitulo, autor, editora, cidade, ano)
%   iv  creditos (direitos, ficha catalografica, ISBN)
% Nenhuma delas leva numero de pagina impresso.

% i - falsa folha de rosto
\thispagestyle{empty}
\begin{vplace}[0.7]
  \centering
  {\LARGE\titulolivro\par}
\end{vplace}
\cleardoublepage

% iii - folha de rosto
\thispagestyle{empty}
\begin{center}
  \vspace*{0.12\textheight}
  {\large\textls[160]{\scshape\MakeLowercase{\autorlivro}}\par}
  \vspace{0.12\textheight}
  {\fontsize{30}{34}\selectfont\titulolivro\par}
  \vspace{1.2\onelineskip}
  {\color{fio}\rule{0.25\textwidth}{0.4pt}\par}
  \vspace{1.2\onelineskip}
  {\large\itshape\subtitulolivro\par}
  \vfill
  {\small\textls[120]{\scshape\MakeLowercase{\editoralivro}}\par}
  \vspace{0.3\onelineskip}
  {\small\cidadelivro, \anolivro\par}
\end{center}
\clearpage

% iv - creditos
\thispagestyle{empty}
{\footnotesize\raggedright
\vspace*{\fill}

Copyright \textcopyright\ \anolivro\ \autorlivro

Todos os direitos reservados. Nenhuma parte deste livro pode ser reproduzida,
por qualquer meio, sem autorização prévia e por escrito do autor, exceto trechos
curtos citados com a fonte.

\bigskip
\edicaolivro, \anolivro

\bigskip
\begin{tabular}{@{}ll@{}}
  Revisão & \pendente{revisor} \\
  Capa & \pendente{capa} \\
  Diagramação & o autor, em \LaTeX \\
\end{tabular}

\bigskip
\fbox{\parbox{0.9\textwidth}{\centering\vspace{0.5\onelineskip}
  Dados Internacionais de Catalogação na Publicação (CIP)\par
  \pendente{ficha catalográfica, feita por bibliotecário com CRB}\par
  \vspace{0.5\onelineskip}
  ISBN \isbnlivro\par\vspace{0.5\onelineskip}}}

\bigskip
Os nomes de produtos e empresas citados pertencem a seus donos e aparecem só
para identificar as ferramentas. As opiniões são do autor e não representam
nenhum empregador, atual ou passado.
\par}
\cleardoublepage

% Dedicatoria e epigrafe do livro (paginas impares, sem numero impresso).
% Apague este arquivo do main.tex se nao quiser nenhuma das duas.

\thispagestyle{empty}
\begin{vplace}[0.6]
  \raggedleft\itshape
  \pendente{dedicatória}
\end{vplace}
\cleardoublepage

\tableofcontents*
\chapter{Prefácio}
\label{cap:prefacio}

\rascunho{
  Por que escrevi este livro e para quem ele é. Ideias:
  \begin{itemize}
    \item A trajetória de trabalho com análise e engenharia de dados, antes e depois de a IA entrar no dia a dia.
    \item O incômodo que motivou o livro: a pergunta \enquote{ainda preciso aprender SQL e Python?} e a sensação de que as trilhas de estudo ficaram velhas.
    \item O que o livro é (um panorama geral, com a prática de quem está no meio do trabalho) e o que ele não é (manual de ferramenta, curso de programação).
    \item Agradecimentos curtos aqui ou numa seção própria no fim.
  \end{itemize}
}

\pendente{escrever o prefácio}

\chapter{Como ler este livro}
\label{cap:como-ler}

\rascunho{
  Mapa do livro e roteiros de leitura por perfil. Ideias:
  \begin{itemize}
    \item As cinco partes em um parágrafo cada (\nameref{parte:origens}, \nameref{parte:trilhas}, \nameref{parte:modelagem}, \nameref{parte:devops}, \nameref{parte:agil}).
    \item Roteiros: quem está começando lê tudo em ordem; quem já trabalha com dados pode ir direto à Parte~\ref{parte:modelagem}, \ref{parte:devops} ou \ref{parte:agil}; gestores podem ler a Parte~\ref{parte:origens} e a Parte~\ref{parte:agil}.
    \item Convenções do livro: as caixas \enquote{Da minha experiência} são relatos de projetos reais (sem nomes de empresas nem clientes); \enquote{Antes e depois da IA} compara o mesmo trabalho nos dois momentos; trechos de código em SQL e Python.
  \end{itemize}
}

\pendente{escrever o guia de leitura}


% ---------------------------------------------------------------- miolo
\mainmatter

\part{De onde viemos}
\label{parte:origens}
% Capitulo: Antes do dashboard
\chapter{Antes do dashboard}
\label{cap:origens}

Antes de existir dashboard, existia o relatório impresso. Ele chegava na mesa do
gerente alguns dias depois do fechamento do mês, em formulário contínuo, com páginas e
páginas de números alinhados em colunas. Quem precisava de uma resposta diferente
daquela que o relatório trazia fazia um pedido ao centro de processamento de dados e
esperava. Às vezes semanas.

Parece pré-história, mas boa parte dos problemas que discutimos hoje já estava ali:
números que chegam tarde, relatórios que ninguém lê, a distância entre quem tem a
pergunta e quem sabe extrair a resposta. Este capítulo volta a essa época por um motivo
prático. Quem entende de onde vêm as ideias do BI reconhece mais rápido quando uma
ferramenta nova está só dando outro nome a uma ideia velha, e quando está trazendo algo
de fato novo.

\section{Relatórios, planilhas e sistemas de apoio à decisão}
\label{sec:sad}

Nos anos 1960, as grandes empresas começaram a usar computadores de grande porte para
processar a folha de pagamento, o faturamento e o estoque. Como esses dados já estavam
na máquina, parecia natural usá-los também para informar a gestão. Surgiram os sistemas
de informação gerencial, que produziam relatórios periódicos a partir dos dados
operacionais.

A promessa era dar ao gestor toda a informação de que ele precisava. Em 1967, Russell
Ackoff publicou uma crítica que continua atual. Para ele, o problema da maioria dos
gestores não era a falta de informação relevante, e sim o excesso de informação
irrelevante \parencite{ackoff1967}. Os sistemas da época partiam do princípio de que mais
dados levariam a decisões melhores, e por isso despejavam relatórios. Ackoff argumentava
o contrário: um sistema de informação útil precisa filtrar e condensar, e para isso
precisa entender como a decisão é tomada.

Quem já viu uma empresa com duzentos dashboards e nenhuma reunião em que eles sejam
abertos reconhece o diagnóstico. A tecnologia mudou várias vezes desde então; a
tentação de produzir mais relatório em vez de melhor informação continua a mesma.

A resposta acadêmica a esse problema veio no começo dos anos 1970. \textcite{gorry1971}
separaram as decisões em estruturadas, semiestruturadas e não estruturadas, e propuseram
o nome \emph{sistemas de apoio à decisão} (SAD) para os sistemas voltados às decisões
que não se resolvem com uma regra fixa. A ideia era apoiar o julgamento do gestor, não
substituí-lo. \textcite{keen1978} desenvolveram o conceito no livro que se tornou
referência da área, e a história completa desse período está bem contada por
\textcite{power2007}.

Os SAD eram, em geral, sistemas construídos para uma decisão específica: planejar a
produção, avaliar uma carteira de investimentos, escolher rotas. Eram caros e dependiam
de especialistas. A virada que levou a análise para a mesa de qualquer pessoa veio de
outro lugar: as planilhas eletrônicas.

Quando as planilhas chegaram aos computadores pessoais, no fim dos anos 1970 e ao longo
dos anos 1980, o gestor passou a poder montar a própria análise sem pedir nada a
ninguém. Foi o primeiro \enquote{self-service} da história do BI, décadas antes de o
termo existir. E foi tão bem-sucedido que até hoje a planilha é a ferramenta de análise
mais usada nas empresas, inclusive naquelas que pagam caro por ferramentas de BI.

A planilha resolveu o problema da espera e criou outro, que vai aparecer várias vezes
neste livro: cada pessoa com a sua cópia dos dados, a sua fórmula e o seu número. O
capítulo~\ref{cap:master-data} trata de uma das consequências mais caras disso, os
cadastros que existem em versões diferentes em cada planilha.

No fim dos anos 1980, havia ainda os sistemas de informação executiva, que levavam
indicadores resumidos para a alta direção em telas simples. Eram, na prática, os avós
dos dashboards. A sopa de siglas (SIG, SAD, SIE) era grande, e foi nesse ambiente que
um nome guarda-chuva pegou.

\rascunho{
  Relato opcional: o primeiro contato com relatórios estáticos ou com uma empresa que
  rodava tudo em planilhas, e o que isso ensinou. Se não houver um caso bom, a seção se
  sustenta sem relato.
}

\section{De onde vem o termo business intelligence}
\label{sec:termo-bi}

A expressão \enquote{business intelligence} é mais antiga do que costuma se imaginar.
Ela já aparece em um livro de anedotas comerciais de 1865, para descrever um banqueiro
que conseguia vantagem sobre os concorrentes porque recebia notícias antes deles
\parencite{devens1865}. A palavra \emph{intelligence}, em inglês, tem esse sentido de
informação obtida para orientar uma ação, como em serviço de inteligência.

O uso que se aproxima do atual vem de 1958. Hans Peter Luhn, pesquisador da IBM,
publicou um artigo descrevendo um \enquote{sistema de business intelligence}: um sistema
automático que resumiria documentos, identificaria o interesse de cada área da empresa e
faria a informação certa chegar a quem precisava agir \parencite{luhn1958}. Luhn definia
inteligência como a capacidade de perceber as relações entre os fatos apresentados de
forma a orientar a ação em direção a um objetivo. \pendente{conferir a definição de Luhn
no artigo original antes de manter a paráfrase}

Vale reparar no que Luhn colocava no centro. Não era o armazenamento de dados nem o
gráfico; era a entrega da informação certa para a pessoa certa, com o objetivo de
orientar uma ação. Quase setenta anos depois, essa ainda é uma boa definição do que um
time de dados deveria fazer.

O termo só se popularizou no fim dos anos 1980. Howard Dresner, então analista do
Gartner, passou a usar \enquote{business intelligence} como guarda-chuva para os
conceitos e métodos de apoio à decisão baseados em fatos, num mercado cheio de siglas
\parencite{computerworld2006}. A autoria é disputada: o termo já aparecia em
publicações anteriores, e o próprio Dresner deu versões diferentes sobre quando e onde
começou a usá-lo. Para este livro, a disputa importa menos do que o fato de que o nome
pegou, e com ele veio uma indústria inteira.

\section{O que ficou dessa época}
\label{sec:legado-origens}

Os relatórios em formulário contínuo sumiram, os SAD viraram nota de rodapé e as siglas
foram esquecidas. Três ideias dessa época, porém, continuam valendo, e vão aparecer ao
longo do livro.

A primeira é que a pergunta de negócio vem antes do dado. Luhn falava em orientar a
ação; Gorry e Scott Morton classificavam sistemas pelo tipo de decisão que apoiavam.
Nenhum dos dois começava pela tecnologia. Projetos de dados que começam pela ferramenta,
e só depois procuram uma pergunta para responder, continuam fracassando pelo mesmo
motivo que fracassavam nos anos 1970.

A segunda é a crítica de Ackoff: mais informação não é melhor informação. Ela reaparece
quando o livro discute com que frequência os dados precisam atualizar
(capítulo~\ref{cap:frequencia}) e por que vale medir o uso dos próprios produtos de dados.
Um relatório que ninguém abre é o mesmo desperdício de 1967, só que mais barato de
produzir e, na era dos créditos e tokens, nem sempre tão barato de manter.

A terceira é que a análise sempre escapa para perto de quem decide. Aconteceu com a
planilha, aconteceu com as ferramentas de self-service (capítulo~\ref{cap:self-service})
e está acontecendo de novo com a IA (capítulo~\ref{cap:ia}). Cada vez que isso acontece,
a empresa ganha agilidade e perde consistência, e alguém precisa reconstruir a
consistência sem matar a agilidade. Esse alguém, quase sempre, é o time de dados.

O passo seguinte da história foi justamente uma tentativa de resolver a consistência de
uma vez: tirar os dados de análise de dentro dos sistemas operacionais e organizá-los num
lugar próprio. É o assunto do próximo capítulo.

\begin{leituras}
  \leitura{ackoff1967}{Um artigo curto e provocador. Os cinco pressupostos errados que
    Ackoff aponta nos sistemas de informação da época continuam aparecendo em projetos de
    dashboard.}
  \leitura{gorry1971}{O artigo que deu nome aos sistemas de apoio à decisão e que
    organizou as decisões gerenciais pelo grau de estrutura.}
  \leitura{power2007}{Um panorama acessível da história dos sistemas de apoio à decisão,
    escrito por quem acompanhou a área de perto.}
  \leitura{luhn1958}{Para ler o primeiro \enquote{sistema de business intelligence} na
    fonte e perceber como a preocupação com a pessoa certa e a ação certa já estava lá.}
\end{leituras}

% Capitulo: O data warehouse e a modelagem dimensional
\chapter{O data warehouse e a modelagem dimensional}
\label{cap:dw}

Nos anos 1990, o BI deixou de ser um conjunto de relatórios e planilhas e virou projeto
corporativo, com orçamento, equipe e arquitetura próprios. O centro desse projeto era o
data warehouse, o armazém de dados. Este capítulo conta como ele surgiu, apresenta as
duas escolas que definiram o campo, a de Bill Inmon e a de Ralph Kimball, e mostra por
que a modelagem dimensional, uma ideia dessa época, continua sendo uma das habilidades
mais úteis para quem trabalha com dados.

\section{Como surgiu o data warehouse}
\label{sec:surgimento-dw}

Os sistemas operacionais de uma empresa (vendas, estoque, contabilidade, folha de
pagamento) foram construídos para registrar transações. Eles precisam gravar um pedido
em milissegundos, garantir que dois vendedores não vendam a mesma unidade e manter o
saldo sempre correto. Para isso, guardam o estado atual das coisas, organizado de forma
a evitar repetição.

Análise pede quase o oposto. O gestor quer saber como as vendas evoluíram nos últimos
três anos, por região e por produto, comparando com o orçamento. Essa consulta lê
milhões de linhas, precisa de histórico que o sistema operacional muitas vezes já
apagou, e cruza dados de sistemas que não conversam entre si. Rodar esse tipo de
consulta no mesmo banco que registra as vendas deixa a operação lenta, e nem assim
resolve o cruzamento entre sistemas.

A saída que as empresas encontraram primeiro foi extrair dados. Cada área pedia um
arquivo do sistema de origem, tratava do seu jeito e montava o seu relatório. Com o
tempo, as extrações se multiplicavam, cada uma com uma regra um pouco diferente, e os
números de áreas diferentes deixavam de bater. \pendente{achar fonte para a
proliferação de extrações; Inmon descreve esse cenário}

Em 1988, Barry Devlin e Paul Murphy descreveram na IBM uma arquitetura diferente: um
armazém de dados de negócio, separado dos sistemas operacionais, alimentado de forma
controlada e usado como fonte única para os relatórios \parencite{devlin1988}. Quatro
anos depois, Bill Inmon publicou o livro que popularizou o termo data warehouse e
transformou a ideia numa disciplina \parencite{inmon1992}.

A definição de Inmon ficou clássica. Um data warehouse é uma coleção de dados orientada
por assunto, integrada, variável no tempo e não volátil, que apoia a tomada de decisão
gerencial \parencite{inmon2005}. \pendente{conferir a citação exata e a página na
edição usada} Cada um dos quatro adjetivos resolve um problema concreto:

\begin{description}[style=nextline, leftmargin=1.5em, font=\normalfont\itshape]
  \item[Orientada por assunto] Os dados se organizam pelos assuntos do negócio
    (cliente, produto, venda), não pelos sistemas de onde vieram. O analista procura
    \enquote{clientes} e não precisa saber se eles estão no ERP ou no CRM.
  \item[Integrada] O mesmo cliente tem o mesmo código e o mesmo nome em todo o armazém,
    mesmo que cada sistema de origem o registre de um jeito. É a primeira aparição, no
    livro, do problema de master data.
  \item[Variável no tempo] O armazém guarda o histórico. Se o cliente mudou de cidade, é
    possível saber onde ele estava quando fez cada compra.
  \item[Não volátil] O dado carregado não é alterado pelas transações do dia a dia. Uma
    consulta de hoje sobre o mês passado dá o mesmo resultado amanhã.
\end{description}

A Figura~\ref{fig:dw-linha-tempo} situa esses marcos numa linha do tempo que vai até a
IA. O resto do livro percorre, de certa forma, a metade de baixo dessa linha.

\begin{figure}[htbp]
  \centering
  % Linha do tempo do data warehouse (cap. 2). Fontes na legenda do capitulo.
\begin{tikzpicture}[diagrama, x=1mm, y=1mm]
  \draw[fio, line width=0.8pt] (18,5) -- (18,-85);
  \foreach \y/\ano/\texto/\peso in {
      0/1958/{Luhn descreve um \enquote{sistema de business intelligence} na IBM}/normal,
    -10/1970/{Codd propõe o modelo relacional, base dos bancos de dados SQL}/normal,
    -20/1988/{Devlin e Murphy descrevem a arquitetura de um \enquote{business data warehouse} na IBM}/normal,
    -30/1992/{Inmon publica \emph{Building the Data Warehouse}: o armazém corporativo, normalizado, construído de cima para baixo}/destaque,
    -40/1996/{Kimball publica \emph{The Data Warehouse Toolkit}: a modelagem dimensional, construída de baixo para cima}/destaque,
    -50/{Anos 2000}/{O Data Vault, de Dan Linstedt, ganha espaço como alternativa para a camada de integração}/normal,
    -60/{Anos 2010}/{Armazéns de dados na nuvem separam armazenamento de processamento, e o ETL dá lugar ao ELT}/normal,
    -70/2021/{Armbrust e colegas propõem o lakehouse, que junta data lake e data warehouse}/normal,
    -80/{Anos 2020}/{A IA generativa entra no fluxo de trabalho com dados}/normal%
  }{
    \ifthenelse{\equal{\peso}{destaque}}
      {\filldraw[draw=tinta, fill=tinta] (18,\y) circle (1.4);}
      {\filldraw[draw=cinza, fill=white, line width=0.6pt] (18,\y) circle (1.1);}
    \node[anchor=east, font=\footnotesize] at (14.5,\y) {\ano};
    \node[anchor=west, text width=92mm, align=left, font=\footnotesize] at (22,\y) {\texto};
  }
\end{tikzpicture}

  \caption{Uma linha do tempo do data warehouse.}
  \label{fig:dw-linha-tempo}
  \fonte{elaborado pelo autor com base em \textcite{luhn1958}, \textcite{codd1970},
    \textcite{devlin1988}, \textcite{inmon1992}, \textcite{kimball1996},
    \textcite{linstedt2015} e \textcite{armbrust2021}.}
\end{figure}

\section{A abordagem de Inmon}
\label{sec:inmon}

Para Inmon, o data warehouse é a base corporativa de informação. Ele deve ser
construído de cima para baixo: primeiro se modela a empresa como um todo, depois se
carrega um armazém central, integrado e normalizado, na terceira forma normal, o mesmo
tipo de modelagem usado nos sistemas transacionais para evitar repetição de dados
\parencite{inmon2005}. Esse armazém é a fonte de verdade.

Os data marts, bases menores voltadas a uma área (vendas, finanças, logística), vêm
depois e são derivados do armazém central. Por isso são chamados de dependentes. Neles,
sim, os dados podem ser reorganizados para facilitar a consulta. A
Figura~\ref{fig:inmon} resume o desenho.

\begin{figure}[htbp]
  \centering
  % Arquitetura de Inmon: de cima para baixo (cap. 2).
\begin{tikzpicture}[diagrama, x=1mm, y=1mm]
  % sistemas de origem
  \foreach \y/\nome/\id in {13.5/ERP/f1, 4.5/CRM/f2, -4.5/Planilhas/f3, -13.5/{Outros}/f4}
    \node[caixa, minimum width=15mm, minimum height=6mm] (\id) at (8,\y) {\nome};
  \node[moldura, fit=(f1)(f4), label={[rotulo]above:Sistemas de origem}] {};

  % armazem corporativo
  \node[bancodestaque, minimum width=26mm, minimum height=26mm] (edw) at (46,0)
    {Data warehouse\\corporativo\\[1mm]{\scriptsize normalizado (3FN)}\\{\scriptsize e integrado}};
  \foreach \id in {f1,f2,f3,f4} \draw[seta] (\id.east) -- (edw.west);
  \node[rotulo] at (25.5,18) {ETL};

  % data marts dependentes
  \foreach \y/\nome/\id in {13/Vendas/m1, 0/Finanças/m2, -13/Logística/m3}
    \node[banco, minimum width=17mm, minimum height=9mm] (\id) at (79,\y) {\nome};
  \node[moldura, fit=(m1)(m3), label={[rotulo]above:Data marts dependentes}] {};
  \foreach \id in {m1,m2,m3} \draw[seta] (edw.east) -- (\id.west);

  % consumo
  \foreach \y/\nome/\id in {13/Relatórios/u1, 0/Dashboards/u2, -13/Análises/u3}
    \node[caixa, minimum width=16mm, minimum height=6mm] (\id) at (108,\y) {\nome};
  \foreach \i in {1,2,3} \draw[seta] (m\i.east) -- (u\i.west);

  \node[rotulo, text width=110mm] at (58,-27)
    {Ordem de construção: primeiro o armazém corporativo, depois os data marts de cada área.};
\end{tikzpicture}

  \caption{A arquitetura de Inmon: armazém corporativo normalizado e data marts dependentes.}
  \label{fig:inmon}
  \fonte{elaborado pelo autor com base em \textcite{inmon2005} e \textcite{inmon2001}.}
\end{figure}

Inmon e coautores ampliaram depois essa visão para a Corporate Information Factory, uma
arquitetura completa em que o data warehouse é uma das peças, ao lado de bases
operacionais integradas, data marts e aplicações de análise \parencite{inmon2001}.

A força dessa abordagem é a consistência. Como tudo passa pelo armazém central, as áreas
partem dos mesmos dados e das mesmas regras. E como o núcleo é normalizado, ele aguenta
bem perguntas que ninguém previu quando o modelo foi desenhado. O ponto fraco é o tempo.
Modelar a empresa inteira antes de entregar o primeiro relatório exige meses, às vezes
anos, e um patrocínio que nem sempre sobrevive à espera.

\section{A abordagem de Kimball}
\label{sec:kimball}

Ralph Kimball propôs o caminho inverso. Em vez de começar pela empresa inteira, o
armazém é construído um processo de negócio por vez: primeiro vendas, depois estoque,
depois compras. Cada processo ganha um modelo dimensional, desenhado para ser fácil de
entender e rápido de consultar \parencite{kimball1996,kimball2013}.

O que impede esse crescimento incremental de virar um amontoado de bases isoladas são as
dimensões conformadas. Data, produto, cliente e loja são as mesmas tabelas, com as mesmas
chaves e os mesmos atributos, em todos os modelos. Assim, uma análise que cruza vendas e
estoque por produto funciona, porque os dois usam exatamente a mesma dimensão de produto.
O conjunto de modelos ligados pelas dimensões conformadas forma o que Kimball chama de
arquitetura de barramento (Figura~\ref{fig:kimball}). Uma dimensão conformada é, na
prática, master data com outro nome, e o capítulo~\ref{cap:master-data} volta a esse
ponto.

\begin{figure}[htbp]
  \centering
  % Arquitetura de Kimball: data marts dimensionais integrados pelo barramento (cap. 2).
\begin{tikzpicture}[diagrama, x=1mm, y=1mm]
  % sistemas de origem
  \foreach \y/\nome/\id in {13.5/ERP/f1, 4.5/CRM/f2, -4.5/Planilhas/f3, -13.5/{Outros}/f4}
    \node[caixa, minimum width=15mm, minimum height=6mm] (\id) at (8,\y) {\nome};
  \node[moldura, fit=(f1)(f4), label={[rotulo]above:Sistemas de origem}] {};

  % area de preparacao
  \node[caixa, minimum width=19mm, minimum height=30mm] (stg) at (36,0)
    {Área de\\preparação\\[1mm]{\scriptsize (ETL)}};
  \foreach \id in {f1,f2,f3,f4} \draw[seta] (\id.east) -- (\id.east -| stg.west);

  % area de apresentacao: data marts dimensionais
  \foreach \x/\nome/\id in {62/Vendas/m1, 80/Estoque/m2, 98/Finanças/m3}
    \node[bancodestaque, minimum width=16mm, minimum height=13mm] (\id) at (\x,5)
      {\nome\\{\scriptsize estrela}};
  % barramento de dimensoes conformadas
  \node[caixa, fill=fundo, draw=cinza, minimum width=56mm, minimum height=6mm, inner sep=1pt] (bus) at (80,-12)
    {\scriptsize Barramento: dimensões conformadas (Data, Produto, Cliente)};
  \foreach \id in {m1,m2,m3} \draw[ligacao, line width=1.2pt] (\id.south) -- (\id.south |- bus.north);
  \node[moldura, fit=(m1)(m3)(bus), label={[rotulo]above:Área de apresentação}] (apr) {};
  \draw[seta] (stg.east) -- (stg.east -| apr.west);

  \node[rotulo, text width=110mm] at (58,-27)
    {Ordem de construção: um processo de negócio por vez, integrados pelas dimensões conformadas.};
\end{tikzpicture}

  \caption{A arquitetura de Kimball: data marts dimensionais integrados pelo barramento de dimensões conformadas.}
  \label{fig:kimball}
  \fonte{elaborado pelo autor com base em \textcite{kimball2013}.}
\end{figure}

O planejamento desse barramento cabe numa tabela simples, a matriz de barramento: os
processos de negócio nas linhas, as dimensões nas colunas, e uma marca onde o processo
usa a dimensão (Tabela~\ref{tab:matriz-barramento}). Ela mostra, de relance, quais
dimensões precisam ser construídas com mais cuidado, porque serão compartilhadas por
muitos processos.

\begin{table}[htbp]
  \centering
  \caption{Uma matriz de barramento: processos de negócio nas linhas, dimensões conformadas nas colunas.}
  \label{tab:matriz-barramento}
  \small
  \begin{tabular}{@{}lccccc@{}}
    \toprule
    Processo de negócio & Data & Produto & Cliente & Loja & Fornecedor \\
    \midrule
    Vendas      & $\bullet$ & $\bullet$ & $\bullet$ & $\bullet$ &           \\
    Compras     & $\bullet$ & $\bullet$ &           & $\bullet$ & $\bullet$ \\
    Estoque     & $\bullet$ & $\bullet$ &           & $\bullet$ &           \\
    Devoluções  & $\bullet$ & $\bullet$ & $\bullet$ & $\bullet$ &           \\
    \bottomrule
  \end{tabular}
  \fonte{exemplo do autor, no formato da matriz de barramento de \textcite{kimball2013}.}
\end{table}

Duas outras ideias de Kimball aparecem em quase todo projeto de dados até hoje. As
chaves substitutas (\emph{surrogate keys}) são códigos gerados pelo próprio armazém para
cada linha de dimensão, independentes dos códigos dos sistemas de origem. Elas protegem o
modelo quando a origem muda a numeração ou quando dois sistemas usam o mesmo código para
coisas diferentes. As dimensões que mudam devagar (\emph{slowly changing dimensions})
são as técnicas para lidar com atributos que mudam com o tempo, como a cidade de um
cliente. A forma mais simples sobrescreve o valor antigo; a mais usada para análise
histórica cria uma nova linha a cada mudança, com datas de início e fim de validade.
Essa segunda forma vai reaparecer, com outro nome, quando o livro tratar de dados para
modelos preditivos (\veja{sec:vazamento}).

A força da abordagem de Kimball é a entrega rápida e o modelo que o usuário entende. O
ponto fraco é que tudo depende da disciplina nas dimensões conformadas. Sem ela, cada
processo ganha a sua própria dimensão de cliente, e o armazém vira um conjunto de silos
com a aparência de um sistema integrado.

\section{Fatos, dimensões e granularidade}
\label{sec:modelagem-dimensional}

O modelo dimensional tem duas peças. As tabelas de fatos guardam as medidas de um
processo de negócio: quantidade vendida, valor, custo. As tabelas de dimensão guardam o
contexto dessas medidas: quando, o quê, quem, onde. Desenhado, o modelo parece uma
estrela, com o fato no centro e as dimensões em volta (Figura~\ref{fig:estrela}).

\begin{figure}[htbp]
  \centering
  % Esquema estrela de vendas (cap. 2). Os mesmos nomes da consulta SQL do capitulo.
\begin{tikzpicture}[diagrama, x=1mm, y=1mm]
  \node[tabela, line width=0.9pt] (f) at (0,0)
    {\textbf{fato\_vendas}\nodepart{second}id\_data\\id\_produto\\id\_cliente\\id\_loja\\[0.6mm]quantidade\\valor\_total};

  \node[tabela] (d1) at (-40,20)
    {\textbf{dim\_data}\nodepart{second}id\_data\\data\\mes\\trimestre\\ano};
  \node[tabela] (d2) at (40,20)
    {\textbf{dim\_produto}\nodepart{second}id\_produto\\nome\\categoria\\marca};
  \node[tabela] (d3) at (-40,-20)
    {\textbf{dim\_cliente}\nodepart{second}id\_cliente\\nome\\cidade\\segmento};
  \node[tabela] (d4) at (40,-20)
    {\textbf{dim\_loja}\nodepart{second}id\_loja\\nome\\cidade\\regiao};

  \foreach \d in {d1,d2,d3,d4} \draw[ligacao, line width=0.7pt] (f) -- (\d);
  \node[rotulo, anchor=north] at (f.south) [yshift=-1mm] {grão: uma linha por item vendido};
\end{tikzpicture}

  \caption{Um esquema estrela de vendas. A tabela de fatos, no centro, guarda as chaves das dimensões e as medidas (quantidade e valor); as dimensões, em volta, guardam o contexto.}
  \label{fig:estrela}
  \fonte{elaborado pelo autor.}
\end{figure}

\textcite{kimball2013} organizam o desenho em quatro passos: escolher o processo de
negócio, declarar o grão, identificar as dimensões e identificar os fatos. A ordem
importa, e o segundo passo é o mais importante e o mais esquecido.

O grão é o que uma linha da tabela de fatos representa. \enquote{Uma linha por item
vendido} é um grão; \enquote{uma linha por pedido} é outro; \enquote{uma linha por loja
por dia} é um terceiro. Cada um responde a perguntas diferentes, e misturar grãos na
mesma tabela é uma das origens mais comuns de número errado em dashboard. Se parte das
linhas é por item e parte é por pedido, qualquer soma conta alguma coisa duas vezes.
Kimball recomenda declarar o grão mais detalhado que os dados permitem, porque agregar
depois é fácil e desagregar é impossível.

\rascunho{
  Relato: um erro de grão visto ou cometido na prática (uma tabela que misturava
  níveis de detalhe e gerava um número errado que alguém levou para uma reunião).
  Nomes e valores só em ordem de grandeza.
}

A recompensa por esse cuidado aparece na hora de consultar. Com o grão claro e as
dimensões bem desenhadas, as perguntas de negócio viram consultas curtas e previsíveis:
junta o fato às dimensões de que a pergunta precisa, filtra, agrupa e soma.

\begin{sql}
-- Receita por mês e por categoria de produto, no esquema estrela de vendas
SELECT
    d.ano,
    d.mes,
    p.categoria,
    SUM(f.valor_total) AS receita
FROM fato_vendas  AS f
JOIN dim_data     AS d ON d.id_data    = f.id_data
JOIN dim_produto  AS p ON p.id_produto = f.id_produto
GROUP BY d.ano, d.mes, p.categoria
ORDER BY d.ano, d.mes, p.categoria;
\end{sql}

Essa previsibilidade também ajuda a IA. Quando o livro tratar de dados como contexto
para modelos de linguagem (capítulo~\ref{cap:contexto-ia}), o modelo dimensional bem
feito vai aparecer como meio caminho andado: grão explícito e nomes claros reduzem a
chance de a máquina somar a coisa errada.

\section{Inmon ou Kimball?}
\label{sec:inmon-ou-kimball}

Durante anos, a pergunta dividiu a área. A Tabela~\ref{tab:inmon-kimball} resume as
diferenças principais.

\begin{table}[htbp]
  \centering
  \caption{Inmon e Kimball lado a lado.}
  \label{tab:inmon-kimball}
  \small
  \begin{tabularx}{\textwidth}{@{}>{\raggedright\arraybackslash}p{0.2\textwidth}
      >{\raggedright\arraybackslash}X>{\raggedright\arraybackslash}X@{}}
    \toprule
    & Inmon & Kimball \\
    \midrule
    Ponto de partida & O armazém corporativo, de cima para baixo & Os processos de negócio, de baixo para cima \\
    \addlinespace
    Modelo do núcleo & Normalizado (3FN) & Dimensional (esquema estrela) \\
    \addlinespace
    Data marts & Dependentes, derivados do armazém & São a própria camada de apresentação \\
    \addlinespace
    Integração & No armazém central & Pelas dimensões conformadas (barramento) \\
    \addlinespace
    Primeira entrega & Mais lenta, pede o modelo corporativo antes & Mais rápida, um processo por vez \\
    \addlinespace
    Prioridade & Consistência e integração na empresa toda & Facilidade de consulta para o usuário \\
    \bottomrule
  \end{tabularx}
  \fonte{elaborado pelo autor com base em \textcite{inmon2005} e \textcite{kimball2013}.}
\end{table}

Na prática, muitas empresas acabaram com uma mistura das duas: uma camada integrada
perto das fontes, que lembra o armazém de Inmon, e modelos dimensionais para consumo,
no estilo de Kimball. Nos anos 2000 surgiu ainda uma terceira opção para a camada de
integração, o Data Vault, que separa as chaves de negócio, os relacionamentos e os
atributos descritivos em estruturas próprias, pensando em fontes que mudam com
frequência \parencite{linstedt2015}.

A nuvem mudou a conta de novo. Com armazenamento barato e processamento que cresce sob
demanda (capítulo~\ref{cap:nuvem}), algumas preocupações de desempenho que pesavam nas
decisões de Inmon e Kimball perderam força. É comum ver tabelas largas, que repetem
atributos que um modelo normalizado evitaria. Mesmo assim, a modelagem dimensional
continua no centro dos projetos de analytics engineering, porque o problema que ela
resolve não é de desempenho: é fazer o dado ser entendido do mesmo jeito por todo mundo
que o consulta.

\rascunho{
  Opinião do autor, com base no que viu funcionar: qual das abordagens aparece na
  prática dos projetos de hoje e por quê. Pode fechar com um relato curto.
}

\section{ETL, cargas noturnas e cubos OLAP}
\label{sec:etl-olap}

Os armazéns de dados dos anos 1990 e 2000 eram abastecidos por processos de ETL:
extrair os dados das fontes, transformá-los fora do banco de destino, num servidor
próprio, e carregá-los já prontos. Como as fontes eram sistemas operacionais que não
podiam ficar lentos durante o expediente, a carga acontecia de madrugada, numa janela
que precisava terminar antes de a empresa abrir. Quando a janela estourava, os
relatórios da manhã chegavam com o dado do dia anterior ao anterior.

Para que a consulta fosse rápida, os dados agregados eram pré-calculados em cubos
multidimensionais: estruturas em que cada combinação de dimensões (mês, região,
produto) já tinha o total guardado. O termo OLAP, de processamento analítico em linha,
foi proposto em 1993 num documento assinado por Edgar Codd, o criador do modelo
relacional, com doze regras para esse tipo de ferramenta \parencite{codd1993}. O
documento foi encomendado por um fornecedor de software de cubos, o que gerou
críticas sobre a independência das regras. O termo, ainda assim, pegou.

Em volta dos armazéns e dos cubos cresceram as suítes corporativas de BI, que
juntavam relatórios, cubos e painéis num mesmo pacote, com licenças caras e equipes
especializadas. No fim dos anos 2000, os maiores fornecedores independentes dessas
suítes foram comprados por grandes empresas de software corporativo.
\pendente{conferir e citar as aquisições de 2007 e 2008}

Desse período, o que envelheceu foi a infraestrutura: janelas noturnas, cubos
pré-calculados, servidores dedicados. O que ficou foi o modelo. Quando o processamento
ficou barato o bastante para calcular os totais na hora, os cubos perderam a razão de
existir, mas os fatos, as dimensões e o grão continuaram exatamente onde estavam. O
próximo capítulo mostra o que aconteceu quando a ferramenta de análise saiu do
departamento de TI e foi parar na mesa do analista.

\begin{leituras}
  \leitura{kimball2013}{A referência da modelagem dimensional. Os primeiros capítulos
    apresentam o processo de quatro passos; os demais aplicam a técnica a casos de
    vários setores.}
  \leitura{inmon2005}{A visão de Inmon na fonte: a definição de data warehouse, o modelo
    corporativo normalizado e o cuidado com a granularidade dos dados.}
  \leitura{inmon2001}{A Corporate Information Factory, a arquitetura completa em que o
    data warehouse é uma das peças.}
  \leitura{linstedt2015}{Para conhecer o Data Vault, alternativa para a camada de
    integração quando as fontes mudam com frequência.}
  \leitura{devlin1988}{Artigo curto e histórico, um dos primeiros a descrever a
    arquitetura de um armazém de dados separado dos sistemas operacionais.}
\end{leituras}

% Capitulo: Self-service e a era dos dashboards
\chapter{Self-service e a era dos dashboards}
\label{cap:self-service}

O data warehouse resolveu a consistência, mas manteve a distância. Para ter um relatório
novo, o analista de negócio ainda abria um chamado, esperava a fila da TI e recebia, semanas
depois, algo parecido com o que tinha pedido. Nos anos 2000 e 2010, uma nova geração de
ferramentas encurtou essa distância de forma radical: o próprio analista passou a
conectar os dados, arrastar campos e montar o seu painel. Este capítulo trata dessa
virada, dos problemas que ela criou e do que vale aprender sobre visualização
independentemente da ferramenta.

\section{O BI sai da TI}
\label{sec:bi-sai-da-ti}

Uma das origens dessa mudança é um projeto acadêmico. Na Universidade Stanford, Chris
Stolte, Diane Tang e Pat Hanrahan desenvolveram o Polaris, um sistema que traduzia
gestos visuais (arrastar uma dimensão para as linhas, uma medida para as colunas) em
consultas ao banco de dados e mostrava o resultado como gráfico
\parencite{stolte2002}. A ideia central era que explorar dados devia ser um processo
visual e iterativo, e não a redação de uma consulta seguida da espera pelo resultado. O
projeto deu origem a uma das ferramentas de visualização mais conhecidas do mercado.

Na mesma época, outras ferramentas apostaram em carregar os dados na memória do
computador e permitir que o usuário navegasse livremente por eles, clicando num valor e
vendo todo o resto se reorganizar em torno dele. E, já na metade dos anos 2010, os
grandes fornecedores de software corporativo entraram com ferramentas próprias, muitas
vezes integradas à planilha que o analista já usava todo dia, a preços que tornaram a
licença acessível para qualquer área.

O efeito somado foi uma mudança de quem fazia BI. O analista de negócio, que antes
consumia relatórios, passou a produzi-los. A TI deixou de ser a fábrica de relatórios e
passou, em tese, a cuidar da infraestrutura e dos dados de base. O tempo entre a
pergunta e o primeiro gráfico caiu de semanas para horas.

Foi uma conquista real. Muita gente que hoje trabalha com dados entrou na área por esse
caminho: começou montando painéis para a própria equipe, descobriu que gostava e foi
aprendendo SQL, modelagem e estatística no meio do caminho. A trilha tradicional de
aprendizado que o capítulo~\ref{cap:trilha-tradicional} descreve tem muito desse
período.

\section{Mil dashboards e nenhuma verdade}
\label{sec:mil-dashboards}

A liberdade trouxe de volta o problema que o data warehouse tinha tentado resolver. Cada
área passou a conectar a ferramenta onde achava melhor: no armazém, numa extração, numa
planilha que alguém mantinha à mão. Cada uma aplicava os próprios filtros e as próprias
regras. Não demorou para aparecerem três versões do faturamento do mês, todas
\enquote{oficiais}, todas diferentes.

O sintoma clássico é a reunião em que a primeira meia hora se gasta discutindo de onde
veio cada número, e a decisão fica para a próxima semana. O problema raramente é a
ferramenta. É que a definição da métrica (o que entra no faturamento, se devolução
desconta, em que data a venda conta) foi feita várias vezes, por pessoas diferentes, sem
que ninguém soubesse.

Há também o problema do volume. Como criar um painel ficou fácil, criar mais um painel
passou a ser a resposta para qualquer pergunta. Empresas acumularam centenas de painéis,
muitos abertos uma vez e esquecidos, outros duplicados com pequenas variações, todos
consumindo atualização, manutenção e, cada vez mais, dinheiro. É a crítica de Ackoff
(\veja{sec:sad}) com ferramenta moderna. O capítulo~\ref{cap:frequencia} volta a esse
ponto com uma proposta concreta: medir o uso dos próprios produtos de dados e desligar o
que não é usado.

\rascunho{
  Relato: uma situação de números divergentes entre áreas, ou de painéis duplicados, e
  como foi resolvida. Sem nomes de empresa nem valores exatos.
}

\section{Governança sem burocracia}
\label{sec:governanca}

A primeira reação de muitas empresas ao caos foi recentralizar: tirar a liberdade das
áreas e voltar a fila da TI. Funciona por algum tempo, até as áreas voltarem às
planilhas, porque a fila continua lenta e a pergunta continua urgente. Governança que só
proíbe empurra o trabalho para fora do alcance dela.

O que costuma funcionar é separar o que precisa ser único do que pode ser livre. O que
precisa ser único são os dados de base e as definições das métricas que a empresa usa
para decidir: faturamento, margem, cliente ativo. Esses ficam num lugar só, com dono e
com regra escrita, e as ferramentas de self-service consomem dali. O que pode ser livre é
a exploração em cima dessa base: cada área monta as análises e os painéis de que
precisa, desde que parta dos dados certificados.

Na prática, isso se traduz em conjuntos de dados certificados, que as ferramentas
marcam como oficiais, e numa camada semântica, onde cada métrica é definida uma vez e
reaproveitada por todos os painéis (\veja{sec:camada-semantica}). Aqui o livro só
apresenta o problema. A governança de dados como disciplina, e a forma como a IA muda o
papel de quem cuida dela, ficam para o capítulo~\ref{cap:governanca}.

\section{Visualização que não depende de ferramenta}
\label{sec:visualizacao}

As ferramentas de self-service tornaram fácil fazer gráfico. Não tornaram fácil fazer
gráfico bom. Algumas ideias, todas anteriores às ferramentas atuais, continuam sendo o
melhor ponto de partida.

Edward Tufte propôs que a maior parte da tinta de um gráfico deveria representar dados,
e não decoração: bordas, sombras, fundos e efeitos tridimensionais disputam a atenção
com a informação \parencite{tufte1983}. O princípio vale para a tela. Um painel limpo,
com poucos elementos e muito espaço, comunica mais do que um painel cheio de enfeites.

Stephen Few levou essas ideias para os dashboards e propôs uma definição exigente: um
dashboard mostra, numa única tela, a informação mais importante para atingir um ou mais
objetivos, de forma que possa ser acompanhada de relance \parencite{few2006}. Se o painel
precisa de rolagem, de dez abas e de um manual, ele pode ser útil, mas não é um
dashboard nesse sentido. Few também insistia em comparar sempre: um número sozinho não
diz se está bom ou ruim; o mesmo número ao lado da meta, do mês anterior ou do mesmo mês
do ano passado diz.

Cole Nussbaumer Knaflic acrescentou a perspectiva da comunicação: antes de escolher o
gráfico, é preciso saber quem vai ler, o que essa pessoa precisa decidir e qual é a
mensagem principal \parencite{knaflic2015}. Na maior parte das vezes, uma barra simples,
com o ponto importante destacado e um título que diz a conclusão, funciona melhor do que
qualquer gráfico sofisticado.

\begin{dica}
Antes de publicar um painel, faça três perguntas: quem vai abrir, com que frequência, e
que decisão essa pessoa toma com o que está vendo. Se não houver resposta para a terceira,
o painel provavelmente não precisa existir.
\end{dica}

Essas três referências têm em comum o fato de que nenhuma depende de ferramenta. É
por isso que valem mais do que qualquer curso de um produto específico: o produto muda a
cada poucos anos, o olho de quem lê o gráfico não.

O próximo capítulo volta para os bastidores. Enquanto as ferramentas de visualização
democratizavam a análise, o volume de dados crescia, o armazenamento ficava barato e o
data warehouse começava a se mudar para a nuvem.

\begin{leituras}
  \leitura{few2006}{O livro que definiu o que é um bom dashboard. Os exemplos de painéis
    ruins, analisados um a um, ensinam tanto quanto as recomendações.}
  \leitura{knaflic2015}{Para quem precisa apresentar dados a gestores: como escolher o
    gráfico, tirar o excesso e conduzir a atenção para a mensagem.}
  \leitura{tufte1983}{O clássico da visualização de dados. Mais conceitual que prático,
    mas forma o olhar para sempre.}
  \leitura{stolte2002}{Para ver, na fonte, a ideia de exploração visual que está por trás
    de boa parte das ferramentas de self-service.}
\end{leituras}

% Capitulo: Nuvem, big data e o modern data stack
\chapter{Nuvem, big data e o modern data stack}
\label{cap:nuvem}

Enquanto o self-service levava a análise para a mesa do analista, os bastidores
mudavam de escala. A internet, os celulares e os sensores passaram a gerar dados num
volume que os armazéns tradicionais não davam conta de guardar, e muito menos de
processar dentro de uma janela noturna. Este capítulo percorre as três ondas que vieram
em seguida: o big data, os armazéns de dados na nuvem e a troca do ETL pelo ELT, que deu
origem ao chamado modern data stack. E termina com uma constatação que vai acompanhar o
resto do livro: muitas das ideias novas são ideias antigas com outra roupa.

\section{Big data e data lakes}
\label{sec:big-data}

Em 2001, Doug Laney, analista do META Group, descreveu os desafios de gestão de dados em
três dimensões: volume, velocidade e variedade \parencite{laney2001}. Os \enquote{três
Vs} viraram, anos depois, a definição informal de big data: dados grandes demais, rápidos
demais ou heterogêneos demais para as ferramentas da época.

A resposta técnica veio das empresas de internet, que enfrentavam o problema antes de
todo mundo. O Google publicou em 2003 a arquitetura do seu sistema de arquivos
distribuído, que espalhava dados por milhares de máquinas baratas
\parencite{ghemawat2003}, e em 2004 o MapReduce, um modelo para processar esses dados
em paralelo dividindo o trabalho em duas etapas simples \parencite{dean2004}. Um projeto
de código aberto inspirado nesses dois artigos, o Hadoop, ganhou força a partir de 2006 e
se tornou sinônimo de big data na década seguinte.

A promessa do Hadoop era guardar tudo, em qualquer formato, a um custo baixo, e decidir
depois o que fazer com os dados. Em 2010, James Dixon cunhou o termo \emph{data lake}
para essa ideia: em vez de um armazém em que a água chega engarrafada e pronta para
consumo, um lago onde ela chega no estado natural e cada um retira o que precisa
\parencite{dixon2010}.

A metáfora era boa, e a prática nem sempre acompanhou. Muitos lagos viraram pântanos:
arquivos sem documentação, sem dono, sem garantia de qualidade, que ninguém tinha
coragem de apagar e ninguém conseguia usar. E a tecnologia exigia especialistas raros,
escrevendo código de baixo nível para tarefas que, num banco relacional, seriam uma
consulta SQL de poucas linhas. Não por acaso, uma das primeiras coisas que o ecossistema
Hadoop ganhou foram ferramentas para consultar os dados com SQL.

O big data deixou duas heranças duradouras. A primeira é a ideia de separar o
armazenamento barato e flexível do processamento, que vai reaparecer na nuvem e no
lakehouse. A segunda é a lição de que guardar tudo não substitui modelar e documentar. O
dado sem contexto, por mais barato que seja guardá-lo, continua sem valor.

\section{O armazém vai para a nuvem}
\label{sec:dw-nuvem}

A segunda onda trouxe o data warehouse de volta, agora como serviço. Em vez de comprar
servidores e licenças, a empresa passou a alugar capacidade de um provedor de nuvem e a
pagar pelo uso.

Três sistemas marcaram essa transição, e todos foram descritos em artigos técnicos que
valem a leitura. O Dremel, do Google, mostrou que era possível consultar bilhões de
linhas em segundos usando armazenamento em colunas e processamento distribuído
\parencite{melnik2010}; ele está por trás de um dos primeiros armazéns oferecidos como
serviço. A Amazon defendeu um armazém de dados mais simples de contratar e operar, que
qualquer time pudesse ligar em minutos \parencite{gupta2015}. E o Snowflake levou ao
extremo a separação entre armazenamento e processamento: os dados ficam num lugar, e
quantos \enquote{motores} de processamento forem necessários podem ser ligados e
desligados sobre eles, de forma independente \parencite{dageville2016}.

Essa separação mudou a economia do armazém de dados. No modelo antigo, a capacidade era
comprada para o pico e ficava ociosa a maior parte do tempo. No modelo novo, o custo
acompanha o uso: uma consulta pesada custa mais, um processo que roda a cada hora custa
mais do que um que roda uma vez por dia. É uma vantagem enorme para quem começa pequeno,
e uma armadilha para quem não acompanha o consumo. Essa mudança, da licença fixa para o
pagamento por uso, foi o primeiro passo de uma tendência que a IA acelerou e que o
capítulo~\ref{cap:custo-consumo} trata em detalhe.

\section{De ETL para ELT}
\label{sec:elt}

Com processamento barato e elástico dentro do próprio armazém, uma inversão passou a
fazer sentido. Em vez de transformar os dados num servidor intermediário antes de
carregá-los (ETL), a empresa carrega os dados brutos no armazém e transforma lá dentro,
com SQL (ELT). A janela noturna deixa de ser um gargalo, o dado bruto fica disponível
para reprocessar quando uma regra muda, e a transformação passa a ser escrita na
linguagem que os analistas já conhecem.

Em meados dos anos 2010 surgiram ferramentas para organizar essa transformação como
código: modelos em SQL, versionados, testados e documentados, com dependências entre si
resolvidas automaticamente. A mais conhecida delas, o dbt, ajudou a consolidar uma forma
de trabalhar que trouxe práticas de engenharia de software para dentro do time de dados.
A Parte~\ref{parte:devops} do livro é, em boa medida, sobre essas práticas.

Com a nova forma de trabalhar veio um papel novo. Em 2019, Michael Kaminsky descreveu o
\emph{analytics engineer}: alguém entre o engenheiro de dados e o analista, que domina SQL
e modelagem, aplica práticas de engenharia de software e entrega conjuntos de dados
limpos e bem documentados para os analistas usarem \parencite{kaminsky2019}. A
comunidade em torno do dbt adotou e popularizou o termo. O
capítulo~\ref{cap:mapas} compara esse papel com o do analista e o do engenheiro de
dados.

Os projetos de analytics engineering costumam se organizar em camadas: uma camada de
preparação, que limpa e padroniza o dado bruto de cada fonte; uma camada intermediária,
que junta e aplica regras de negócio quando necessário; e uma camada final, com os
modelos prontos para consumo, muitas vezes dimensionais. Quem leu o capítulo anterior
reconhece aqui ecos de Inmon e de Kimball. O capítulo~\ref{cap:medalhao} mostra como
essa organização em camadas ganhou um nome próprio, a arquitetura medalhão.

Ao conjunto de ferramentas dessa época (ingestão gerenciada, armazém na nuvem,
transformação como código, BI por cima) deu-se o nome de \emph{modern data stack}. O
nome é de marketing, mas a mudança foi real: um time pequeno passou a conseguir montar,
em semanas, uma plataforma que antes exigia um projeto de anos.

\rascunho{
  Relato: a experiência de montar ou migrar uma plataforma nesse modelo (armazém na
  nuvem, transformação em camadas, BI por cima), com o que deu certo e o que surpreendeu,
  inclusive no custo. Sem nomes de empresa e com números só em ordem de grandeza.
}

\section{Ideias antigas, nomes novos}
\label{sec:nomes-novos}

A área de dados tem o hábito de redescobrir as próprias ideias. O lakehouse, proposto em
2021, junta o armazenamento aberto e barato do data lake com as garantias de um data
warehouse: transações, esquemas, desempenho de consulta \parencite{armbrust2021}. É uma
evolução técnica real, e também uma volta ao ponto em que o data warehouse já estava,
agora sobre outra base.

A camada semântica, em que cada métrica é definida uma vez e reaproveitada por todas as
ferramentas, voltou à moda nos últimos anos. As suítes corporativas de BI dos anos 1990
já tinham camadas desse tipo. A organização em camadas de refinamento lembra a separação
entre a área de preparação e a área de apresentação de Kimball. E o
\enquote{dado como produto} lembra muito o data mart com dono.

Reconhecer isso não é desmerecer o novo. Cada volta acontece sobre uma tecnologia
melhor, mais barata e mais acessível, e por isso chega a mais gente. Mas reconhecer o
padrão ajuda a separar o que é de fato novo do que é rótulo, e a aproveitar décadas de
experiência acumulada em vez de reaprender tudo do zero. O próximo capítulo trata da
onda mais recente, e talvez da primeira em muito tempo que muda não só a infraestrutura,
mas também o trabalho de quem analisa os dados: a chegada da IA generativa.

\begin{leituras}
  \leitura{dageville2016}{O artigo que descreve a separação entre armazenamento e
    processamento no Snowflake. Ajuda a entender por que o custo da nuvem se comporta como
    se comporta.}
  \leitura{dean2004}{O artigo do MapReduce, base do Hadoop. Curto, claro e importante
    para entender a década do big data.}
  \leitura{kaminsky2019}{O texto que descreveu o papel do analytics engineer, ainda uma
    boa leitura para quem quer entender onde ele se encaixa entre análise e engenharia.}
  \leitura{reis2022}{Uma visão atual e independente de ferramenta do ciclo de vida da
    engenharia de dados, da geração ao consumo.}
\end{leituras}

% Capitulo: A chegada da IA
\chapter{A chegada da IA}
\label{cap:ia}

\rascunho{
  \textbf{Ideia do capítulo.} O que a IA generativa mudou no trabalho com dados e o que ela não mudou.
  \begin{itemize}
    \item Modelos de linguagem como copilotos: escrever SQL, explicar código, documentar.
    \item Text-to-SQL e a promessa de \enquote{conversar com os dados}: onde funciona e onde engana.
    \item Agentes que executam tarefas inteiras e a nova pergunta: quem revisa?
    \item O que não mudou: definição de métrica, qualidade do dado, contexto de negócio.
    \item O que a IA passou a pedir dos dados (treino e contexto) fica para a
      Parte~\ref{parte:modelagem}; o custo em tokens, para o capítulo~\ref{cap:custo-consumo}.
  \end{itemize}
}

\section{O copiloto na mesa do analista}
\label{sec:copiloto}

\pendente{escrever}

\section{Conversar com os dados}
\label{sec:text-to-sql}

\pendente{escrever}

\section{Agentes e a pergunta de quem revisa}
\label{sec:agentes}

\pendente{escrever}

\section{O que continua igual}
\label{sec:continua-igual}

\pendente{escrever}

% Capitulo: Os quatro estágios da IA no trabalho com dados
\chapter{Os quatro estágios da IA no trabalho com dados}
\label{cap:estagios-ia}

A adoção da IA no trabalho com dados não aconteceu de uma vez. Ela veio em estágios, cada
um dando à IA um pouco mais de autonomia e mudando um pouco mais o papel de quem trabalha
com ela. Primeiro, a IA era uma conversa numa janela à parte. Depois, passou a trabalhar
dentro das ferramentas de código. Em seguida, passou a construir, publicar e manter
aplicações inteiras. E agora começa a se tornar proativa: em vez de esperar que alguém
pergunte, ela acompanha os dados, decide dentro de limites combinados e avisa quando uma
decisão humana é necessária.

Este capítulo descreve esses quatro estágios (Figura~\ref{fig:estagios-ia}). Eles não são
uma escada que toda empresa sobe na mesma ordem, e convivem no mesmo time: a mesma pessoa
pode conversar com a IA de manhã e supervisionar um agente à tarde. Mas separá-los ajuda a
entender o que muda em cada um e, principalmente, o que cada estágio exige dos dados e das
pessoas. A ideia de níveis de automação não é nova. Os carros autônomos, por exemplo, são
classificados em níveis que dizem quem faz o quê, o sistema ou o motorista, em cada situação
\parencite{sae2021}. A pergunta que organiza este capítulo é a mesma.

\begin{figure}[htbp]
  \centering
  % Os quatro estagios da IA no trabalho com dados (cap. 6).
\begin{tikzpicture}[diagrama, x=1mm, y=1mm,
    degrau/.style={caixa, anchor=south west, minimum width=27mm, text width=24mm,
      minimum height=17mm, align=center, font=\scriptsize},
    humano/.style={rotulo, anchor=north, text width=27mm}]
  \node[degrau] (e1) at (0,0)  {\textbf{1. Conversa}\\[0.5mm]a IA responde numa janela à parte};
  \node[degrau] (e2) at (30,7) {\textbf{2. Código}\\[0.5mm]a IA trabalha dentro do editor e do repositório};
  \node[degrau] (e3) at (60,14) {\textbf{3. Aplicações}\\[0.5mm]a IA constrói, publica e mantém sistemas inteiros};
  \node[degrau, fill=fundo, line width=0.7pt] (e4) at (90,21)
    {\textbf{4. Proativa}\\[0.5mm]a IA monitora, decide nos limites e avisa};
  \node[humano] at (e1.south) [yshift=-1mm] {a pessoa pergunta, copia e integra};
  \node[humano] at (e2.south) [yshift=-1mm] {a pessoa dirige e revisa};
  \node[humano] at (e3.south) [yshift=-1mm] {a pessoa especifica e supervisiona};
  \node[humano] at (e4.south) [yshift=-1mm] {a pessoa define as regras e decide nas exceções};
  \draw[seta, draw=cinza] (0,-12) -- (117,-12);
  \node[rotulo, anchor=north] at (58.5,-12.5) {mais autonomia para a IA, mais importância para as regras e para a revisão};
\end{tikzpicture}

  \caption{Os quatro estágios da IA no trabalho com dados. A cada estágio, a IA ganha
    autonomia e o papel humano muda de fazer para definir regras e decidir nas exceções.}
  \label{fig:estagios-ia}
  \fonte{elaborado pelo autor.}
\end{figure}

\section{Estágio 1: a IA como conversa}
\label{sec:estagio-conversa}

O primeiro contato da maioria das pessoas com a IA generativa foi uma janela de conversa. A
pessoa escreve uma pergunta, recebe uma resposta e decide o que fazer com ela. Para quem
trabalha com dados, isso já foi uma mudança grande: pedir uma consulta, colar um erro e pedir
explicação, pedir ajuda para estruturar uma análise ou para escrever o e-mail com as
conclusões (capítulo~\ref{cap:ia}).

A característica que define esse estágio é o isolamento. A IA não vê o banco de dados, não
conhece o código do projeto nem os padrões do time. Tudo o que ela sabe é o que a pessoa
cola na conversa, e tudo o que ela produz precisa ser copiado, adaptado e integrado à mão.
A pessoa é a ponte entre a IA e o trabalho real.

Esse isolamento tem um lado bom: o controle é total, porque nada acontece sem que a pessoa
execute. E tem um lado ruim: o contexto é pobre, a integração é manual e cansativa, e é fácil
colar na conversa o que não deveria sair da empresa (\veja{sec:seguranca}). O estágio da
conversa continua útil, e não vai desaparecer, mas é o que menos aproveita o que a IA pode
fazer.

\section{Estágio 2: a IA conectada ao código}
\label{sec:estagio-codigo}

O segundo estágio levou a IA para dentro das ferramentas em que o trabalho acontece. Primeiro
vieram os assistentes integrados ao editor de código, que sugerem a próxima linha ou o
próximo bloco enquanto a pessoa escreve. Depois vieram os agentes de programação: assistentes
que leem o repositório inteiro, entendem a estrutura do projeto, executam comandos, rodam os
testes, corrigem o que falhou e preparam a alteração para revisão. A pessoa deixa de pedir
trechos de código e passa a pedir tarefas.

A diferença em relação ao estágio anterior é o contexto e a ação. O agente vê o código real,
os testes, as convenções do time, e consegue agir sobre eles. Uma tarefa como \enquote{crie
o modelo de receita por cliente seguindo o padrão das outras tabelas, com os testes de
unicidade e de frete} pode ser executada do começo ao fim, com a pessoa revisando o resultado.

Esse estágio trouxe também um vocabulário novo e um debate útil. Em 2025, a expressão
\emph{vibe coding} se popularizou para descrever um jeito de programar em que a pessoa
descreve o que quer, aceita o código gerado sem lê-lo e segue em frente enquanto parece
funcionar. Ela foi rapidamente distinguida do uso responsável dos assistentes: se a pessoa
revisou, testou e entendeu o código que a IA escreveu, isso não é \emph{vibe coding}, é
programar com um assistente \parencite{willison2025}. Para protótipos descartáveis, aceitar
sem ler pode ser aceitável. Para um modelo de dados que vai alimentar a decisão de alguém,
não é. A forma que tem ganhado espaço para o trabalho profissional costuma ser chamada de
engenharia com agentes: a IA executa, a pessoa dirige e verifica. É o que os
capítulos~\ref{cap:programar} e~\ref{cap:testes} defendem com outras palavras.

\section{Estágio 3: construir, publicar e gerenciar aplicações}
\label{sec:estagio-aplicacoes}

No terceiro estágio, a IA deixa de trabalhar só sobre partes de um projeto e passa a
construir aplicações inteiras: a interface, a lógica, o banco de dados, a publicação e, cada
vez mais, a manutenção. Uma pessoa descreve uma ferramenta interna, como um painel de
acompanhamento com alertas ou uma aplicação de conciliação de cadastros
(capítulo~\ref{cap:aplicacao-propria}), e em pouco tempo tem uma versão funcionando,
publicada e acessível ao time.

Para times de dados, isso muda a fronteira do que é possível. Ferramentas internas que antes
nunca entravam na fila de prioridades da TI, porque eram pequenas demais, agora podem ser
construídas pelo próprio time. Uma aplicação para coletar dados que hoje chegam por planilha,
uma tela para o steward revisar sugestões, um painel com uma lógica que a ferramenta de BI
não suporta: tudo isso ficou ao alcance.

O risco também mudou de lugar. Construir ficou fácil; manter continua difícil. Uma aplicação
publicada precisa de alguém que cuide de segurança, de permissões, de atualização de
dependências, de dúvidas dos usuários e de mudanças no negócio. Uma empresa que deixa cada
pessoa publicar a sua aplicação com ajuda da IA, sem nenhum critério, repete o problema dos
mil dashboards (\veja{sec:mil-dashboards}), agora com sistemas inteiros. As mesmas práticas
que o livro propõe para os dados valem aqui: versionamento, testes, um dono para cada
aplicação, e um critério claro do que pode ir para produção (Parte~\ref{parte:devops}).

\section{Estágio 4: a IA proativa}
\label{sec:estagio-proativa}

Nos três primeiros estágios, a IA espera ser chamada. Alguém faz uma pergunta, pede uma
tarefa ou descreve uma aplicação. O quarto estágio inverte essa relação. A IA passa a
acompanhar os dados continuamente, com as regras de negócio no contexto, percebe quando um
número mudou de forma relevante, analisa a situação e, conforme o caso, decide sozinha dentro
de limites combinados ou avisa quem precisa decidir, já com a análise pronta.

É uma mudança profunda na natureza do dado dentro da empresa. Durante toda a história contada
nos capítulos anteriores, o dado foi passivo: ficava no armazém, ou no painel, esperando que
alguém o consultasse, percebesse algo e decidisse. A crítica de Ackoff aos relatórios que
ninguém lia (\veja{sec:sad}) e o problema dos painéis abertos uma vez por semana
(capítulo~\ref{cap:frequencia}) são versões do mesmo limite: a informação existe, mas só gera
ação se uma pessoa estiver olhando no momento certo. Com a IA proativa, o dado deixa de
esperar.

Uma referência clássica da engenharia de automação ajuda a pensar esse estágio com cuidado.
\textcite{parasuraman2000} propõem que a automação pode atuar sobre quatro funções: obter a
informação, analisá-la, escolher a decisão e executar a ação, cada uma num nível diferente,
do totalmente manual ao totalmente automático. Os painéis tradicionais automatizam a primeira
função e parte da segunda. A IA proativa avança sobre a terceira e, em alguns casos, sobre a
quarta. A pergunta de projeto deixa de ser \enquote{a IA consegue?} e passa a ser \enquote{em
que nível cada função deve ficar, para cada tipo de decisão?}.

Uma proposta recente para agentes de IA organiza a resposta pelo papel que a pessoa assume:
operadora, colaboradora, consultora, aprovadora ou observadora. No nível do aprovador, o agente
age sozinho e só envolve a pessoa em situações de risco ou previamente combinadas; no do
observador, age com autonomia total e a pessoa apenas acompanha \parencite{feng2025}. Para
decisões de negócio, o nível do aprovador costuma ser o ponto de equilíbrio, e a forma de
chegar lá é parecida com o que as empresas já fazem com pessoas: definir alçadas. A
Tabela~\ref{tab:decisoes-proativa} mostra um exemplo de como esse desenho pode ficar.

\begin{table}[htbp]
  \centering
  \caption{Um desenho de alçadas para uma IA proativa.}
  \label{tab:decisoes-proativa}
  \small
  \begin{tabularx}{\textwidth}{@{}>{\raggedright\arraybackslash}p{0.22\textwidth}
      >{\raggedright\arraybackslash}X>{\raggedright\arraybackslash}p{0.2\textwidth}
      >{\raggedright\arraybackslash}p{0.2\textwidth}@{}}
    \toprule
    Tipo de situação & Exemplo & O que a IA faz & Papel humano \\
    \midrule
    Rotineira, baixo impacto & Repor um item de estoque dentro do limite combinado & Decide, executa e registra & Audita por amostragem \\
    \addlinespace
    Fora do padrão, impacto médio & Custo de frete muito acima do contratado numa rota & Analisa e recomenda & Aprova ou rejeita \\
    \addlinespace
    Alto impacto ou ambígua & Suspender as compras de um fornecedor & Avisa com a análise pronta & Decide \\
    \addlinespace
    Fora das regras conhecidas & Uma situação que nenhuma regra cobre & Avisa que não sabe o que fazer & Define a regra nova \\
    \bottomrule
  \end{tabularx}
  \fonte{elaborado pelo autor com base em \textcite{parasuraman2000} e \textcite{feng2025}.}
\end{table}

A última linha da tabela é talvez a mais importante. Uma IA proativa bem desenhada sabe
reconhecer quando a situação está fora das regras que recebeu, e nesse caso a ação correta é
avisar, e não improvisar. É a velha ideia de gestão por exceção: as pessoas deixam de
acompanhar tudo e passam a cuidar do que foge do normal.

\rascunho{
  Relato ou exemplo da minha experiência: um caso em que um alerta proativo (com ou sem IA)
  mudou uma decisão, ou um caso em que eu gostaria de ter tido esse tipo de alerta. Pode ser
  também uma visão de como isso se aplicaria a um dos domínios do livro (agricultura,
  finanças ou logística).
}

\section{O que cada estágio exige}
\label{sec:estagios-exigem}

Cada estágio exige mais dos dados e das pessoas do que o anterior. No estágio da conversa,
basta saber perguntar e conferir. No estágio do código, é preciso saber revisar, o que exige
saber programar (capítulo~\ref{cap:programar}). No estágio das aplicações, é preciso
disciplina de engenharia para que o que foi construído possa ser mantido.

O estágio proativo é o mais exigente, porque nele a IA age sobre os dados sem que alguém
esteja olhando a cada passo. Para isso funcionar, os dados precisam estar prontos de uma forma
que poucas empresas já estão: regras de negócio explícitas e escritas, métricas definidas uma
vez (\veja{sec:camada-semantica}), cadastros conciliados (capítulo~\ref{cap:master-data}),
testes que garantam que os números estão certos (\veja{sec:tdd}), atualização na frequência
certa e monitorada (capítulo~\ref{cap:frequencia}), e documentação que sirva de contexto
(capítulo~\ref{cap:contexto-ia}). Não é coincidência que essa lista resuma boa parte do resto
do livro. A IA proativa é o estágio em que todo o trabalho de base deixa de ser boa prática e
passa a ser condição.

Há também um risco humano, bem conhecido na pesquisa sobre automação. Quando um sistema
automático acerta quase sempre, as pessoas tendem a confiar demais nele e a prestar menos
atenção justamente nos casos em que ele erra \parencite{parasuraman2010}. Uma IA proativa que
pede aprovação dezenas de vezes por dia, com recomendações quase sempre certas, corre o risco
de ser aprovada no automático. O desenho das alçadas precisa levar isso em conta: poucas
aprovações, bem escolhidas, com a informação necessária para decidir, valem mais do que muitas
aprovações que viram rotina.

\begin{destaque}
No último estágio, o dado deixa de esperar ser analisado. A IA acompanha, decide dentro de
limites combinados e chama uma pessoa quando a decisão é dela. O trabalho humano passa a ser
definir as regras e cuidar das exceções.
\end{destaque}

Com este capítulo, termina a primeira parte do livro: a história de como chegamos até aqui e
de para onde a IA está levando o trabalho com dados. A próxima parte muda de pergunta: com
tudo isso acontecendo, o que vale a pena estudar, e em que ordem?

\begin{leituras}
  \leitura{parasuraman2000}{O modelo clássico de tipos e níveis de automação. Escrito para
    sistemas como aviação e controle industrial, serve muito bem para pensar a IA nas
    decisões de negócio.}
  \leitura{feng2025}{Uma proposta atual de níveis de autonomia para agentes de IA, organizada
    pelo papel que a pessoa assume em cada nível.}
  \leitura{parasuraman2010}{Para entender por que as pessoas passam a confiar demais em
    sistemas automáticos, e como desenhar a interação para evitar isso.}
\end{leituras}


\part{Trilhas de aprendizado}
\label{parte:trilhas}
% Capitulo: A trilha tradicional
\chapter{A trilha tradicional}
\label{cap:trilha-tradicional}

Quem começou a trabalhar com dados antes da IA generativa provavelmente percorreu um
caminho parecido: planilha, depois SQL, depois uma ferramenta de BI, e, para alguns,
estatística e programação. Ninguém desenhou essa trilha. Ela surgiu porque cada etapa
resolvia uma dor concreta que a anterior deixava sem resposta. Este capítulo descreve
esse caminho, mostra como ele era aprendido na prática e faz um balanço honesto do que
ele tinha de bom e de ruim. É o ponto de comparação para o próximo capítulo, que
pergunta o que muda quando a IA está disponível desde o primeiro dia.

\section{A sequência clássica}
\label{sec:sequencia-classica}

A porta de entrada quase sempre era a planilha. Tabela dinâmica, procura de valores entre
abas, gráficos, algumas fórmulas condicionais. Com isso já era possível responder boa parte
das perguntas de uma área: quanto vendemos, quem são os maiores clientes, como estamos em
relação à meta. A planilha ensinava, sem dizer, as primeiras noções de modelagem: que uma
tabela precisa de uma linha por registro, que juntar duas listas exige uma chave em comum,
que somar coisas de níveis diferentes dá errado.

O limite da planilha aparecia quando os dados cresciam ou quando a mesma análise precisava
ser refeita todo mês. Abrir um arquivo de um milhão de linhas, copiar e colar a extração do
sistema, refazer as mesmas fórmulas: em algum momento, alguém sugeria \enquote{pedir
acesso ao banco}. Era a entrada do SQL.

SQL é uma das linguagens mais duráveis da computação. Foi proposta em 1974, com o nome de
SEQUEL, como uma linguagem de consulta que se parecesse com inglês estruturado
\parencite{chamberlin1974}, e meio século depois continua sendo a forma padrão de conversar
com bancos de dados. Aprender \codigo{SELECT}, \codigo{JOIN} e \codigo{GROUP BY} mudava a
relação do analista com os dados. Ele deixava de depender da extração que alguém fazia e
passava a buscar exatamente o que precisava, direto na fonte.

O passo seguinte era a ferramenta de BI. Com os dados acessíveis via SQL, o analista
passava a montar painéis que se atualizavam sozinhos, distribuídos para a área inteira.
Para muitos, era aqui que a carreira se definia, com cursos e certificações da ferramenta
escolhida pela empresa.

Estatística vinha depois, quando vinha. Média, mediana, desvio, alguma noção de amostra e de
teste de hipótese. Muita gente fazia análise de dados por anos sem nunca ter estudado
estatística formalmente, o que nem sempre aparecia nos resultados, mas às vezes aparecia nas
conclusões. E, por fim, para quem seguia em direção à ciência ou à engenharia de dados,
vinha uma linguagem de programação: Python ou R, ou antes delas as macros de planilha.
Livros como os de \textcite{mckinney2022} e \textcite{wickham2023} formaram uma geração de
analistas que aprendeu a programar a partir de problemas de dados, e não o contrário.

\section{Aprender no trabalho}
\label{sec:aprender-no-trabalho}

Pouca gente aprendeu essa trilha numa sala de aula. O mais comum era aprender no trabalho,
empurrado por um problema: uma consulta que demorava demais, um painel que ninguém
conseguia manter, um número que não batia com o da outra área. A pessoa pesquisava, perguntava
para um colega mais experiente, procurava a mesma dúvida em fóruns e, depois de algumas
horas, resolvia. E, na próxima vez, já sabia.

Esse jeito de aprender tinha uma propriedade que só ficou evidente depois: ele era lento
de propósito, sem que ninguém tivesse planejado isso. Para resolver o problema, a pessoa
precisava entender o problema. O erro na consulta obrigava a ler a consulta linha por linha.
A resposta do fórum raramente servia sem adaptação, e adaptar exigia compreender. A pesquisa
sobre aprendizagem tem um nome para esse tipo de esforço que parece atrapalhar, mas
consolida o aprendizado: dificuldades desejáveis \parencite{bjork1994}. O próximo capítulo
volta a essa ideia, porque é justamente ela que a IA põe em risco.

\rascunho{
  Relato: como eu aprendi. Por onde entrei na área, qual foi a primeira ferramenta, o
  problema que me levou ao SQL, o que eu aprendi sozinho e o que aprendi com alguém mais
  experiente. Curto, um ou dois parágrafos, e com o que eu faria diferente hoje.
}

\section{Acertos e limites}
\label{sec:acertos-limites}

A trilha tradicional formou profissionais com fundamentos sólidos. Quem digitou cada
consulta e depurou cada erro desenvolveu uma intuição difícil de ensinar: olha um resultado
e desconfia que o total está alto demais, lê uma junção e percebe que ela vai duplicar
linhas. Também formou paciência com dado sujo, porque não havia atalho para limpar um
cadastro problemático.

Os limites também eram claros. O primeiro era o tempo: chegar à autonomia levava anos, e
muito desse tempo era gasto em tarefas repetitivas que não ensinavam nada novo depois da
décima vez. O segundo era a dependência de ferramenta. Como a formação se organizava em
torno do produto que a empresa usava, a carreira ficava presa a ele, e trocar de
ferramenta parecia recomeçar do zero, mesmo quando os conceitos eram os mesmos.

O terceiro limite era a separação entre análise e engenharia. O analista sabia SQL e
ferramenta de BI, mas raramente sabia versionar código, escrever testes ou automatizar uma
rotina. O engenheiro sabia montar pipelines, mas raramente entendia a pergunta de negócio
que os dados iam responder. Essa separação gerava as passagens de bastão que o
capítulo~\ref{cap:estrutura-times} discute, e foi uma das razões para o surgimento do
analytics engineer (\veja{sec:elt}).

O quarto limite era o mais silencioso: o conhecimento de negócio era aprendido por acaso.
Ninguém ensinava como a empresa ganhava dinheiro, quais números a diretoria olhava, onde a
operação costumava falhar. Quem aprendia isso, aprendia porque estava perto das pessoas
certas. Como o capítulo~\ref{cap:negocio} vai argumentar, essa é hoje a parte da trilha
que mais diferencia um profissional, e continua sendo a menos estruturada.

\begin{destaque}
A trilha tradicional era lenta e presa a ferramentas, mas formava intuição. O desafio de
hoje é manter a intuição sem manter a lentidão.
\end{destaque}

\begin{leituras}
  \leitura{mckinney2022}{Escrito pelo criador da biblioteca pandas, é a referência para
    aprender Python a partir de problemas de dados.}
  \leitura{wickham2023}{O equivalente em R, com uma abordagem muito didática do ciclo de
    importar, arrumar, transformar, visualizar e comunicar dados. Disponível gratuitamente
    na internet.}
  \leitura{bjork1994}{Para entender por que o esforço durante o aprendizado, que parece
    ineficiência, é o que faz o conhecimento ficar.}
\end{leituras}

% Capitulo: A trilha hoje
\chapter{A trilha hoje}
\label{cap:trilha-hoje}

Quem começa hoje tem algo que nenhuma geração anterior teve: um tutor disponível a
qualquer hora, que explica um erro com paciência infinita, gera exercícios sob medida e
revisa uma consulta em segundos. É uma vantagem enorme, e também uma armadilha. O mesmo
assistente que explica pode simplesmente fazer, e quem deixa a IA fazer durante o
aprendizado termina o curso sem ter aprendido. Este capítulo trata das duas coisas: como
usar a IA para aprender mais rápido, e como não usá-la para deixar de aprender. No fim,
propõe uma ordem de estudo para quem está começando agora.

\section{A IA como tutor}
\label{sec:ia-tutor}

Os usos da IA que mais ajudam no aprendizado têm algo em comum: neles, quem está
aprendendo continua fazendo o trabalho de pensar. Alguns exemplos:

\begin{itemize}
  \item \textbf{Explicar um erro.} Em vez de pedir a consulta corrigida, colar a mensagem
    de erro e pedir uma explicação do que ela significa e de onde costuma vir.
  \item \textbf{Explicar código alheio.} Pedir que a IA explique, linha por linha, uma
    consulta ou um script herdado, e depois tentar reescrevê-lo sem olhar.
  \item \textbf{Gerar exercícios.} Pedir dez exercícios de junção, em dificuldade
    crescente, sobre um conjunto de dados conhecido, com as respostas guardadas para
    depois.
  \item \textbf{Revisar uma solução.} Escrever a consulta primeiro e só então pedir uma
    revisão: o que pode dar errado, que caso de borda ficou de fora, como ficaria mais
    legível.
  \item \textbf{Ser questionado.} Pedir que a IA faça perguntas sobre um assunto recém
    estudado, sem dar as respostas, e corrija o que estiver errado.
\end{itemize}

A evidência mais clara de que a forma de usar faz diferença vem de um experimento com
quase mil estudantes de ensino médio, em aulas de matemática. Um grupo usou um assistente
parecido com os de uso geral; outro usou um assistente configurado como tutor, que conhecia
a solução de cada exercício e os erros comuns, mas tinha instrução de dar dicas
progressivas sem entregar a resposta. Os dois grupos foram muito melhor que o grupo de
controle nos exercícios feitos com ajuda. Na prova, sem ajuda, o grupo do assistente comum
foi pior do que quem nunca teve acesso, e o grupo do tutor ficou no mesmo nível do controle
\parencite{bastani2025}. O assistente sem limites virou muleta; o tutor, não.

\begin{dica}
Ao estudar com IA, comece o pedido com algo como: \enquote{Estou aprendendo. Não me dê a
resposta. Me ajude a chegar nela com perguntas e dicas, uma de cada vez.} A diferença no
que fica aprendido é grande.
\end{dica}

\section{O risco do atalho}
\label{sec:risco-atalho}

O risco não é hipotético nem restrito a estudantes de ensino médio. Em 2026, um experimento
com programadores, na maioria em início de carreira, pediu que aprendessem uma biblioteca de
Python que não conheciam. Metade podia usar um assistente de IA; metade programou à mão. No
teste de compreensão feito depois, o grupo com IA acertou, em média, metade das questões;
o grupo sem IA, cerca de dois terços. A maior diferença apareceu justamente nas questões de
depuração, a habilidade necessária para perceber quando o código gerado está errado
\parencite{shen2026}. E o grupo com IA nem chegou a terminar a tarefa significativamente
mais rápido.

O mesmo estudo traz a parte animadora. Entre os que usaram IA, os que se saíram melhor
foram os que a usaram para entender, fazendo perguntas sobre o código, e não só para gerar
o código. A ferramenta era a mesma; o resultado dependeu de como foi usada.

Uma pesquisa com profissionais do conhecimento aponta na mesma direção, por outro caminho:
quanto mais as pessoas confiavam na IA para uma tarefa, menos esforço de pensamento crítico
diziam aplicar a ela; quanto mais confiavam na própria capacidade, mais esforço aplicavam
\parencite{lee2025}. O atalho é confortável, e o conforto é exatamente o problema.

A explicação para tudo isso é antiga. A pesquisa sobre aprendizagem mostra há décadas que
o esforço de recuperar uma informação, errar e corrigir é o que consolida o conhecimento
\parencite{bjork1994}, e que a habilidade de especialista vem de prática deliberada: tarefas
um pouco além da capacidade atual, com feedback, repetidas com atenção
\parencite{ericsson1993}. Quando a IA faz a parte difícil, ela remove justamente o esforço
que geraria o aprendizado. O resultado parece bom na hora e não fica.

\begin{destaque}
Usar IA para fazer é ótimo para quem já sabe. Para quem está aprendendo, é preciso usá-la
para entender, ou ela ensina só a pedir.
\end{destaque}

Isso não significa estudar sem IA. Significa separar dois momentos. No momento de
aprender uma coisa nova, a IA explica, pergunta e revisa, mas quem escreve é a pessoa. No
momento de produzir algo que já se sabe fazer, a IA pode escrever, e a pessoa revisa. O
erro é usar o segundo modo no primeiro momento.

\section{Uma ordem de estudo para hoje}
\label{sec:ordem-estudo}

Com essas ideias em mente, esta é a ordem que eu recomendaria para quem começa agora em
análise ou engenharia de dados. Ela não é muito diferente da trilha tradicional no
conteúdo. Muda na ordem, na profundidade de cada etapa e na forma de usar a IA em cada uma.

\begin{enumerate}
  \item \textbf{Dados em tabela e SQL, à mão.} Seleção, filtro, agregação, junções, funções
    de janela. É a base de tudo e a linguagem que mais tempo vai durar. A IA entra como
    tutor e gerador de exercícios, nunca como quem escreve as consultas.
  \item \textbf{Modelagem.} Grão, chaves, fatos e dimensões (capítulo~\ref{cap:dw}). Na
    trilha tradicional isso vinha tarde, quando vinha. Hoje vem cedo, porque é o que
    permite avaliar se o que a IA produziu faz sentido.
  \item \textbf{Python para dados.} Ler, limpar, transformar e visualizar dados com
    código. Aqui a IA já pode ajudar mais, desde que cada trecho gerado seja lido e
    entendido antes de ser usado.
  \item \textbf{Estatística aplicada e pensamento crítico.} Distribuições, comparação,
    incerteza, correlação e causalidade (\veja{sec:fund-estatistica}). É o que impede
    conclusões erradas tiradas de números certos.
  \item \textbf{Visualização e comunicação.} Escolher o gráfico, escrever a conclusão
    primeiro, apresentar para quem decide (\veja{sec:visualizacao}).
  \item \textbf{Práticas de engenharia.} Versionamento, testes, automação
    (Parte~\ref{parte:devops}). Antes eram território do engenheiro; hoje qualquer
    analista que trabalha com IA precisa delas para revisar e manter o que produz.
  \item \textbf{IA como ferramenta de produção.} Usar assistentes e agentes para
    acelerar o que já se sabe fazer, com critério para revisar e saber de que lado da
    fronteira está cada tarefa (\veja{sec:copiloto}).
\end{enumerate}

Em paralelo a tudo isso, desde o primeiro dia, aprender o negócio onde os dados nascem.
O capítulo~\ref{cap:negocio} explica por que essa é a parte que mais diferencia um
profissional, e o capítulo~\ref{cap:mapas} detalha a trilha para cada carreira.

\rascunho{
  Opcional: um parágrafo meu sobre como eu oriento hoje quem entra no time, ou o que
  observo de diferente entre quem aprendeu antes e quem está aprendendo com IA.
}

\begin{leituras}
  \leitura{bastani2025}{O experimento com estudantes que mostra a diferença entre usar a
    IA como muleta e como tutor. Leitura rápida e muito clara.}
  \leitura{shen2026}{Um experimento controlado sobre como a IA afeta a formação de
    habilidades de programação, com atenção especial à depuração.}
  \leitura{ericsson1993}{O artigo clássico sobre prática deliberada, a base de boa parte do
    que se sabe sobre como alguém se torna especialista.}
\end{leituras}

% Capitulo: Ainda vale aprender a programar?
\chapter{Ainda vale aprender a programar?}
\label{cap:programar}

Se a IA escreve SQL e Python a partir de uma descrição em português, por que alguém
deveria gastar meses aprendendo essas linguagens? A pergunta é legítima, e a resposta deste
capítulo é que sim, vale, mas por motivos diferentes dos de antes. Antes, se aprendia a
programar principalmente para escrever código. Hoje, se aprende principalmente para ler,
revisar, depurar e decidir sobre código que outra coisa escreveu. Essa diferença muda o
que estudar e com que profundidade, e é o argumento central deste livro.

\section{Ler antes de assinar}
\label{sec:ler-antes-assinar}

Quem coloca em produção uma consulta gerada por IA está assinando embaixo dela. Se o
número do painel estiver errado, a pergunta na reunião não vai ser \enquote{qual modelo
gerou isso?}, e sim \enquote{quem publicou isso?}. Quem não consegue ler o código que
publica está assinando sem ler.

Ler código bem é mais difícil do que parece. Brian Kernighan e P. J. Plauger observaram,
ainda nos anos 1970, que depurar é bem mais difícil do que escrever, e que quem escreve o
código mais esperto que consegue fica, por definição, sem capacidade de depurá-lo
\parencite{kernighan1978}. \pendente{conferir a formulação e a página no livro} Com a IA,
essa observação ganha um sentido novo. O código chega pronto, muitas vezes mais elaborado do
que a pessoa escreveria sozinha, e precisa ser entendido por alguém que não passou pelo
processo de escrevê-lo. É depuração desde o primeiro minuto.

O experimento sobre formação de habilidades citado no capítulo anterior encontrou a maior
perda justamente aí: quem aprendeu com IA foi pior nas questões de depuração
\parencite{shen2026}. É um círculo vicioso possível. Quem depende da IA para escrever
aprende menos a depurar, e quem não sabe depurar não consegue revisar o que a IA escreve.

\section{Onde a IA ajuda e onde ela não decide}
\label{sec:ajuda-nao-decide}

A IA ajuda muito em quase todas as tarefas de código em dados: escrever a primeira versão
de uma consulta, explicar um plano de execução, sugerir um índice, converter código de um
dialeto de SQL para outro, propor testes. Mas há uma diferença entre ajudar e decidir, e
algumas decisões continuam precisando de alguém que entende o que está acontecendo.

Uma é a de desempenho e custo. Num armazém de dados cobrado por uso, a diferença entre uma
consulta que lê a tabela inteira e uma que lê só a partição necessária pode ser grande na
fatura do mês. A IA sugere as duas com a mesma confiança. Saber qual é a certa exige
entender como o banco executa a consulta, o que é uma habilidade de programação, não de
redação de pedidos.

Outra é a de adequação ao contexto. A IA não sabe que aquela tabela tem duplicidade
conhecida, que aquele campo deixou de ser preenchido no ano passado, que a área financeira
fecha o mês com dois dias de atraso. Quem conhece o código e os dados da empresa sabe, e é
essa pessoa que decide se a solução proposta serve.

E há a decisão de quando parar de pedir e escrever. O estudo com desenvolvedores
experientes citado no capítulo~\ref{cap:ia} mostrou que, em projetos grandes e bem
conhecidos, parte do tempo ganho com a IA era perdido revisando e corrigindo sugestões
\parencite{becker2025}. Para quem sabe programar, às vezes é mais rápido escrever do que
explicar. Para quem não sabe, essa opção simplesmente não existe.

\section{Segurança e privacidade}
\label{sec:seguranca}

Código gerado por IA herda os padrões do código com que o modelo foi treinado, inclusive os
inseguros. Um estudo que analisou programas gerados por um assistente de código em cenários
sensíveis encontrou vulnerabilidades em cerca de 40\% deles \parencite{pearce2022}. Outro,
com usuários reais, mostrou um efeito mais preocupante: quem tinha acesso a um assistente
escreveu código menos seguro do que quem não tinha, e ao mesmo tempo acreditava mais que o
próprio código era seguro \parencite{perry2023}. O mesmo estudo observou que os
participantes que desconfiavam mais da IA e trabalhavam mais os próprios pedidos produziam
código mais seguro.

Em dados, os riscos mais comuns são concretos: consultas montadas pela concatenação de
texto vindo do usuário, credenciais escritas direto no código, permissões amplas demais
concedidas \enquote{para funcionar}. E há riscos novos, próprios de sistemas que usam
modelos de linguagem, como instruções maliciosas escondidas nos dados que o modelo lê, que
o fazem agir de forma diferente da esperada. A lista de riscos mantida pela OWASP para
aplicações com modelos de linguagem é um bom ponto de partida \parencite{owasp2025}.

A privacidade merece um cuidado à parte. Tudo o que vai num pedido para um modelo externo
sai do ambiente da empresa. No Brasil, dado pessoal é protegido pela Lei Geral de Proteção
de Dados \parencite{lgpd2018}, e mandar um cadastro de clientes para um serviço externo
pode configurar compartilhamento com terceiro, com as obrigações que isso traz.
\pendente{conferir com referência jurídica antes de manter a afirmação} A regra prática é
simples: dado pessoal, credencial e informação sensível do negócio não entram em pedido
para modelo nenhum sem que a empresa tenha decidido, por escrito, que aquilo é permitido.

\begin{dica}
Antes de colar qualquer coisa num assistente de IA, pergunte: eu mandaria isso por e-mail
para alguém de fora da empresa? Se a resposta for não, não cole.
\end{dica}

\section{A linguagem como forma de pensar}
\label{sec:linguagem-pensar}

Há um motivo para aprender a programar que não tem a ver com revisar código. Uma
linguagem de programação é também uma forma de pensar sobre o problema.

Quem aprendeu SQL de verdade pensa em conjuntos: não \enquote{para cada pedido, some os
itens}, e sim \enquote{agrupe os itens por pedido e some}. Essa forma de pensar ajuda a
formular perguntas que os dados conseguem responder, a perceber de antemão onde uma junção
vai multiplicar linhas, a desconfiar de um total que cresceu demais. Quem aprendeu Python
para dados aprende a decompor um problema em etapas pequenas e verificáveis. São hábitos
mentais que continuam úteis mesmo quando é a IA que escreve.

E há a questão da durabilidade. SQL tem meio século (\veja{sec:sequencia-classica}) e
continua sendo a língua franca dos dados. As ferramentas de BI mudam a cada poucos anos, os
assistentes de IA mudam a cada poucos meses, e o SQL que alguém aprendeu bem há vinte anos
continua servindo. Poucas coisas na área têm esse retorno sobre o tempo de estudo.

\section{Um exemplo de consulta errada}
\label{sec:consulta-errada}

Para tornar o argumento concreto, considere um pedido simples a um assistente: \enquote{a
receita por cliente, incluindo o frete}. Os dados estão em duas tabelas: \codigo{pedidos},
com uma linha por pedido e o valor do frete, e \codigo{itens\_pedido}, com uma linha por item
de cada pedido. Uma resposta muito plausível é esta:

\begin{sql}
-- Receita por cliente, incluindo frete (versão com erro)
SELECT
    p.id_cliente,
    SUM(i.valor_item) + SUM(p.valor_frete) AS receita
FROM pedidos AS p
JOIN itens_pedido AS i ON i.id_pedido = p.id_pedido
GROUP BY p.id_cliente;
\end{sql}

A consulta roda, devolve um número para cada cliente, e o número parece razoável. Está
errado. A junção repete cada pedido uma vez por item, e com ele repete o frete. Um pedido
com três itens e frete de R\$~30 entra na soma com R\$~90 de frete. Quanto mais itens por
pedido, maior o erro, e ele não aparece em nenhuma mensagem: a consulta simplesmente
devolve um número alto demais.

A correção é somar os itens no grão do pedido antes de juntar com a tabela de pedidos:

\begin{sql}
-- Receita por cliente, incluindo frete (versão corrigida)
WITH itens_por_pedido AS (
    SELECT id_pedido, SUM(valor_item) AS valor_itens
    FROM itens_pedido
    GROUP BY id_pedido
)
SELECT
    p.id_cliente,
    SUM(COALESCE(i.valor_itens, 0) + p.valor_frete) AS receita
FROM pedidos AS p
LEFT JOIN itens_por_pedido AS i ON i.id_pedido = p.id_pedido
GROUP BY p.id_cliente;
\end{sql}

Quem percebe o erro na primeira versão não é quem sabe pedir melhor. É quem olha para a
junção e se pergunta qual é o grão de cada tabela, a pergunta que o
capítulo~\ref{cap:dw} chamou de a mais importante da modelagem (\veja{sec:modelagem-dimensional}).
E quem tem o hábito de conferir: o total de frete na consulta é igual ao total de frete na
tabela de pedidos? Essa conferência leva um minuto e pegaria o erro. Ela também pode virar
um teste automático, escrito antes da consulta, como propõe o capítulo~\ref{cap:testes}.

\rascunho{
  Relato opcional: um erro parecido, gerado por IA ou por uma pessoa, que chegou a um
  painel ou a uma reunião antes de ser percebido.
}

\begin{destaque}
Aprender a programar deixou de ser principalmente sobre escrever código. Passou a ser
sobre saber se o código está certo.
\end{destaque}

\begin{leituras}
  \leitura{perry2023}{O estudo que mostra que assistentes de IA podem levar a código menos
    seguro e a mais confiança nele. Vale ler a parte sobre o comportamento de quem se saiu
    melhor.}
  \leitura{owasp2025}{A lista de riscos de segurança específicos de aplicações com modelos
    de linguagem, mantida e atualizada pela comunidade.}
  \leitura{kernighan1978}{Um clássico curto sobre escrever código claro. Os exemplos são
    antigos; as lições não.}
\end{leituras}

% Capitulo: Fundamentos que não envelhecem
\chapter{Fundamentos que não envelhecem}
\label{cap:fundamentos}

Ferramentas envelhecem rápido. A ferramenta de BI que alguém dominava há dez anos
provavelmente já foi substituída, e o assistente de IA de hoje vai parecer antigo em dois
anos. Algumas habilidades, porém, atravessam todas as ondas que a Parte~\ref{parte:origens}
descreveu sem perder valor. Este capítulo reúne cinco delas: modelagem, estatística e
pensamento crítico, qualidade de dados, comunicação e conhecimento do negócio. Nenhuma
depende de produto, e todas ficaram mais importantes com a IA, porque são justamente o que
permite avaliar o que a máquina produz.

\section{Modelagem}
\label{sec:fund-modelagem}

Modelar dados é decidir como o mundo vai ser representado em tabelas: o que é uma
entidade, o que identifica cada registro, como as coisas se relacionam, o que uma linha
representa. O capítulo~\ref{cap:dw} apresentou duas formas clássicas de fazer isso, a
normalizada e a dimensional. Aqui interessa o que elas têm em comum.

Toda boa modelagem responde a três perguntas antes de qualquer outra. Qual é o grão, ou
seja, o que exatamente uma linha representa? Qual é a chave, ou seja, o que torna cada
linha única? E como esta tabela se liga às outras? Quem pergunta isso por hábito evita a
maior parte dos erros de análise, inclusive os que a IA comete com toda a confiança, como o
frete multiplicado do capítulo anterior (\veja{sec:consulta-errada}).

Modelagem é também a habilidade que mais se valoriza com a IA. Quando escrever código fica
barato, o que diferencia um projeto bom de um ruim é o desenho: as tabelas certas, com o
grão certo e os nomes certos. A Parte~\ref{parte:modelagem} mostra que esse desenho muda
conforme o uso, mas as três perguntas continuam as mesmas.

\section{Estatística e pensamento crítico}
\label{sec:fund-estatistica}

Não é preciso ser estatístico para trabalhar com dados, mas é preciso pensar como um em
alguns momentos. Na maior parte do tempo, isso se resume a poucas ideias, aplicadas com
disciplina.

A primeira é olhar a distribuição antes do resumo. Uma média esconde tudo o que importa
quando os dados são assimétricos: o tempo médio de entrega pode estar ótimo enquanto um
em cada dez pedidos atrasa uma semana. Mediana, percentis e um histograma simples contam a
história que a média esconde.

A segunda é sempre comparar. Um número sozinho não diz nada. Cresceu em relação a quê? Está
alto em comparação com o quê? A variação está dentro do normal ou é sinal de alguma
mudança? Boa parte do trabalho de análise é escolher a comparação certa.

A terceira é desconfiar de relações simples. Correlação não é causalidade, e a frase é tão
repetida que perdeu a força. Um exemplo clássico ajuda a recuperá-la. Em 1973, os dados de
admissão na pós-graduação da Universidade da Califórnia em Berkeley mostravam que homens
eram aceitos em proporção maior do que mulheres. Olhando departamento por departamento, o
viés desaparecia, e em vários deles as mulheres tinham taxa de aceitação um pouco maior. A
explicação: as mulheres se candidatavam mais aos departamentos mais concorridos, que
recusavam a maioria dos candidatos de qualquer sexo \parencite{bickel1975}. O agregado
contava uma história; os grupos, outra. Esse fenômeno, em que uma tendência se inverte
quando os dados são separados em grupos, é conhecido como paradoxo de Simpson.

Num painel de vendas, o mesmo efeito aparece quando a margem média cai porque a empresa
vendeu mais de um produto de margem baixa, mesmo que a margem de cada produto tenha subido.
Quem não pensa em composição tira a conclusão errada de um número correto.

A quarta ideia é lembrar que números podem ser apresentados de forma enganosa, por
intenção ou por descuido. O pequeno livro de \textcite{huff1954}, com mais de setenta anos,
continua sendo a melhor introdução ao assunto: eixos cortados, médias escolhidas a dedo,
amostras viciadas. E vale conhecer os atalhos do nosso próprio raciocínio, que tendem a
construir uma história coerente com o pouco que está à vista e a ignorar o que não está
\parencite{kahneman2011}. Um assistente de IA que responde com segurança e fluência explora
exatamente essa tendência.

\section{Qualidade de dados}
\label{sec:fund-qualidade}

Qualidade de dados costuma ser tratada como exatidão: o valor está certo ou errado. Um
estudo clássico mostrou que, para quem usa os dados, qualidade é bem mais do que isso.
\textcite{wang1996} perguntaram a consumidores de dados o que tornava um dado bom e
organizaram as respostas em quatro grupos: qualidade intrínseca (o dado é correto e
confiável), qualidade contextual (é relevante, atual e completo para a tarefa), qualidade
de representação (é compreensível e consistente) e acessibilidade (dá para obter, com
segurança).

A consequência prática é que qualidade depende do uso. Um dado atualizado uma vez por mês
tem ótima qualidade para o fechamento mensal e péssima para o acompanhamento diário. Um
cadastro com o nome do cliente escrito de três formas é aceitável para contar pedidos e
inaceitável para somar a receita por cliente. Perguntar \enquote{este dado é bom?} sem
perguntar \enquote{bom para quê?} não leva a lugar nenhum.

Na prática do dia a dia, isso vira hábitos simples: conferir se as chaves são únicas,
contar nulos, comparar totais com a fonte, olhar a data da última atualização antes de
confiar num número. Esses hábitos viram testes automáticos no capítulo~\ref{cap:testes} e
uma disciplina de governança no capítulo~\ref{cap:governanca}.

\section{Comunicação}
\label{sec:fund-comunicacao}

Uma análise que ninguém entende não muda nenhuma decisão. Comunicar é parte do trabalho, e
não um enfeite no fim dele.

A recomendação mais útil sobre escrita para quem trabalha com dados vem da consultoria:
comece pela conclusão. Barbara Minto propôs organizar qualquer texto como uma pirâmide,
com a ideia principal no topo e os argumentos que a sustentam embaixo
\parencite{minto2009}. Um e-mail que começa com \enquote{as vendas caíram 8\% por causa da
ruptura de estoque em duas regiões} e só depois mostra os números é lido; um que começa
pela metodologia e chega à conclusão no quinto parágrafo, não.

O mesmo vale para gráficos e apresentações (\veja{sec:visualizacao}): saber quem vai ler,
o que essa pessoa precisa decidir, e colocar a mensagem no título, não numa legenda que
ninguém lê. E vale, cada vez mais, para a forma como se escreve para a IA. Um pedido claro,
com contexto, objetivo e critério de qualidade, é um pequeno exercício de comunicação, e
quem escreve bem para pessoas costuma escrever bem para máquinas.

\section{Conhecer o negócio}
\label{sec:fund-negocio}

O quinto fundamento é o que menos se ensina e o que mais diferencia: entender o negócio em
que os dados nascem. Saber como a empresa ganha dinheiro, quais decisões dependem dos
números e onde a operação costuma falhar transforma uma análise correta numa análise útil.

Os quatro fundamentos anteriores ficaram mais acessíveis com a IA, que ajuda a modelar, a
calcular, a conferir e a escrever. Este é o único em que ela pouco ajuda, porque o
conhecimento de negócio está nas pessoas e nos processos da empresa, e não nos dados de
treinamento de um modelo. Por isso ele ganha o próximo capítulo inteiro.

\begin{leituras}
  \leitura{huff1954}{Curto, divertido e ainda atual. A melhor vacina contra gráficos e
    números enganosos.}
  \leitura{wang1996}{O estudo que mostrou que qualidade de dados, para quem usa, vai muito
    além da exatidão.}
  \leitura{minto2009}{O método da pirâmide para organizar textos e apresentações. Muda a
    forma de escrever qualquer e-mail de análise.}
  \leitura{kahneman2011}{Para entender os atalhos do raciocínio que levam a conclusões
    erradas, inclusive diante de números corretos.}
\end{leituras}

% Capitulo: Conhecer o negócio
\chapter{Conhecer o negócio}
\label{cap:negocio}

Durante muito tempo, o que separava um bom profissional de dados de um mediano era a
técnica: quem escrevia SQL melhor, quem dominava a ferramenta, quem montava o pipeline mais
robusto. Com a IA acelerando a parte técnica, essa diferença diminuiu. Quem sabe pouco
consegue produzir algo razoável, e quem sabe muito produz mais rápido. O que passa a
separar o bom do excelente é outra coisa: a capacidade de antecipar o que o negócio precisa
e entregar o que ainda não foi pedido. Essa capacidade vem de conhecer a fundo um negócio
específico, e é o assunto deste capítulo.

\section{O gargalo mudou de lugar}
\label{sec:gargalo}

Todo processo tem um gargalo, a etapa que limita a velocidade do todo. No trabalho com
dados, o gargalo costumava ser a construção: escrever a consulta, montar o modelo, desenhar
o painel. Com a IA, construir ficou muito mais rápido, e o gargalo se deslocou para antes:
saber o que construir.

A IA responde muito bem ao que foi perguntado. Ela não sabe o que deixou de ser perguntado.
Se o gestor pede a receita por região, ela entrega a receita por região. Não pergunta se a
queda numa região é efeito de um cliente grande que mudou de prazo de pagamento, nem
sugere olhar a margem junto com a receita porque o problema daquela empresa costuma ser
desconto, e não volume. Essas perguntas exigem saber como aquele negócio funciona.

Esse ponto não é novo na literatura de ciência de dados. O modelo de processo mais usado
para projetos de mineração de dados, o CRISP-DM, começa por uma fase chamada entendimento
do negócio, antes de qualquer contato com os dados \parencite{chapman2000}. E
\textcite{provost2013} defendem que o mais importante para um projeto de dados dar certo é
o pensamento analítico orientado ao problema de negócio, mais do que o domínio de qualquer
algoritmo. O que muda com a IA é a proporção: quando a técnica fica barata, o entendimento
do negócio passa a ser quase todo o diferencial.

\section{Do bom ao excelente: entregar o que não foi pedido}
\label{sec:bom-excelente}

O bom profissional entrega o pedido com qualidade. O número está certo, o painel está
claro, o prazo foi cumprido. O excelente faz isso e mais uma coisa: percebe a pergunta por
trás do pedido e chega com a resposta para a próxima.

Um pedido como \enquote{preciso do relatório de vendas por vendedor} raramente é sobre o
relatório. Pode ser sobre a revisão das metas do semestre, sobre um vendedor que o gestor
desconfia que está concedendo desconto demais, ou sobre a preparação de uma reunião com a
diretoria. Cada uma dessas perguntas pede um recorte diferente. Quem conhece o negócio
pergunta o que está por trás, ou já sabe, e entrega o que vai ser útil, não só o que foi
pedido.

Antecipar também significa perceber necessidades que ninguém verbalizou. Notar que o
fechamento do mês sempre atrasa porque duas áreas conciliam números à mão, e propor uma
base única antes que alguém peça. Perceber que a operação vai precisar acompanhar um
indicador novo porque mudou a regra de um contrato. Esse tipo de entrega não aparece em
nenhuma lista de requisitos e é, muitas vezes, a que mais gera valor.

\begin{destaque}
A IA entrega bem o que foi pedido. Quem conhece o negócio entrega o que era preciso.
\end{destaque}

Há um cuidado aqui. Antecipar não é decidir pelo negócio, nem encher o gestor de análises
que ele não pediu. É fazer as perguntas certas cedo e propor o que tem chance real de ser
útil, sabendo que a decisão continua sendo de quem é dono do problema.

\section{Três domínios: agricultura, finanças e logística}
\label{sec:tres-dominios}

Cada setor tem os seus ciclos, o seu vocabulário e as suas armadilhas, e é isso que o
profissional de dados precisa aprender para antecipar alguma coisa. Três setores ajudam a
ilustrar o que muda de um para outro.

Na \textbf{agricultura}, o tempo é o da safra. As decisões importantes acontecem em
janelas curtas (plantio, aplicação de insumos, colheita) e dependem de clima, de solo e de
preço, coisas que mudam de um ano para outro. A unidade de análise muitas vezes não é o
cliente nem o produto, e sim a área: a fazenda, o talhão, o ponto de amostragem. Comparar
um mês com o mesmo mês do ano anterior pode não fazer sentido se a safra atrasou. E a
variabilidade espacial é parte do problema: a média de uma fazenda esconde as diferenças
entre os talhões que justificam um manejo diferente em cada um.

Nas \textbf{finanças}, o tempo é o do fechamento. Há a data em que a venda aconteceu, a data
em que foi faturada e a data em que o dinheiro entrou, e confundir competência com caixa
gera discussões intermináveis sobre qual é o número certo. A precisão importa ao centavo, a
reconciliação entre sistemas é rotina, e um dado que muda depois do fechamento é um
problema sério. Quem conhece esse mundo sabe que o painel financeiro precisa dizer em que
critério cada número está e se o período já foi fechado.

Na \textbf{logística}, o tempo é o da operação, muitas vezes medido em horas. Os
indicadores giram em torno de prazo, ocupação e custo: entregas no prazo e completas, uso
da capacidade dos veículos e dos armazéns, custo por entrega ou por quilômetro. A demanda
tem sazonalidade forte, a operação depende de muitos parceiros, e um atraso de um dia pode
custar um cliente. Aqui, mais do que nos outros dois, a frequência de atualização dos dados
é decisiva (capítulo~\ref{cap:frequencia}).

\rascunho{
  Relatos: um caso de cada domínio em que eu antecipei uma necessidade, ou seja, entreguei
  algo que ainda não tinha sido pedido e que se mostrou importante. Para cada caso: o
  contexto (sem nome de empresa), o que eu percebi, o que entreguei e o que aconteceu
  depois. Usar uma caixa \codigo{relato} por caso.
}

O ponto de mostrar os três lado a lado é que o conhecimento técnico é o mesmo nos três, e
o que torna alguém excelente em cada um é completamente diferente. Quem muda de setor leva
a técnica consigo e precisa reconstruir o resto.

\section{Como estudar um domínio de negócio}
\label{sec:estudar-dominio}

Diferente da técnica, o conhecimento de negócio não vem em curso estruturado. Mas dá para
estudá-lo de forma deliberada, e alguns hábitos aceleram muito o processo.

\begin{itemize}
  \item \textbf{Conversar com a operação.} Passar um tempo com quem gera os dados: o
    vendedor, o operador do armazém, o agrônomo, o analista do contas a pagar. Entender
    como o dado nasce explica a maior parte das inconsistências que aparecem depois.
  \item \textbf{Ler o que a diretoria lê.} Os relatórios de resultado, as apresentações de
    planejamento, as atas de reunião que estiverem acessíveis. Eles mostram o que de fato
    importa para quem decide.
  \item \textbf{Aprender o vocabulário.} Cada setor tem termos próprios, e muitos
    significam coisas diferentes em áreas diferentes da mesma empresa. Um glossário pessoal,
    mantido ao longo do tempo, vale muito.
  \item \textbf{Acompanhar um ciclo completo.} Uma safra, um ano fiscal, um pico de
    demanda. Muitas regras só fazem sentido quando se vê o ciclo inteiro.
  \item \textbf{Estudar o setor fora da empresa.} Publicações setoriais, dados públicos,
    associações, cursos técnicos de curta duração. Ajudam a entender o que é particular da
    empresa e o que é do setor.
  \item \textbf{Perguntar o porquê.} Diante de qualquer pedido, perguntar que decisão ele
    vai apoiar. Com o tempo, as respostas formam um mapa de como o negócio funciona.
\end{itemize}

A IA pode ajudar nesse estudo, explicando conceitos de um setor, resumindo publicações e
sugerindo perguntas para fazer à operação. O que ela não substitui é o contato com as
pessoas e com o dia a dia da empresa, onde está o conhecimento que nenhum modelo tem.

Esse conhecimento também se conecta com o resto do livro. É ele que permite escrever os
testes certos antes do modelo, porque quem conhece o negócio sabe quais regras não podem
quebrar (\veja{sec:tdd}). E é ele que permite levar ao quadro do time uma história já bem
definida, em vez de um pedido vago (Parte~\ref{parte:agil}).

\begin{leituras}
  \leitura{provost2013}{Um livro sobre ciência de dados escrito para quem precisa pensar
    em problemas de negócio, não só em algoritmos.}
  \leitura{chapman2000}{O guia original do CRISP-DM. Vale ler ao menos a fase de
    entendimento do negócio, que costuma ser a mais pulada.}
  \leitura{davenport2007}{Uma visão gerencial de como empresas usam dados para competir, útil
    para entender o que a diretoria espera do time de dados.}
\end{leituras}

% Capitulo: Mapas de trilhas
\chapter{Mapas de trilhas}
\label{cap:mapas}

Os capítulos anteriores trataram do que estudar em termos gerais. Este fecha a
Parte~\ref{parte:trilhas} com roteiros mais concretos para três carreiras: analista de
dados, analytics engineer e engenheiro de dados. Para cada uma, o que continua igual em
relação à trilha tradicional e o que muda com a IA. No fim, algumas recomendações sobre
como mostrar o que se sabe, porque, numa época em que qualquer um gera código, um portfólio
precisa mostrar outra coisa.

\section{Três carreiras}
\label{sec:tres-carreiras}

As fronteiras entre os papéis variam muito de empresa para empresa, e o
capítulo~\ref{cap:estrutura-times} discute como elas estão se movendo. Ainda assim, há um
centro reconhecível em cada um.

O \textbf{analista de dados} está mais perto do negócio. Recebe perguntas, investiga os
dados, constrói análises e painéis e comunica conclusões. As ferramentas centrais são SQL,
uma ferramenta de visualização e, cada vez mais, Python. O que define um bom analista é a
qualidade das perguntas que faz e a clareza com que comunica as respostas.

O \textbf{analytics engineer} fica no meio. Transforma os dados brutos em modelos limpos,
testados e documentados, que os analistas e as ferramentas consomem
\parencite{kaminsky2019}. Trabalha com SQL, modelagem dimensional, versionamento, testes e
ferramentas de transformação como código. O que o define é a capacidade de criar bases
confiáveis que outras pessoas usam sem precisar perguntar como foram feitas.

O \textbf{engenheiro de dados} está mais perto da infraestrutura. Constrói e mantém a
ingestão dos dados, as plataformas, a orquestração, o desempenho e o custo do ambiente
\parencite{reis2022}. Usa Python, SQL, ferramentas de nuvem, orquestradores e as práticas
de engenharia de software em geral. O que o define é a confiabilidade: os dados chegam,
chegam no prazo e chegam inteiros.

O cientista de dados, que constrói modelos preditivos e prescritivos, compartilha muito com
o analista e com o engenheiro, e a forma de organizar dados para o trabalho dele está no
capítulo~\ref{cap:modelagem-cd}.

\section{Tradicional contra hoje}
\label{sec:comparativo-trilhas}

A Tabela~\ref{tab:comparativo-trilhas} resume, para cada carreira, o que continua igual e o
que muda. O padrão é o mesmo nas três: a base técnica continua necessária, porque é o que
permite revisar o que a IA produz, e o tempo liberado pela IA é cobrado em outras
habilidades, principalmente modelagem, revisão e negócio.

\begin{table}[htbp]
  \centering
  \caption{O que continua e o que muda em cada carreira com a IA.}
  \label{tab:comparativo-trilhas}
  \small
  \begin{tabularx}{\textwidth}{@{}>{\raggedright\arraybackslash}p{0.17\textwidth}
      >{\raggedright\arraybackslash}X>{\raggedright\arraybackslash}X@{}}
    \toprule
    Carreira & O que continua & O que muda com a IA \\
    \midrule
    Analista de dados & SQL sólido, visualização, estatística básica, comunicação &
      Menos tempo montando consultas e painéis; mais tempo formulando perguntas,
      conferindo resultados e conhecendo o negócio \\
    \addlinespace
    Analytics engineer & Modelagem dimensional, SQL, versionamento, testes &
      Testes escritos antes do modelo viram regra; documentação passa a ser insumo da IA;
      mais revisão de código gerado \\
    \addlinespace
    Engenheiro de dados & Fundamentos de sistemas, Python, SQL, nuvem, orquestração &
      Mais atenção a custo por consumo, segurança e dados para treino e contexto de IA;
      agentes executando parte da operação \\
    \bottomrule
  \end{tabularx}
  \fonte{elaborado pelo autor.}
\end{table}

Duas observações valem para as três carreiras. A primeira é que as fronteiras estão se
estreitando. O analista que aprende versionamento e testes, ou o engenheiro que aprende a
conversar com o negócio, ocupa um espaço que antes exigia duas pessoas. A IA facilita
justamente esse movimento, porque ajuda cada um nas partes que domina menos.

A segunda é que a ordem de estudo do capítulo~\ref{cap:trilha-hoje} vale para as três,
mudando só a profundidade de cada etapa. O analista vai mais fundo em estatística e
comunicação; o analytics engineer, em modelagem e testes; o engenheiro, em sistemas e
infraestrutura. Mas todos passam por SQL à mão, modelagem, Python, práticas de engenharia e
negócio.

\rascunho{
  Opcional: o que eu procuro hoje ao contratar para cada um desses papéis, e o que mudou
  em relação a alguns anos atrás.
}

\section{Portfólio}
\label{sec:portfolio}

Durante anos, um bom portfólio de dados mostrava domínio técnico: um painel bonito, um
modelo com boa acurácia, um pipeline funcionando. Hoje, qualquer pessoa gera algo
parecido com a ajuda da IA em poucas horas. Um portfólio que se destaca precisa mostrar o
que a IA não faz sozinha: o raciocínio.

Algumas recomendações práticas:

\begin{itemize}
  \item \textbf{Começar por uma pergunta, não por um conjunto de dados.} \enquote{Fiz uma
    análise do conjunto X} diz pouco. \enquote{Queria saber se Y, e o que encontrei foi Z}
    mostra como a pessoa pensa.
  \item \textbf{Usar dados públicos e reais.} Órgãos de estatística, bancos centrais,
    agências reguladoras e portais de dados abertos oferecem dados sujos e incompletos de
    verdade, que ensinam mais do que os conjuntos limpos de tutorial.
  \item \textbf{Mostrar as decisões.} Por que este grão, por que esta métrica, por que
    este gráfico, o que foi descartado e por quê. Um texto curto explicando as escolhas vale
    mais do que o código.
  \item \textbf{Ser reprodutível.} Código versionado, instruções para rodar do zero,
    testes que mostram que os números estão certos. É uma forma de mostrar práticas de
    engenharia sem precisar dizer que as conhece.
  \item \textbf{Ser honesto sobre a IA.} Dizer onde e como a IA foi usada, e o que foi
    revisado. Quem avalia sabe que ela foi usada; o que interessa é se a pessoa sabe
    avaliar o que ela fez.
  \item \textbf{Escrever.} Um texto que explica uma análise para um público leigo mostra
    comunicação e entendimento ao mesmo tempo. É a parte do portfólio que mais diferencia
    e a que menos gente faz.
\end{itemize}

Um projeto pequeno, com uma pergunta clara, dados reais, decisões explicadas e números
conferidos, mostra mais do que dez projetos grandes e genéricos. É, aliás, um bom resumo de
toda esta parte do livro: técnica suficiente para fazer e para revisar, fundamentos para
saber se o resultado faz sentido, e entendimento do problema para saber se ele é útil.

A próxima parte muda de assunto e de nível de detalhe. Sai da pergunta \enquote{o que
estudar} e entra na pergunta \enquote{como organizar os dados}, agora que eles precisam
servir a muito mais do que painéis.

\begin{leituras}
  \leitura{reis2022}{A referência mais completa e independente de ferramenta para quem
    segue em direção à engenharia de dados.}
  \leitura{kaminsky2019}{Para entender o papel do analytics engineer a partir do texto que
    o descreveu.}
  \leitura{provost2013}{Para analistas e cientistas de dados que querem fortalecer o
    pensamento analítico orientado ao negócio.}
\end{leituras}


\part{Modelar e governar dados}
\label{parte:modelagem}
% Capitulo: Arquitetura medalhão e modelagem tradicional
\chapter{Arquitetura medalhão e modelagem tradicional}
\label{cap:medalhao}

Esta parte do livro trata de como organizar os dados, agora que eles precisam servir a
muito mais do que painéis. O ponto de partida é uma arquitetura que se popularizou na era
do lakehouse e virou vocabulário comum em times de dados: a arquitetura medalhão, com suas
camadas bronze, prata e ouro. Este capítulo explica o que ela é, em que se parece e em que
difere da modelagem tradicional que o capítulo~\ref{cap:dw} apresentou, e por que ela é um
bom ponto de partida para pensar em vários usos do mesmo dado. A tese é simples: medalhão
organiza camadas de refinamento; não substitui modelagem.

\section{Bronze, prata e ouro}
\label{sec:bronze-prata-ouro}

O nome vem do ecossistema do lakehouse, em que o termo foi difundido para descrever uma
organização dos dados em camadas de qualidade crescente
\parencite{databricksmedallion,armbrust2021}. Cada camada tem um papel claro.

A camada \textbf{bronze} guarda o dado como chegou da fonte, sem transformação. Se o
sistema de origem mandou um arquivo com colunas a mais, datas em formato estranho e
registros duplicados, é assim que ele fica no bronze, junto com o momento em que chegou.
Parece desperdício guardar dado sujo, mas o bronze tem duas funções importantes: permite
reprocessar tudo quando uma regra de negócio muda, e permite auditar, mostrando exatamente
o que a fonte mandou e quando.

A camada \textbf{prata} guarda o dado limpo e integrado. Os tipos estão corretos, as
duplicidades foram tratadas, os nomes de colunas seguem um padrão, e as chaves de sistemas
diferentes foram conciliadas, de forma que o mesmo cliente tem o mesmo identificador
independentemente da fonte. A prata ainda não está organizada para um uso específico; é
uma base confiável a partir da qual vários usos podem ser construídos.

A camada \textbf{ouro} guarda o dado pronto para um consumo específico. Para o BI, modelos
dimensionais e métricas de negócio. Para a ciência de dados, tabelas de variáveis. Para
um assistente de IA, documentos preparados e índices de busca. Pode haver muitas tabelas
ouro construídas sobre a mesma prata, cada uma desenhada para quem vai usá-la.

Quem trabalha com transformação de dados como código vai reconhecer aqui as camadas de
preparação, intermediária e de consumo descritas no capítulo~\ref{cap:nuvem}
(\veja{sec:elt}). A correspondência não é exata, mas a ideia é a mesma: separar o trabalho
de limpar do trabalho de modelar para um uso.

\begin{figure}[htbp]
  \centering
  % Arquitetura medalhao e os consumidores de cada camada (cap. 12).
\begin{tikzpicture}[diagrama, x=1mm, y=1mm]
  % fontes
  \foreach \y/\nome/\id in {12/Sistemas/f1, 4/APIs/f2, -4/Arquivos/f3, -12/Documentos/f4}
    \node[caixa, minimum width=15mm, minimum height=6mm] (\id) at (7.5,\y) {\nome};
  \node[moldura, fit=(f1)(f4), label={[rotulo]above:Fontes}] {};

  % camadas
  \node[banco, minimum width=19mm, minimum height=26mm] (b) at (31,0)
    {\textbf{Bronze}\\[1mm]{\scriptsize bruto, como}\\{\scriptsize chegou}};
  \node[banco, minimum width=19mm, minimum height=26mm] (p) at (53,0)
    {\textbf{Prata}\\[1mm]{\scriptsize limpo, tipado,}\\{\scriptsize integrado}};
  \node[bancodestaque, minimum width=19mm, minimum height=26mm] (o) at (75,0)
    {\textbf{Ouro}\\[1mm]{\scriptsize pronto para}\\{\scriptsize o consumo}};
  \foreach \id in {f1,f2,f3,f4} \draw[seta] (\id.east) -- (b.west);
  \draw[seta] (b) -- (p);
  \draw[seta] (p) -- (o);

  % consumidores
  \foreach \y/\nome/\id in {15/{BI descritivo}/c1, 5/{Ciência de dados}/c2, -5/{Treino de IA}/c3, -15/{Contexto p/ IA}/c4}
    \node[caixa, minimum width=21mm, minimum height=6mm] (\id) at (104.5,\y) {\nome};
  \foreach \id in {c1,c2,c3,c4} \draw[seta] (o.east) -- (\id.west);
  \node[rotulo, text width=110mm] at (58,-23)
    {Ciência de dados e treino de IA muitas vezes leem também da prata, que guarda mais detalhe que o ouro.};
\end{tikzpicture}

  \caption{A arquitetura medalhão e os consumidores que leem de cada camada.}
  \label{fig:medalhao}
  \fonte{elaborado pelo autor com base em \textcite{databricksmedallion}.}
\end{figure}

A Figura~\ref{fig:medalhao} mostra as três camadas e os consumidores típicos. Repare que
nem todo consumo vem do ouro. Ciência de dados e treino de IA muitas vezes leem direto da
prata, porque precisam de detalhe que o ouro, agregado para o BI, já perdeu.

\section{O que muda em relação a Inmon e Kimball}
\label{sec:medalhao-tradicional}

É comum ver a arquitetura medalhão apresentada como alternativa a Inmon e Kimball. Não é.
As duas escolas do capítulo~\ref{cap:dw} falam de \emph{modelo}: como as tabelas devem ser
desenhadas. A arquitetura medalhão fala de \emph{camadas de refinamento}: em que estágio de
tratamento cada dado está. As duas coisas se combinam.

Na prática, a prata lembra muito a camada integrada de Inmon: um lugar onde os dados das
várias fontes são limpos e conciliados antes de servir a qualquer uso. E o ouro, quando
serve ao BI, quase sempre é um modelo dimensional de Kimball, com fatos, dimensões e grão
declarado (\veja{sec:modelagem-dimensional}). O bronze é a novidade mais clara em relação
aos armazéns dos anos 1990, que raramente guardavam o dado bruto, porque armazenamento era
caro. Com armazenamento barato, guardar o bruto deixou de ser luxo e passou a ser
recomendável.

O risco está em tratar a arquitetura medalhão como receita completa. Três camadas não
garantem nada sobre a qualidade do desenho. É possível ter bronze, prata e ouro e, ainda
assim, um ouro sem grão definido, com a mesma métrica calculada de três formas em três
tabelas, e uma prata em que a conciliação de clientes nunca foi feita. Nesse caso, o
resultado é a mesma bagunça de antes, só que organizada em pastas com nomes de metais.

\begin{destaque}
A arquitetura medalhão diz em que estágio cada dado está. A modelagem diz como ele deve ser
desenhado. Uma não substitui a outra.
\end{destaque}

\rascunho{
  Relato: como eu organizei as camadas num projeto real (o que ficou em cada camada, o que
  foi difícil de decidir, que regra de nomenclatura funcionou). Sem nomes de empresa.
}

\section{Modelagem para análise descritiva}
\label{sec:modelagem-descritiva}

A análise descritiva responde ao que aconteceu e por quê. É o uso que a modelagem
tradicional sempre atendeu, e continua sendo, de longe, o maior volume de trabalho na
maioria das empresas: acompanhamento de vendas, custos, operação, metas.

Na camada ouro, a modelagem para análise descritiva segue os princípios do
capítulo~\ref{cap:dw}. Esquemas estrela, com o grão declarado e as dimensões conformadas.
Métricas definidas uma vez, preferencialmente numa camada semântica, e reaproveitadas por
todos os painéis (\veja{sec:camada-semantica}). Tabelas agregadas para os painéis mais
acessados, quando o custo de calcular na hora não compensa.

Algumas práticas da tradição valem especialmente aqui. Chaves substitutas protegem o modelo
de mudanças nas fontes. O histórico das dimensões, guardado com datas de validade, permite
responder como as coisas eram em qualquer momento do passado. E nomes claros, que um
analista novo entende sem perguntar, são talvez o investimento mais barato e mais
subestimado de todos. Eles ajudam o analista novo e, como mostra o
capítulo~\ref{cap:contexto-ia}, ajudam também a IA.

\section{Um dado, vários usos}
\label{sec:um-dado-varios-usos}

A modelagem tradicional tinha um consumidor principal: o analista, usando uma ferramenta de
BI. Hoje, o mesmo dado alimenta modelos preditivos, decisões automatizadas, o treino e o
ajuste de modelos de IA, e assistentes que respondem perguntas em linguagem natural. Cada
um desses usos pede uma forma diferente de organizar os dados, e um dos erros mais comuns é
tentar servir a todos com o mesmo modelo, geralmente o que foi feito para o BI.

A Tabela~\ref{tab:usos-modelagem} resume as diferenças e serve de mapa para os próximos
capítulos desta parte. O capítulo~\ref{cap:master-data} trata da base que todos os usos
compartilham, os cadastros mestres. Os capítulos~\ref{cap:governanca}
e~\ref{cap:aplicacao-propria} tratam de como cuidar dessa base. E os
capítulos~\ref{cap:modelagem-cd}, \ref{cap:treinar-ia} e~\ref{cap:contexto-ia} tratam de
um uso cada um.

\begin{table}[htbp]
  \centering
  \caption{Cada uso pede uma forma diferente de organizar os dados.}
  \label{tab:usos-modelagem}
  \small
  \begin{tabularx}{\textwidth}{@{}>{\raggedright\arraybackslash}p{0.22\textwidth}
      >{\raggedright\arraybackslash}X>{\raggedright\arraybackslash}X@{}}
    \toprule
    Uso & Forma dos dados & Cuidado principal \\
    \midrule
    Descritiva \newline{\footnotesize\itshape o que aconteceu?} & Esquema estrela e tabelas agregadas, na camada ouro & Métrica definida uma vez; grão explícito \\
    \addlinespace
    Preditiva \newline{\footnotesize\itshape o que vai acontecer?} & Tabela larga: uma linha por entidade e momento, com variáveis e alvo & Não deixar o futuro vazar para o passado \\
    \addlinespace
    Prescritiva \newline{\footnotesize\itshape o que fazer?} & Previsões junto com custos, restrições e cenários do negócio & Registrar a decisão tomada e o resultado \\
    \addlinespace
    Treino de IA & Conjunto curado, rotulado, versionado e documentado & Qualidade e representatividade dos exemplos, privacidade \\
    \addlinespace
    Contexto para IA & Documentos em trechos, metadados, glossário, camada semântica & Atualização, dono de cada fonte, permissões \\
    \bottomrule
  \end{tabularx}
  \fonte{elaborado pelo autor.}
\end{table}

A boa notícia é que a arquitetura em camadas facilita essa diversidade. Com uma prata
confiável, construir um ouro novo para um uso novo é muito mais barato do que voltar às
fontes. O investimento que mais rende, portanto, é na prata: limpeza, integração e,
principalmente, a conciliação dos cadastros, que é o assunto do próximo capítulo.

\begin{leituras}
  \leitura{armbrust2021}{O artigo que propõe o lakehouse, o contexto em que a
    arquitetura medalhão se popularizou.}
  \leitura{kleppmann2017}{Para entender os sistemas por baixo das camadas:
    armazenamento, replicação, lote e fluxo contínuo, sem depender de ferramenta.}
  \leitura{kimball2013}{Para a modelagem da camada ouro voltada ao BI, que continua sendo a
    dimensional.}
\end{leituras}

% Capitulo: Master data: de projeto eterno a oportunidade
\chapter{Master data: de projeto eterno a oportunidade}
\label{cap:master-data}

Toda empresa tem algumas entidades que aparecem em quase todos os processos: clientes,
produtos, fornecedores, locais, funcionários. Os dados que descrevem essas entidades são
chamados de dados mestres, ou master data. Quando eles estão bem cuidados, quase ninguém
repara. Quando não estão, tudo dá errado: a receita por cliente não bate, o mesmo produto
aparece com três códigos, o relatório de compras não conversa com o de estoque.

Por muito tempo, cuidar disso foi visto como um projeto caro, lento e difícil de
justificar, e muitas empresas preferiram conviver com as planilhas. Este capítulo
argumenta que essa conta mudou. Com a abordagem certa, e com a ajuda da IA, master data
pode deixar de ser um projeto eterno e virar uma oportunidade.

\section{O que é master data}
\label{sec:o-que-e-master-data}

Dados mestres são as entidades centrais do negócio, compartilhadas por vários processos e
sistemas \parencite{dama2017,loshin2009}. Eles se distinguem de dois outros tipos de dado.
O dado transacional registra eventos: a venda, a viagem, o lançamento contábil. Ele é
volumoso, cresce todo dia e quase nunca muda depois de registrado. O dado de referência
são listas de valores usadas para classificar outros dados: estados, moedas, códigos de
unidade de medida. Ele é pequeno, muda pouco e muitas vezes vem de fora da empresa.

O dado mestre fica no meio. É menos volumoso que o transacional, mas muda com frequência
(o cliente muda de endereço, o produto ganha uma nova embalagem) e é usado por quase todas
as transações. E, principalmente, costuma existir em vários sistemas ao mesmo tempo. O
ERP tem o seu cadastro de clientes, o CRM tem outro, a área comercial mantém uma planilha
com um terceiro. Cada um com os seus códigos e a sua forma de escrever os nomes
(Tabela~\ref{tab:registro-mestre}).

\begin{table}[htbp]
  \centering
  \caption{O mesmo cliente em três fontes e o registro mestre que as concilia.}
  \label{tab:registro-mestre}
  \small
  \begin{tabular}{@{}lll@{}}
    \toprule
    Fonte & ID na fonte & Nome como está na fonte \\
    \midrule
    ERP                & 4471          & AGRO SILVA LTDA \\
    CRM                & C-0932        & Agro Silva \\
    Planilha comercial & (sem ID)      & Agrosilva Faz. Boa Vista \\
    \midrule
    \textbf{Registro mestre} & \textbf{1001} & \textbf{Agro Silva Ltda} \\
    \bottomrule
  \end{tabular}
  \fonte{exemplo fictício elaborado pelo autor.}
\end{table}

Gerir master data é, no essencial, manter duas coisas. A primeira é o registro mestre
(\emph{golden record}): a versão única e confiável de cada entidade, com um identificador
próprio. A segunda é a tabela de correspondência, que liga o identificador de cada fonte ao
identificador mestre. Com as duas, qualquer análise pode juntar dados de sistemas
diferentes sabendo que está falando do mesmo cliente.

Quem leu o capítulo~\ref{cap:dw} já viu esse problema com outro nome. As dimensões
conformadas de Kimball, compartilhadas por todos os processos de negócio, são master data
aplicado ao data warehouse (\veja{sec:kimball}). Sem um cadastro mestre por trás, não há
dimensão conformada que se sustente.

\section{Como master data era visto}
\label{sec:master-data-antes}

O jeito tradicional de resolver master data era um programa grande, separado de qualquer
entrega de negócio. Primeiro se mapeavam todas as fontes de todos os domínios. Depois se
conciliavam os nomes e os códigos, registro por registro. Depois se definia a governança:
quem cria, quem altera, quem aprova. Só então havia alguma coisa para entregar.

Esse caminho exigia gente e tempo. Era preciso contratar analistas de dados para
conversar com os analistas de negócio de cada área, entender como cada sistema registrava
os dados, negociar regras que nenhuma área queria ceder. Até que o resultado aparecesse,
passavam meses, muitas vezes anos (parte de cima da Figura~\ref{fig:master-data-abordagens}).
E o resultado, quando aparecia, era difícil de mostrar: \enquote{dados mais consistentes}
não é algo que a diretoria vê numa reunião.

Diante disso, muitas empresas fizeram uma escolha racional: continuar com as planilhas. A
planilha com o \enquote{de-para} mantido à mão já existia, funcionava mais ou menos e não
pedia orçamento. O custo dela era real, mas espalhado: retrabalho a cada fechamento, cada
área com a sua versão, números que não batiam e consumiam a primeira meia hora de cada
reunião. Ninguém via esse custo numa linha do orçamento, e por isso ele nunca competia com
o custo visível de um programa de master data.

\rascunho{
  Relato: uma situação em que a falta de master data causou um problema concreto
  (números divergentes, retrabalho no fechamento, decisão tomada com base errada). Sem
  nomes e com valores só em ordem de grandeza.
}

\section{Master data puxado por um projeto de negócio}
\label{sec:master-data-puxado}

A abordagem que tem funcionado melhor na minha experiência inverte a lógica. Em vez de
começar um programa de master data e procurar quem vai usá-lo, começa-se por um projeto de
negócio que já foi aprovado, já tem patrocinador e orçamento, e identificam-se os braços
desse projeto que dependem de master data.

Um exemplo típico: um projeto comercial precisa juntar a carteira de clientes de três
sistemas para acompanhar o desempenho de cada vendedor. Sem um cadastro mestre de clientes,
o projeto não fica de pé, porque o mesmo cliente aparece três vezes com nomes diferentes. É
esse o momento de abrir, em paralelo, um projeto de master data só para o domínio de
clientes. A necessidade já está justificada, porque sem ele o projeto de negócio não
entrega. O recurso já está garantido, porque faz parte do projeto aprovado. E o escopo é
pequeno, porque é um domínio só.

O projeto de negócio vira o primeiro cliente do master data, e é ele que valida o
resultado na prática: se os vendedores reconhecem as próprias carteiras no painel, a
conciliação funcionou. Quando o próximo projeto de negócio chegar e precisar de clientes,
o domínio já está pronto, e ele só precisa puxar o domínio que falta, digamos, produtos. A
cada projeto, um domínio novo; a cada domínio, o próximo projeto fica mais barato (parte de
baixo da Figura~\ref{fig:master-data-abordagens}).

\begin{figure}[htbp]
  \centering
  % Master data: programa unico contra dominios puxados por projetos de negocio (cap. 13).
\begin{tikzpicture}[diagrama, x=1mm, y=1mm,
    titulo/.style={anchor=west, font=\footnotesize\scshape, text=cinza},
    linha/.style={anchor=west, font=\footnotesize, text=cinza}]

  % --- antes
  \node[titulo] at (0,40) {antes: um programa de master data};
  \node[caixa, anchor=west, minimum width=26mm] at (20,31) {Mapear fontes};
  \node[caixa, anchor=west, minimum width=34mm] at (46,31) {Conciliar nomes e IDs};
  \node[caixa, anchor=west, minimum width=20mm] at (80,31) {Governança};
  \node[destaquecaixa, anchor=west, minimum width=16mm] at (100,31) {Resultado};
  \draw[seta, draw=cinza] (20,24.5) -- (116,24.5);
  \node[rotulo, anchor=north] at (68,24) {meses ou anos até o primeiro resultado};

  % --- agora
  \node[titulo] at (0,12) {agora: puxado por um projeto de negócio};
  \node[linha] at (0,3) {Negócio};
  \node[linha] at (0,-9) {Master data};
  \node[caixa, anchor=west, minimum width=44mm] (pa) at (20,3) {Projeto de negócio A};
  \node[caixa, anchor=west, minimum width=46mm] (pb) at (70,3) {Projeto de negócio B};
  \node[destaquecaixa, anchor=west, minimum width=28mm] (cl) at (26,-9) {Clientes};
  \node[destaquecaixa, anchor=west, minimum width=28mm] (pr) at (80,-9) {Produtos};
  \draw[seta] (cl.north) -- (cl.north |- pa.south);
  \draw[seta] (pr.north) -- (pr.north |- pb.south);
  \draw[seta, dashed, draw=cinza] (cl.east) -- ($(pb.south west)+(4,0)$)
    node[rotulo, pos=0.45, below right] {reuso};
  \draw[seta, draw=cinza] (20,-16.5) -- (116,-16.5);
  \node[rotulo, anchor=north] at (68,-17)
    {um domínio por vez, cada um justificado por um projeto que precisa dele};
\end{tikzpicture}

  \caption{Dois jeitos de construir master data: um programa único, que entrega no
    fim, ou um domínio por vez, puxado por projetos de negócio que já precisam dele.}
  \label{fig:master-data-abordagens}
  \fonte{elaborado pelo autor.}
\end{figure}

A abordagem tem cuidados próprios. O primeiro é não deixar o prazo do projeto de negócio
atropelar a qualidade do domínio: um cadastro mestre feito às pressas, só para destravar
uma entrega, vira mais uma versão da verdade. O segundo é desenhar o domínio pensando no
reuso, e não só no primeiro projeto. O cadastro de clientes que nasce para o projeto
comercial vai ser usado pelo financeiro e pela logística, e precisa ter os atributos e as
regras que eles vão precisar. O terceiro é dar visibilidade ao projeto paralelo: ele entra
no quadro do time como um épico próprio, com histórias amarradas às do projeto de negócio
(\veja{sec:projeto-sustentacao}), para que o trabalho não fique escondido.

\rascunho{
  Relato completo de um caso real dessa abordagem: o projeto de negócio, o domínio que
  ele puxou, quanto tempo levou, quem validou e o que foi reaproveitado depois. É o relato
  mais importante do capítulo.
}

\section{O que a IA muda na conciliação}
\label{sec:master-data-ia}

A etapa que fazia o master data levar meses era a conciliação: decidir, registro por
registro, se dois cadastros representam a mesma entidade. É um problema com teoria antiga.
Em 1969, Ivan Fellegi e Alan Sunter formalizaram a conciliação probabilística de registros,
que compara os campos de cada par e calcula a probabilidade de os dois serem o mesmo
\parencite{fellegi1969}. As técnicas que vieram depois, de normalização, comparação e
bloqueio, estão reunidas em \textcite{christen2012}, e há bibliotecas de código aberto
que implementam esse modelo em escala \parencite{linacre2022}.

A IA acrescenta a essas técnicas a capacidade de lidar com o que as regras não pegam: nomes
abreviados, razões sociais escritas de jeitos diferentes, endereços incompletos, a filial
que aparece com o nome da fazenda. Ela pode sugerir pares prováveis, padronizar textos e,
principalmente, explicar por que dois registros parecem ser o mesmo, o que acelera muito
a revisão humana. A pessoa continua decidindo; a IA tira o trabalho braçal.

\rascunho{
  Validar com a minha experiência: em que medida a IA de fato encurtou a conciliação nos
  projetos de master data que eu acompanhei, e onde ela errou.
}

O efeito sobre a abordagem por domínio é direto. Conciliar um domínio inteiro dentro do
prazo de um projeto de negócio ficou viável, o que antes exigiria um time dedicado por
meses. Um motor simples de conciliação, com regras, similaridade e sugestões da IA, está
descrito no capítulo~\ref{cap:aplicacao-propria} (\veja{sec:motor-conciliacao}).

Há um cuidado que não pode ser esquecido. Cadastro de cliente tem dado pessoal, e o que vai
para um modelo de IA externo sai do ambiente da empresa. A limpeza do que pode e do que não
pode ser enviado precisa acontecer antes (\veja{sec:seguranca}).

\section{Governança mínima que funciona}
\label{sec:governanca-master-data}

Um domínio de master data sem governança volta a se degradar em poucos meses: alguém cria
um cliente novo direto no sistema de origem, outro corrige o nome na planilha, e a
correspondência deixa de valer. A governança não precisa ser grande, mas precisa existir
desde o primeiro domínio.

\begin{itemize}
  \item \textbf{Um dono de negócio por domínio}, que decide as regras: o que é um cliente
    ativo, quando dois cadastros são o mesmo, quais atributos são obrigatórios.
  \item \textbf{Regras de criação e alteração}: quem pode criar um cliente novo e em qual
    sistema, e como essa criação chega ao cadastro mestre.
  \item \textbf{Testes automáticos} de unicidade e de correspondência, rodando a cada carga
    (\veja{sec:tdd}).
  \item \textbf{A tabela de correspondência versionada e auditável}, num banco de dados e
    não numa planilha, com o registro de quem decidiu cada vínculo.
  \item \textbf{Crescimento junto com os domínios}: governança do tamanho do que já foi
    entregue, ampliada a cada domínio novo.
\end{itemize}

A governança completa, com o papel do data steward na era da IA, é o assunto do próximo
capítulo. Como construir uma aplicação própria para tudo isso, do motor de conciliação à
fila de revisão, fica para o capítulo~\ref{cap:aplicacao-propria}.

\begin{leituras}
  \leitura{dama2017}{O guia de referência de gestão de dados. O capítulo sobre dados
    mestres e de referência dá o vocabulário e as responsabilidades.}
  \leitura{loshin2009}{Um livro inteiro sobre master data, do caso de negócio à
    arquitetura e à governança.}
  \leitura{christen2012}{Para quem vai implementar a conciliação: técnicas de
    comparação, bloqueio e avaliação da qualidade dos pares encontrados.}
  \leitura{fellegi1969}{O artigo clássico que formalizou a conciliação probabilística
    de registros, ainda base de muitas ferramentas.}
\end{leituras}

% Capitulo: Governança de dados na era da IA
\chapter{Governança de dados na era da IA}
\label{cap:governanca}

Governança de dados sempre teve fama de burocracia: comitês, políticas longas,
dicionários de dados em planilha e um responsável cobrando todo mundo por cadastros mal
preenchidos. Não por acaso, era a primeira coisa cortada quando um projeto apertava. Este
capítulo argumenta que a IA muda essa equação. Boa parte do trabalho braçal da governança
(documentar, classificar, conciliar, avisar) pode ser proposta pela máquina, e o papel de
quem cuida dos dados passa de executor a tutor e revisor. E construir uma ferramenta interna
que faça parte desse trabalho ficou ao alcance de qualquer time.

\section{O que é governança de dados}
\label{sec:gov-o-que-e}

Governança de dados é o conjunto de papéis, regras e processos que garantem que os dados
são confiáveis, bem descritos, protegidos e usados de forma adequada \parencite{dama2017}.
Ela abrange várias frentes que já apareceram neste livro: o catálogo e os metadados, que
dizem o que cada dado é; a qualidade, que diz se ele está certo para o uso
(\veja{sec:fund-qualidade}); o master data (capítulo~\ref{cap:master-data}); o acesso e
a privacidade; a linhagem, que mostra de onde cada dado veio; e o ciclo de vida, que define
quanto tempo cada dado fica guardado.

Os papéis clássicos são três, embora a nomenclatura varie entre autores e empresas. O
\emph{dono do dado} é alguém do negócio que responde por um domínio e decide as suas regras.
O \emph{data steward} cuida do dia a dia: a qualidade, a documentação, a resolução de
problemas. O \emph{custodiante} cuida da parte técnica: onde o dado fica, como é protegido,
como é copiado. \pendente{conferir a nomenclatura dos papéis no DMBOK}
\textcite{eryurek2021} mostram como esses papéis se reorganizam quando os dados vão para a
nuvem, com muito mais automação e muito mais gente com acesso.

\section{Como a governança era vista}
\label{sec:gov-antes}

Na forma tradicional, governança era um programa corporativo com comitês, políticas
extensas e um conjunto de documentos que ninguém lia. O dicionário de dados ficava numa
planilha, desatualizado no dia seguinte à publicação. E o data steward fazia papel de
fiscal: mantinha tudo à mão, cobrava dos outros e estava sempre atrás dos dados, que mudavam
mais rápido do que ele conseguia documentar.

O resultado era previsível. A governança virava sinônimo de atraso, as áreas aprendiam a
contorná-la, e, no primeiro aperto de prazo, era cortada. Não por má vontade, mas porque o
custo era visível e imediato, e o benefício, difuso e futuro, como acontecia com o master
data (\veja{sec:master-data-antes}).

Uma das críticas mais úteis a esse modelo é a da governança não invasiva. Em vez de criar
uma estrutura nova, com cargos e comitês novos, a proposta é reconhecer e formalizar as
responsabilidades que as pessoas já exercem sobre os dados no trabalho do dia a dia
\parencite{seiner2014}. Quem cadastra clientes já é, na prática, responsável pela qualidade
desse cadastro. A governança deixa de ser algo imposto de fora e passa a ser a organização
do que já acontece.

\section{O data steward como tutor e revisor da IA}
\label{sec:steward-tutor}

A IA não muda o que a governança precisa fazer. Muda quem faz o trabalho braçal. Uma lista
do que ela já pode propor, com boa qualidade, num ambiente de dados:

\begin{itemize}
  \item descrições de tabelas e colunas a partir dos nomes, dos dados e do código que as gera;
  \item classificação de dados sensíveis, como dado pessoal ou financeiro, para orientar
    permissões;
  \item pares suspeitos de duplicidade em cadastros mestres;
  \item inconsistências entre fontes, como o mesmo produto com unidades diferentes;
  \item regras de qualidade a partir do comportamento dos dados, como faixas de valores
    esperadas.
\end{itemize}

Em todos esses casos, a IA propõe e uma pessoa decide. O data steward deixa de escrever
cada descrição e passa a revisar descrições propostas; deixa de procurar duplicidades e
passa a aprovar ou rejeitar pares sugeridos. E, a cada decisão, ensina: a correção vira
exemplo ou regra, e a próxima proposta vem melhor (Figura~\ref{fig:governanca-ciclo}). O
steward passa a ser tutor e revisor da IA.

\begin{figure}[htbp]
  \centering
  % Ciclo de governanca com o data steward como tutor e revisor da IA (cap. 14).
\begin{tikzpicture}[diagrama, x=1mm, y=1mm,
    etapa/.style={caixa, text width=34mm, minimum height=13mm, font=\scriptsize}]
  \node[destaquecaixa, etapa] (prop) at (55,28)
    {\textbf{A IA propõe}\\descrições, classificação de dado sensível, pares suspeitos};
  \node[etapa] (rev) at (96,10)
    {\textbf{O steward revisa}\\aprova, corrige ou rejeita};
  \node[etapa] (reg) at (55,-8)
    {\textbf{A decisão é registrada}\\auditoria, correspondência, catálogo};
  \node[etapa] (apr) at (14,10)
    {\textbf{Regras e exemplos melhoram}\\cada correção ensina a próxima proposta};
  \draw[seta] (prop.east) to[bend left=20] (rev.north);
  \draw[seta] (rev.south) to[bend left=20] (reg.east);
  \draw[seta] (reg.west) to[bend left=20] (apr.south);
  \draw[seta] (apr.north) to[bend left=20] (prop.west);
  \node[rotulo] at (55,10) {o steward como\\tutor e revisor};
\end{tikzpicture}

  \caption{O ciclo de governança com a IA propondo e o data steward revisando. Cada
    decisão registrada melhora a próxima proposta.}
  \label{fig:governanca-ciclo}
  \fonte{elaborado pelo autor.}
\end{figure}

É o mesmo movimento que aparece em outros pontos do livro. Testes escritos antes do modelo
eram recomendados e raramente feitos, e a IA os tornou baratos o bastante para virar regra
(\veja{sec:tdd}). A documentação era boa prática que sempre ficava para depois, e virou o
insumo principal dos assistentes que respondem perguntas sobre os dados
(\veja{sec:metadados-contexto}). Com a governança acontece o mesmo: o que antes não cabia no
tempo de ninguém agora cabe, desde que alguém revise.

A mudança tem uma consequência para a carreira de quem cuida dos dados. O valor do data
steward deixa de estar no volume de trabalho que ele executa e passa a estar no julgamento
que ele aplica. Conhecer o negócio, saber quando duas coisas parecidas são diferentes,
perceber que uma descrição tecnicamente correta vai confundir quem lê: é isso que a IA não
faz, e é isso que o steward ensina a ela.

\rascunho{
  Relato: onde eu já vi ou apliquei esse modelo de steward como revisor da IA (por exemplo,
  descrições de tabelas geradas e revisadas, ou classificação de dados sensíveis), e o que
  foi preciso corrigir. Sem nomes.
}

\section{Uma aplicação interna ao alcance de qualquer time}
\label{sec:gov-aplicacao-interna}

Um exemplo concreto mostra como o cenário mudou. Imagine um programa simples, em Python, de
uso interno, que roda depois de cada carga de dados. Ele procura, nas dimensões do armazém,
valores diferentes que provavelmente representam a mesma coisa: o mesmo cliente escrito de
três jeitos, o mesmo produto com duas descrições. Quando encontra um caso, registra numa
fila de revisão e manda uma mensagem para o dono do domínio, com o link para decidir.

Há poucos anos, isso seria um projeto, com especificação, orçamento e meses de
desenvolvimento. Hoje, alguém que programa em Python monta uma primeira versão em poucos
dias, com a IA ajudando a escrever o código. O que continua difícil não é o código: é
combinar com o negócio quem decide, o que conta como \enquote{o mesmo cliente} e o que
acontece depois da decisão. O capítulo~\ref{cap:aplicacao-propria} mostra o que considerar
para levar uma aplicação dessas a sério, com um trecho de código do motor de conciliação
(\veja{sec:motor-conciliacao}).

\section{Ferramentas de mercado}
\label{sec:gov-ferramentas}

Construir não é a única opção. O mercado de governança, catálogo e master data é grande, e
a Tabela~\ref{tab:ferramentas-governanca} organiza as alternativas por categoria. É um
retrato de 2026: nomes, donos e preços mudam rápido, e por isso o que importa são as
categorias, não os nomes.

\begin{table}[htbp]
  \centering
  \caption{Ferramentas de governança, catálogo e master data por categoria.}
  \label{tab:ferramentas-governanca}
  \small
  \begin{tabularx}{\textwidth}{@{}>{\raggedright\arraybackslash}p{0.25\textwidth}
      >{\raggedright\arraybackslash}X>{\raggedright\arraybackslash}X@{}}
    \toprule
    Categoria & Exemplos & Quando faz sentido \\
    \midrule
    Catálogo e governança independentes & Alation, Collibra, Atlan, Informatica (Salesforce) & Muitas plataformas de dados, setor regulado, programa corporativo \\
    \addlinespace
    Nativas da plataforma & Databricks Unity Catalog, Snowflake Horizon, Microsoft Purview & Dados concentrados numa plataforma só \\
    \addlinespace
    Código aberto & DataHub, OpenMetadata & Time com engenharia para hospedar e manter \\
    \addlinespace
    Master data & Informatica MDM (Salesforce), Reltio (SAP), Profisee, Semarchy, Stibo Systems & Vários domínios, muitas fontes, regras complexas de conciliação \\
    \addlinespace
    Aplicação própria & Python, um banco de dados e um canal de mensagens & Começar pequeno, um domínio por vez, regras muito específicas do negócio \\
    \bottomrule
  \end{tabularx}
  \fonte{elaborado pelo autor com base nos sites dos fornecedores, em novembro de 2026.
    \pendente{conferir a tabela antes da impressão}}
\end{table}

O mercado também se movimenta. Nos últimos anos, dois dos maiores fornecedores
independentes de master data foram comprados por gigantes de software corporativo: a
Informatica pela Salesforce e a Reltio pela SAP. \pendente{conferir datas e citar fontes
oficiais antes da impressão} Para quem escolhe uma ferramenta, isso reforça a importância
de avaliar não só os recursos, mas também a facilidade de tirar os dados de lá se um dia
for preciso.

Uma orientação geral ajuda a navegar a tabela. As ferramentas nativas da plataforma
resolvem boa parte dos casos de quem concentra os dados num único armazém, com pouco custo
adicional. Os catálogos independentes fazem mais sentido quando há muitas plataformas, ou em
setores regulados que exigem trilhas de auditoria formais. As opções de código aberto são
boas para quem tem engenharia para hospedar e manter. E nada impede começar com uma
aplicação própria, pequena, num domínio, e migrar quando a escala pedir.

\section{O que não se delega à IA}
\label{sec:gov-nao-delegar}

Com tanto trabalho podendo ser proposto pela máquina, vale deixar claro o que continua com
as pessoas, por definição.

\begin{itemize}
  \item \textbf{A responsabilidade por cada domínio.} Alguém precisa responder pelo cadastro
    de clientes, e esse alguém não pode ser um modelo.
  \item \textbf{As regras de negócio.} O que é um cliente ativo, quando dois cadastros
    representam a mesma empresa, qual atributo é obrigatório. A IA pode sugerir; o negócio
    decide.
  \item \textbf{O acesso e as obrigações legais.} Quem pode ver o quê, e como a empresa
    cumpre a legislação de proteção de dados pessoais \parencite{lgpd2018}.
  \item \textbf{A decisão final em caso de dúvida.} A IA propõe, uma pessoa decide e assina.
\end{itemize}

A governança na era da IA, portanto, não é menor. É diferente: menos tempo executando, mais
tempo decidindo e ensinando. E, pela primeira vez, ela tem chance de andar na mesma
velocidade que os dados.

\begin{leituras}
  \leitura{eryurek2021}{Uma visão atual e prática de governança, com atenção a dados na
    nuvem, qualidade, privacidade e ao papel de cada pessoa.}
  \leitura{seiner2014}{A ideia de governança que formaliza o que as pessoas já fazem, em
    vez de criar uma burocracia nova. Combina bem com o steward como revisor.}
  \leitura{dama2017}{O guia de referência para o vocabulário e para as áreas da gestão de
    dados.}
\end{leituras}

% Capitulo: Construir a própria aplicação de master data e governança
\chapter{Construir a própria aplicação de master data e governança}
\label{cap:aplicacao-propria}

Há muito material sobre ferramentas de mercado para master data e governança, e muito
pouco sobre o que considerar quando uma empresa decide construir a sua. Até pouco tempo,
fazia sentido: construir era caro demais para a maioria. Com a IA ajudando a escrever
código, isso mudou, e mais times vão se perguntar se vale a pena. Este capítulo é um guia
para essa decisão e para a construção. Não é um tutorial de uma tecnologia específica; é a
lista de decisões que vale tomar antes de escrever a primeira linha de código.

\section{Construir ou comprar}
\label{sec:construir-comprar}

A IA baixou muito o custo de construir. Não baixou o custo de manter. Uma aplicação interna
precisa de alguém que corrija erros, atualize dependências, atenda dúvidas e adapte o
sistema quando o negócio mudar. A pergunta que decide entre construir e comprar, portanto,
não é \enquote{conseguimos fazer?}, e sim \enquote{quem vai cuidar disso daqui a dois
anos?}. A Tabela~\ref{tab:construir-comprar} reúne os outros critérios.

\begin{table}[htbp]
  \centering
  \caption{Quando construir e quando comprar uma solução de master data e governança.}
  \label{tab:construir-comprar}
  \small
  \begin{tabularx}{\textwidth}{@{}>{\raggedright\arraybackslash}p{0.17\textwidth}
      >{\raggedright\arraybackslash}X>{\raggedright\arraybackslash}X@{}}
    \toprule
    Critério & Construir faz sentido quando & Comprar faz sentido quando \\
    \midrule
    Escopo & Um ou dois domínios, com regras muito próprias do negócio & Muitos domínios e fontes, com regras comuns de mercado \\
    \addlinespace
    Equipe & Há quem programe e quem mantenha depois & Não há quem mantenha código interno \\
    \addlinespace
    Regulação & Um histórico de mudanças simples basta & Auditoria formal, certificações, setor regulado \\
    \addlinespace
    Prazo & É preciso algo em semanas, um domínio por vez & É um programa corporativo de longo prazo \\
    \addlinespace
    Custo & A licença é desproporcional ao problema & Manter código próprio custaria mais que a licença \\
    \addlinespace
    Integração & É preciso encaixar em fluxos internos específicos & A ferramenta já conecta com toda a pilha de dados \\
    \bottomrule
  \end{tabularx}
  \fonte{elaborado pelo autor.}
\end{table}

Construir também não precisa ser para sempre. Começar com uma aplicação própria num
domínio, aprender com ela o que a empresa realmente precisa e migrar para uma ferramenta de
mercado quando a escala pedir é um caminho legítimo. A condição é desenhar a aplicação de
forma que os dados saiam fácil: registro mestre, tabela de correspondência e histórico de
decisões em tabelas simples, num banco de dados comum, e não presos à lógica da aplicação.

\section{Começar pelo problema, não pela plataforma}
\label{sec:comecar-pequeno}

O erro mais comum de quem decide construir é começar pela plataforma: um sistema genérico de
master data, que serviria para qualquer domínio, com telas de configuração e regras
parametrizáveis. Esse caminho reproduz o programa eterno que o capítulo~\ref{cap:master-data}
criticou, agora com código próprio.

O ponto de partida melhor é um problema concreto, com um domínio, uma regra e um dono.
Algo como: \enquote{clientes com nomes diferentes no ERP e no CRM, revisados pelo time
comercial, para o projeto de acompanhamento de carteira}. A plataforma, se vier, nasce do
segundo e do terceiro problema, quando os padrões começam a se repetir.

O primeiro usuário da aplicação é o projeto de negócio que puxou o domínio
(\veja{sec:master-data-puxado}). É ele quem diz se a aplicação está resolvendo, e é a
opinião dele que deve orientar as primeiras melhorias, não uma lista de recursos imaginada
de antemão.

\section{Uma arquitetura mínima}
\label{sec:arquitetura-minima}

Uma aplicação de master data e governança, mesmo pequena, tem algumas peças que não podem
faltar (Figura~\ref{fig:aplicacao-propria}).

\begin{figure}[htbp]
  \centering
  % Arquitetura minima de uma aplicacao propria de master data e governanca (cap. 15).
\begin{tikzpicture}[diagrama, x=1mm, y=1mm]
  % fontes e ingestao
  \foreach \y/\nome/\id in {8/ERP/f1, 0/CRM/f2, -8/Planilhas/f3}
    \node[caixa, minimum width=14mm, minimum height=6mm] (\id) at (7,\y) {\nome};
  \node[moldura, fit=(f1)(f3), label={[rotulo]above:Fontes}] {};
  \node[caixa, minimum width=17mm, minimum height=10mm] (ing) at (29,0) {Ingestão\\{\scriptsize (prata)}};
  \foreach \id in {f1,f2,f3} \draw[seta] (\id.east) -- (ing.west);

  % motor de conciliacao
  \node[destaquecaixa, minimum width=24mm, minimum height=18mm] (mot) at (54,0)
    {Motor de\\conciliação\\[0.5mm]{\scriptsize regras,}\\{\scriptsize similaridade, IA}};
  \draw[seta] (ing) -- (mot);

  % saidas
  \node[bancodestaque, minimum width=24mm, minimum height=13mm] (reg) at (102,9)
    {\scriptsize Registro mestre e\\\scriptsize correspondência};
  \node[caixa, minimum width=20mm, minimum height=9mm] (fila) at (81,-10) {Fila de\\revisão};
  \node[caixa, minimum width=18mm, minimum height=9mm] (stew) at (104,-10) {Steward\\revisa};
  \node[caixa, draw=cinza, dashed, text=cinza, minimum width=20mm, minimum height=8mm] (notif) at (81,-24)
    {\scriptsize Mensagem\\\scriptsize aos donos};
  \draw[seta] (mot.east) ++(0,3) -- node[rotulo, above, sloped] {alta confiança} (reg.west);
  \draw[seta] (mot.east) ++(0,-3) -- node[rotulo, pos=0.55, below left=0.5mm] {dúvida} (fila.west);
  \draw[seta] (fila) -- (stew);
  \draw[seta] (stew) -- (reg);
  \draw[seta, draw=cinza] (fila) -- (notif);
  \node[rotulo, anchor=south] at (reg.north) {lido por BI, sistemas e IA};

  % aprendizado
  \draw[seta, dashed, draw=cinza] (stew.south) -- ++(0,-19) -| (mot.south);
  \node[rotulo, anchor=north] at (68,-33) {as decisões do steward viram regras e exemplos para o motor};
\end{tikzpicture}

  \caption{Uma arquitetura mínima para uma aplicação própria de master data e
    governança.}
  \label{fig:aplicacao-propria}
  \fonte{elaborado pelo autor.}
\end{figure}

A \textbf{ingestão} traz os registros das fontes. Sempre que possível, ela deve ler da
camada prata (capítulo~\ref{cap:medalhao}), onde os dados já estão limpos e tipados, em vez
de ir direto aos sistemas de origem. O \textbf{motor de conciliação} compara os registros e
atribui a cada par uma nota de confiança. A partir da nota, cada par segue um de três
caminhos: com confiança alta, vai direto para o registro mestre; com confiança baixa, é
descartado; no meio, vai para a \textbf{fila de revisão}, onde uma pessoa decide. O
\textbf{registro mestre}, com a tabela de correspondência, é o produto final, lido pelo BI,
pelos sistemas e pela IA. As \textbf{notificações} avisam quem precisa decidir ou saber.

A seta tracejada na parte de baixo da figura é a que transforma uma ferramenta em governança:
as decisões do steward voltam para o motor, como regras novas ou exemplos, e a próxima
rodada sugere melhor. É o ciclo da Figura~\ref{fig:governanca-ciclo} implementado em
código.

Para a tecnologia, a recomendação é usar o mínimo que o time já domina. Um banco de dados,
Python, um agendador que rode o processo depois de cada carga e um canal de mensagens que as
pessoas já usem resolvem o começo. Complexidade adicional só quando um problema real pedir.

\section{O motor de conciliação}
\label{sec:motor-conciliacao}

O motor de conciliação funciona melhor em camadas, da comparação mais barata para a mais
cara. Primeiro, identificadores fortes: se dois cadastros têm o mesmo documento de
registro da empresa, ou o mesmo código de produto, são a mesma entidade, sem discussão.
Depois, regras de normalização e similaridade de texto para o que sobrou. Por fim, só para os
casos que as regras não resolvem, uma sugestão da IA, sempre acompanhada da explicação do
porquê.

A base teórica dessa comparação é a conciliação probabilística de registros
\parencite{fellegi1969,christen2012}, e há bibliotecas de código aberto maduras que a
implementam em escala, como o Splink \parencite{linacre2022}, além de bibliotecas rápidas de
similaridade de texto. Para começar, porém, a biblioteca padrão do Python é suficiente. O
trecho abaixo normaliza os nomes (caixa baixa, sem acento, sem pontuação, sem sufixos
societários) e compara cada par pela similaridade dos textos.

\begin{python}
import unicodedata
from difflib import SequenceMatcher
from itertools import combinations

SUFIXOS = {"ltda", "sa", "me", "eireli"}

def normalizar(nome: str) -> str:
    """Minúsculas, sem acento, pontuação ou sufixo."""
    nome = unicodedata.normalize("NFKD", nome)
    nome = nome.encode("ascii", "ignore").decode().lower()
    nome = "".join(c for c in nome if c.isalnum() or c == " ")
    partes = nome.split()
    return " ".join(p for p in partes if p not in SUFIXOS)

def pares_suspeitos(nomes, limiar=0.85):
    """Pares de nomes que parecem a mesma entidade."""
    for a, b in combinations(sorted(set(nomes)), 2):
        na, nb = normalizar(a), normalizar(b)
        nota = SequenceMatcher(None, na, nb).ratio()
        if nota >= limiar:
            yield a, b, round(nota, 2)

clientes = ["AGRO SILVA LTDA", "Agro Silva", "Agrosilva",
            "Cerealista Norte S.A."]
for a, b, nota in pares_suspeitos(clientes):
    print(f"Revisar: {a!r} x {b!r} ({nota})")
\end{python}

\begin{terminal}
Revisar: 'AGRO SILVA LTDA' x 'Agro Silva' (1.0)
Revisar: 'AGRO SILVA LTDA' x 'Agrosilva' (0.95)
Revisar: 'Agro Silva' x 'Agrosilva' (0.95)
\end{terminal}

O resultado encontra as três grafias do mesmo cliente e ignora a empresa diferente. Num
sistema real, cada par deixaria de ser impresso na tela e viraria um item na fila de revisão,
com uma mensagem para o dono do domínio. Duas observações importam mais do que o código. A
primeira: comparar todos os pares entre si cresce rápido demais quando o cadastro tem
milhares de registros, e é por isso que as técnicas de bloqueio, que só comparam registros
com alguma característica em comum, existem \parencite{christen2012}. A segunda: o limiar
de 0,85 é arbitrário. Ele deve ser calibrado com uma amostra revisada pelo negócio e
ajustado ao longo do tempo com as decisões do steward.

\section{Revisão humana e notificações}
\label{sec:revisao-notificacao}

A fila de revisão é o coração da aplicação. É ali que a governança acontece, e por isso ela
precisa ser fácil de usar. Cada item deve mostrar os dois registros lado a lado, a nota de
confiança, o motivo da sugestão e três ações: aprovar, corrigir ou rejeitar. Uma interface
simples basta, e existem bibliotecas em Python que permitem montar telas internas desse tipo
em poucas horas.

As notificações devem ir para quem precisa decidir, pelo canal que a pessoa já usa (e-mail
ou o chat da empresa), com o link direto para a revisão. Um cuidado: mensagem demais vira
ruído e passa a ser ignorada. Agrupar as pendências por domínio, num resumo diário ou
semanal, costuma funcionar melhor do que avisar cada caso.

Por fim, a revisão precisa de um prazo combinado e de um indicador de itens parados. Uma
fila que ninguém revisa é pior do que nenhuma fila, porque dá a impressão de que o problema
está sendo cuidado.

\section{Histórico, permissões e privacidade}
\label{sec:auditoria-privacidade}

Toda mudança no registro mestre precisa ficar registrada: o que mudou, quem decidiu,
quando e a partir de qual sugestão. Sem esse histórico, a aplicação é só um banco de dados
editável, e não há governança. Com ele, é possível responder a perguntas que sempre
aparecem, como \enquote{por que este cliente foi unido àquele?}, e desfazer uma decisão
errada sem perder o resto.

As permissões devem ser definidas por domínio: quem vê, quem pode sugerir, quem pode
aprovar. O time comercial aprova clientes; o time de compras aprova fornecedores. E, como
cadastros costumam ter dados pessoais, o que for enviado a um modelo de IA externo precisa
ser limpo antes: nomes de pessoas, documentos e contatos não têm por que sair do ambiente da
empresa para que a IA sugira se duas razões sociais são a mesma (\veja{sec:seguranca}).

\section{Operar como produto}
\label{sec:operar-produto}

A aplicação de master data é um produto de dados como qualquer outro, e as práticas da
Parte~\ref{parte:devops} valem para ela: código versionado, testes, entrega contínua e
monitoramento. Ela também precisa de métricas próprias, que mostram se está funcionando:

\begin{itemize}
  \item porcentagem de pares conciliados automaticamente;
  \item tempo médio entre a sugestão e a decisão;
  \item número de itens parados na fila e há quanto tempo;
  \item duplicidades encontradas por carga, que devem cair com o tempo;
  \item custo por execução, inclusive em tokens, se houver IA no motor
    (capítulo~\ref{cap:custo-consumo}).
\end{itemize}

E precisa de um dono, separado dos donos dos domínios. O dono do domínio de clientes decide
quem é cliente; o dono da aplicação garante que ela funciona. Misturar os dois papéis é um
jeito rápido de ter uma aplicação que ninguém mantém.

\rascunho{
  Relato opcional: uma aplicação interna desse tipo que eu construí ou acompanhei, o que
  funcionou e o que eu faria diferente.
}

\section{Lista de verificação}
\label{sec:checklist-aplicacao}

A Tabela~\ref{tab:checklist-aplicacao} reúne, como referência rápida, as perguntas que vale
responder antes de começar. Nenhuma delas é técnica no sentido estrito, e isso diz muito
sobre onde estão os riscos de um projeto desses.

\begin{table}[htbp]
  \centering
  \caption{Perguntas para responder antes de construir a própria aplicação.}
  \label{tab:checklist-aplicacao}
  \small
  \begin{tabularx}{\textwidth}{@{}>{\raggedright\arraybackslash}p{0.22\textwidth}
      >{\raggedright\arraybackslash}X@{}}
    \toprule
    Tema & Pergunta \\
    \midrule
    Domínio e dono & Qual domínio vem primeiro, e quem decide as regras dele? \\
    Fontes & De onde vêm os registros, e com que frequência mudam? \\
    Identidade & O que faz dois registros serem o mesmo? Existe um identificador forte? \\
    Limiares & Acima de que nota aprova sozinho? Abaixo de qual descarta? \\
    Revisão & Quem revisa, em quanto tempo e por qual canal? \\
    Notificação & Quem precisa saber de cada mudança, e com que frequência? \\
    Histórico & Toda mudança fica registrada, com quem decidiu e quando? \\
    Privacidade & O que pode ir para um modelo de IA, e o que precisa ser mascarado? \\
    Testes & Unicidade, correspondência e regras de negócio rodam a cada carga? \\
    Custo & Quanto custa cada execução, inclusive em tokens? \\
    Manutenção & Quem cuida da aplicação daqui a dois anos? \\
    Saída & Se um dia migrar para uma ferramenta de mercado, os dados saem fácil? \\
    \bottomrule
  \end{tabularx}
  \fonte{elaborado pelo autor.}
\end{table}

\begin{leituras}
  \leitura{christen2012}{Para aprofundar o motor de conciliação: normalização,
    bloqueio, comparação e avaliação dos pares encontrados.}
  \leitura{linacre2022}{A descrição de uma biblioteca de código aberto de conciliação
    probabilística usada em escala por órgãos públicos.}
  \leitura{loshin2009}{Para os componentes de uma solução completa de master data, útil
    para saber o que se está deixando de fora ao começar pequeno.}
\end{leituras}

% Capitulo: Modelagem para ciência de dados
\chapter{Modelagem para ciência de dados}
\label{cap:modelagem-cd}

O esquema estrela é ótimo para somar, filtrar e comparar. Para treinar um modelo que prevê o
futuro, ele quase nunca serve do jeito que está. Modelos preditivos precisam de outra forma
de dado, com outro grão e outros cuidados, e modelos prescritivos precisam ainda de um tipo
de dado que a análise descritiva nunca guardou: o registro das decisões tomadas. Este
capítulo trata dessas diferenças e do erro mais caro da modelagem para ciência de dados, o
vazamento de informação do futuro para o passado.

\section{Descritiva, preditiva e prescritiva}
\label{sec:tres-niveis}

É comum separar a análise de dados em três níveis de ambição. A descritiva responde ao que
aconteceu. A preditiva responde ao que vai acontecer. A prescritiva responde ao que fazer a
respeito. Uma das formulações mais conhecidas dessa escada, do relatório padrão à
otimização, vem de \textcite{davenport2007}, que associam cada degrau a uma vantagem
competitiva maior para quem chega lá.

Um exemplo de logística ajuda a ver a diferença. A análise descritiva mostra que, no último
trimestre, as entregas de uma certa região atrasaram mais do que as das outras, e
principalmente às segundas-feiras. A preditiva estima, para cada entrega programada para
amanhã, a probabilidade de ela atrasar. A prescritiva usa essas probabilidades, junto com o
custo de cada veículo, a capacidade disponível e as janelas de cada cliente, para recomendar
como montar as rotas de forma a minimizar atrasos e custo ao mesmo tempo.

Cada nível depende do anterior. Não se prevê bem o que não se descreve bem: se o atraso é
registrado de formas diferentes em cada filial, o modelo aprende a diferença entre as
filiais, e não o atraso. E não se prescreve bem sem uma previsão confiável. Por isso a
qualidade da camada descritiva, com tudo o que os capítulos anteriores discutiram, é
pré-requisito para os outros dois.

\rascunho{
  Opcional: ligar com a minha experiência de pesquisa, que segue os três níveis
  (descritivo, preditivo, prescritivo), sem detalhar dados ou empresa.
}

\section{A tabela de variáveis}
\label{sec:tabela-features}

Um modelo preditivo aprende a partir de exemplos. Cada exemplo é uma linha com as
informações disponíveis num certo momento, as variáveis (\emph{features}), e o que aconteceu
depois, o alvo. Para prever se uma entrega vai atrasar, cada linha seria uma entrega, com
variáveis como a região, o dia da semana, a distância, o histórico recente de atrasos do
motorista e da rota, e o alvo: atrasou ou não.

Isso muda o grão. No modelo dimensional, a tabela de fatos tem uma linha por evento (cada
entrega realizada), e as perguntas são respondidas agregando. Na tabela de variáveis, o grão
é \enquote{uma linha por entidade e por momento de previsão}: um cliente por mês, um veículo
por dia, uma entrega no momento em que foi programada. E as variáveis muitas vezes já são
agregações calculadas para aquele momento, como o número de atrasos da rota nos trinta dias
anteriores.

O resultado são tabelas largas, com dezenas ou centenas de colunas, que um modelador
dimensional evitaria. Não há problema nisso: a tabela de variáveis é um produto da camada
ouro feito para um consumidor específico, o modelo, e deve ser desenhada para ele. Na
arquitetura medalhão, ela costuma ser construída a partir da prata, e não do ouro do BI,
porque o ouro já vem agregado e perdeu o detalhe de que o modelo precisa
(\veja{sec:bronze-prata-ouro}).

\section{O futuro não pode vazar para o passado}
\label{sec:vazamento}

O erro mais comum, e mais caro, na construção de tabelas de variáveis é o vazamento
(\emph{data leakage}): usar, para treinar o modelo, uma informação que não estaria disponível
no momento em que a previsão precisa ser feita \parencite{kaufman2012}. O modelo fica
excelente nos testes, porque está, na prática, olhando a resposta. Em produção, onde essa
informação não existe ainda, ele erra.

O vazamento raramente é óbvio. Uma variável como \enquote{número de reclamações do cliente}
parece inocente, mas, se calculada com todos os dados disponíveis hoje, inclui as
reclamações feitas depois da data que se quer prever, inclusive aquelas causadas justamente
pelo atraso que o modelo deveria antecipar. Um campo de status que é atualizado quando a
entrega termina vaza o alvo diretamente. E até a forma de separar treino e teste pode vazar,
se dados do futuro entram no treino e dados do passado no teste \parencite{huyen2022}.

A defesa é a correção no tempo (\emph{point-in-time}): cada variável é calculada só com o
que se sabia na data de referência do exemplo. A consulta abaixo mostra a ideia, calculando
para cada cliente e cada data de referência as compras dos noventa dias anteriores, e nada
depois.

\begin{sql}
-- Variáveis de cada cliente calculadas só com o passado de cada data de referência
SELECT
    r.id_cliente,
    r.data_referencia,
    COUNT(v.id_venda)                AS compras_90d,
    COALESCE(SUM(v.valor_total), 0)  AS valor_90d,
    r.cancelou_no_mes_seguinte       AS alvo
FROM referencias_clientes AS r
LEFT JOIN vendas AS v
       ON v.id_cliente = r.id_cliente
      AND v.data_venda <  r.data_referencia                       -- nada do futuro
      AND v.data_venda >= r.data_referencia - INTERVAL '90' DAY
GROUP BY r.id_cliente, r.data_referencia, r.cancelou_no_mes_seguinte;
\end{sql}

A correção no tempo também vale para os atributos. Se o cliente mudou de segmento no meio do
ano, o modelo precisa saber o segmento que ele tinha na data de referência, e não o atual.
Aqui reaparece uma técnica do capítulo~\ref{cap:dw}: as dimensões que mudam devagar, com
uma linha por versão e datas de validade (\veja{sec:kimball}). O que no BI é um recurso para
análise histórica, na ciência de dados vira requisito.

\begin{destaque}
Se o modelo parece bom demais nos testes, a primeira hipótese não é que ele é ótimo. É que
algum dado do futuro vazou para o treino.
\end{destaque}

\section{Reuso de variáveis e dívida técnica}
\label{sec:feature-store}

Quando cada cientista de dados calcula as próprias variáveis, o mesmo conceito ganha várias
definições. \enquote{Cliente ativo} é calculado de um jeito no modelo de cancelamento e de
outro no de recomendação, e ninguém sabe qual está certo. É o mesmo problema das métricas
divergentes que o data warehouse tentou resolver nos anos 1990, agora do lado da ciência
de dados.

A resposta também é parecida: definir uma vez e reaproveitar. Os repositórios de variáveis
(\emph{feature stores}) guardam as definições, calculam as variáveis com correção no tempo
para o treino e servem os mesmos valores para o modelo em produção, garantindo que ele veja
na operação exatamente o que viu no treino \parencite{huyen2022}. Nem todo time precisa de
uma ferramenta dedicada para isso; uma camada bem organizada de tabelas de variáveis, com
testes e documentação, resolve muitos casos.

Vale lembrar uma observação que ficou famosa entre quem constrói sistemas de aprendizado de
máquina: o código do modelo propriamente dito é uma parte pequena do sistema. Em volta dele
há coleta, preparação, validação e monitoramento de dados, e é nessa parte maior que a dívida
técnica se acumula \parencite{sculley2015}. Para quem vem da engenharia e da análise de
dados, é uma boa notícia. A maior parte do trabalho de colocar um modelo em produção é
trabalho de dados, e as práticas da Parte~\ref{parte:devops} valem integralmente.

\section{Da previsão à decisão}
\label{sec:prescritiva}

A análise prescritiva junta a previsão com o que o negócio pode fazer. Isso exige dados que
as camadas anteriores raramente têm: os custos de cada alternativa, as capacidades (quantos
veículos, quantas pessoas, quanto estoque), as restrições (janelas de entrega, contratos,
regras trabalhistas) e os objetivos, muitas vezes em conflito, como custo e nível de
serviço. Com isso, a recomendação costuma sair de um problema de otimização, e não só de um
modelo estatístico.

Há um dado novo que a prescritiva exige e que quase nenhuma empresa guarda: o registro das
decisões. Para saber se a recomendação funciona, é preciso registrar o que foi recomendado,
o que de fato foi decidido (as pessoas nem sempre seguem a recomendação, e muitas vezes têm
bons motivos) e o que aconteceu depois. Sem esse registro, não há como avaliar o sistema nem
melhorá-lo. Com ele, cada decisão vira um exemplo para a próxima.

Vale também modelar cenários como dados, e não como planilhas soltas. Quando o gestor
pergunta \enquote{e se a demanda crescer 20\%?}, a resposta deveria vir de um conjunto de
premissas versionado, comparável com outros cenários, e não de uma cópia da planilha com um
número alterado.

Por fim, as restrições certas vêm de quem conhece a operação. Um modelo prescritivo que
ignora que um cliente não recebe entregas depois das 16 horas, ou que um motorista não pode
dirigir mais de um certo número de horas seguidas, recomenda o impossível. É mais um lugar
em que o conhecimento de negócio (capítulo~\ref{cap:negocio}) é o que separa o modelo
elegante do modelo útil.

\begin{leituras}
  \leitura{huyen2022}{A visão mais completa e prática de como dados, variáveis e
    modelos se organizam num sistema de aprendizado de máquina em produção.}
  \leitura{kaufman2012}{O artigo de referência sobre vazamento de dados: como ele acontece,
    como detectar e como evitar.}
  \leitura{sculley2015}{Artigo curto que mostra por que o modelo é a menor parte do
    problema e onde a dívida técnica se esconde.}
\end{leituras}

% Capitulo: Dados para treinar IA
\chapter{Dados para treinar IA}
\label{cap:treinar-ia}

Quando o consumidor do dado é um modelo que vai aprender com ele, a modelagem muda de
natureza. Deixa de ser sobre tabelas e métricas e passa a ser sobre exemplos: quais entram,
como são rotulados, como são versionados e documentados, e o que nunca pode entrar. Este
capítulo trata dessas decisões. Ele parte de uma constatação prática: a maioria das empresas
não vai treinar um modelo de linguagem do zero, mas quase todas vão ajustar modelos
existentes ou, no mínimo, precisar avaliar se a IA que usam responde certo. As três
coisas pedem dados organizados com cuidado.

\section{Quando o consumidor é um modelo}
\label{sec:consumidor-modelo}

Há três necessidades de dados diferentes quando se fala em treinar IA, e elas aparecem nas
empresas em ordem crescente de frequência.

A primeira é treinar um modelo do zero. Para modelos de linguagem de uso geral, isso exige
volumes de texto e de processamento que estão fora do alcance de quase todas as empresas.
Para modelos menores e específicos, como um classificador de documentos ou um detector de
anomalias, continua sendo comum, e os cuidados do capítulo~\ref{cap:modelagem-cd} valem
integralmente.

A segunda é ajustar um modelo existente (\emph{fine-tuning}): partir de um modelo pronto e
treiná-lo um pouco mais com exemplos da empresa, para que aprenda um vocabulário, um formato
de resposta ou uma tarefa específica. Aqui, poucas centenas ou milhares de exemplos bons
podem fazer muita diferença, e a qualidade de cada exemplo pesa mais do que a quantidade.

A terceira, e mais comum, é avaliar. Toda empresa que usa IA para responder perguntas,
classificar documentos ou gerar consultas precisa saber se ela acerta, e para isso precisa de
um conjunto de avaliação: perguntas ou tarefas com a resposta certa conhecida. Esse conjunto
é, na prática, um dado de treino que todo time de dados vai precisar montar, mesmo sem
treinar nada. Ele é a forma de saber se uma mudança no modelo, no contexto ou nos dados
melhorou ou piorou as respostas, e o capítulo~\ref{cap:contexto-ia} volta a ele.

\section{Curadoria, rotulagem e versionamento}
\label{sec:curadoria}

Em dados para treino, poucos exemplos bons valem mais do que muitos ruins. Exemplos
duplicados fazem o modelo decorar em vez de generalizar. Exemplos contraditórios, em que a
mesma situação tem respostas diferentes, ensinam o modelo a ser inconsistente. Exemplos fora
do domínio desviam o aprendizado. A curadoria (selecionar, limpar e equilibrar os exemplos)
é, por isso, a parte mais importante do trabalho, e a mais parecida com o trabalho de
qualidade de dados que os capítulos anteriores descreveram.

Muitos exemplos precisam de rótulo: a categoria correta de um documento, a resposta certa de
uma pergunta, se uma consulta gerada está correta ou não. Rotular é um trabalho humano, e
envolve decisões que precisam ser explícitas. Quem rotula? Com que critério escrito? O que
fazer nos casos ambíguos? E como saber se os rotuladores concordam entre si? Para essa
última pergunta há medidas clássicas de concordância, que descontam o acordo que
aconteceria por acaso \parencite{cohen1960}. Se dois especialistas do negócio discordam em
um terço dos casos, o problema não é o modelo: é que a regra não está clara, e nenhum
modelo vai aprender uma regra que as próprias pessoas não aplicam de forma consistente.

O conjunto de dados precisa ser versionado como se versiona código (capítulo~\ref{cap:dados-codigo}).
Cada modelo treinado ou avaliado deve estar associado a uma versão exata do conjunto, para
que seja possível reproduzir o resultado, comparar versões e entender por que um modelo novo
se comporta diferente do anterior.

E a separação entre treino, validação e teste deve acontecer antes de qualquer outra coisa,
e ser respeitada. Se exemplos do teste aparecem no treino, ou se o teste é usado
repetidamente para escolher entre versões, a avaliação deixa de medir o que o modelo faria
com dados novos. É uma forma de vazamento, parente da que o capítulo anterior descreveu
(\veja{sec:vazamento}).

\section{Documentar o conjunto de dados}
\label{sec:datasheets}

Um conjunto de dados para treino sem documentação é um risco: ninguém sabe de onde veio, o
que contém, para que serve e para que não deve ser usado. Em 2018, um grupo de
pesquisadores propôs que todo conjunto de dados viesse acompanhado de uma ficha técnica, nos
moldes das fichas de componentes eletrônicos, respondendo a perguntas sobre motivação,
composição, processo de coleta, pré-processamento, usos recomendados, distribuição e
manutenção \parencite{gebru2021}. Uma proposta parecida, para os próprios modelos, descreve
o desempenho de cada modelo em grupos diferentes e as situações em que ele não deve ser
usado \parencite{mitchell2019}.

Para um time de dados, a ficha não precisa ser longa. Algumas perguntas já evitam a maior
parte dos problemas: por que este conjunto foi criado, de quais fontes e em que período veio,
quem rotulou e com que critério, que limpezas foram feitas, o que ele não cobre, e quem é o
responsável por ele. É o equivalente, para dados de treino, do que a documentação do modelo
dimensional sempre foi para o BI: o que permite que outra pessoa use o dado sem precisar
perguntar a quem o fez.

\section{Privacidade, direitos e viés}
\label{sec:privacidade-vies}

Dado pessoal que entra num treino é difícil de tirar depois, porque o modelo pode ter
aprendido algo a partir dele. Por isso a limpeza precisa acontecer antes. A Lei Geral de
Proteção de Dados distingue o dado anonimizado, que deixa de ser considerado dado pessoal
quando a identificação não é mais possível por meios razoáveis, do dado pseudonimizado, em
que a identificação continua possível com uma informação adicional guardada à parte
\parencite{lgpd2018}. Na prática, isso significa que trocar nomes por códigos não é
suficiente para tratar o dado como não pessoal. \pendente{conferir os artigos da LGPD sobre
anonimização e pseudonimização e citar os números}

Além da privacidade, há os direitos sobre o conteúdo. Documentos internos, textos de
clientes, conteúdo de terceiros e dados comprados têm termos de uso, e nem todos permitem o
uso para treinar ou ajustar modelos. Vale verificar antes, com a área jurídica, e registrar a
verificação na ficha do conjunto.

Por fim, a representatividade. Um conjunto que só tem exemplos de uma região, de um tipo de
produto ou de um perfil de cliente ensina o modelo a acertar nesses casos e a errar nos
outros. Um classificador de solicitações treinado só com dados da matriz pode ir mal nas
filiais, onde o vocabulário e os problemas são outros. A ficha técnica deve dizer o que o
conjunto cobre, e a avaliação deve medir o desempenho separadamente em cada grupo relevante,
e não só na média.

\begin{destaque}
Os modelos de IA são, em grande parte, aquilo com que foram treinados e avaliados. Cuidar
desses dados é cuidar do comportamento do modelo.
\end{destaque}

\rascunho{
  Relato opcional: uma experiência de montar um conjunto de avaliação ou de ajuste para
  uma IA usada no trabalho, e o que deu trabalho (rotulagem, casos ambíguos, privacidade).
}

\begin{leituras}
  \leitura{gebru2021}{A proposta de documentar cada conjunto de dados como se documenta
    um componente eletrônico: origem, composição, usos recomendados e limites.}
  \leitura{mitchell2019}{A proposta equivalente para modelos, com atenção ao desempenho em
    grupos diferentes.}
  \leitura{huyen2022}{Os capítulos sobre dados de treino tratam amostragem, rotulagem e
    desbalanceamento de forma prática.}
\end{leituras}

% Capitulo: Dados como contexto para IA
\chapter{Dados como contexto para IA}
\label{cap:contexto-ia}

O uso mais comum de IA com dados nas empresas não é treinar modelos. É dar contexto: no
momento da pergunta, o assistente recebe as definições, os documentos e os dados de que
precisa para responder. A qualidade da resposta passa a depender de como esse contexto está
organizado, e organizar contexto é, no fundo, modelar e documentar dados para um consumidor
novo. Este capítulo fecha a Parte~\ref{parte:modelagem} tratando desse consumidor: o que ele
precisa, onde ele erra e como preparar os dados para ele.

\section{O modelo só sabe o que recebe}
\label{sec:modelo-so-sabe}

Um modelo de linguagem sabe muito sobre o mundo em geral e nada sobre a sua empresa. Ele não
conhece as tabelas, as métricas, as siglas internas nem os documentos de procedimento. Quando
recebe uma pergunta sobre esses assuntos, responde com o que está no contexto que recebeu
junto com a pergunta. Quando o contexto não basta, ele não costuma dizer que não sabe: completa
a lacuna com o que parece plausível, e responde com a mesma segurança de sempre.

Há duas formas principais de dar contexto. A primeira é buscar, a cada pergunta, os trechos de
documentos mais relevantes e entregá-los ao modelo junto com a pergunta. A técnica se chama
geração aumentada por recuperação, ou RAG, e foi proposta em 2020 como forma de combinar o
conhecimento geral do modelo com uma base de documentos consultada na hora
\parencite{lewis2020}. A segunda é dar ao modelo acesso a dados estruturados (tabelas,
métricas, consultas) junto com as definições de negócio de que ele precisa para usá-los
corretamente. Padrões abertos para conectar assistentes a fontes de dados e ferramentas, como
o Model Context Protocol \parencite{anthropic2024mcp}, tornaram essa segunda forma muito mais
fácil de montar.

Uma tentação comum é resolver tudo colocando mais contexto: a documentação inteira, todas as
tabelas, todo o histórico. Não funciona tão bem quanto parece. Estudos mostraram que modelos de
linguagem usam pior a informação que fica no meio de contextos muito longos do que a que fica
no começo ou no fim \parencite{liu2024}. E contexto grande custa caro, porque cada token
enviado é cobrado (capítulo~\ref{cap:custo-consumo}). O objetivo é dar o contexto certo, não o
maior possível.

\section{Metadados e camada semântica como contexto}
\label{sec:metadados-contexto}

Para responder perguntas sobre dados estruturados, o assistente precisa do que um analista
novo precisaria: a descrição de cada tabela e coluna, a definição de cada métrica, como as
tabelas se relacionam, quais filtros aplicar por padrão e alguns exemplos de perguntas com a
consulta correta. Uma definição de métrica escrita pensando nesse leitor pode ser simples
assim:

\begin{terminal}
metrica: receita_liquida
descricao: Soma do valor dos itens faturados, menos devoluções e descontos.
  Não inclui frete nem impostos. Usa a data de faturamento.
tabela: fato_vendas
grão: uma linha por item vendido
filtro_padrao: status = 'faturado'
dono: Controladoria
perguntas_exemplo:
  - Qual foi a receita líquida de março?
  - Receita líquida por região no último trimestre
\end{terminal}

Nada nesse trecho é novo para quem trabalha com dados. É a documentação que sempre foi
considerada boa prática e que sempre ficava para depois. A diferença é que agora ela tem um
consumidor que a usa a cada pergunta, e a falta dela aparece imediatamente na forma de
respostas erradas. É o mesmo movimento que transformou os testes em requisito
(\veja{sec:tdd}) e o data steward em revisor (\veja{sec:steward-tutor}): a IA tornou barato,
e necessário, o que antes não cabia no tempo de ninguém. E a própria IA pode ajudar a
escrever a primeira versão dessas descrições, desde que alguém que conhece o negócio revise.

Um modelo dimensional bem feito já é meio caminho andado. Grão explícito, nomes claros,
métricas definidas uma vez numa camada semântica (\veja{sec:camada-semantica}): tudo isso
reduz a chance de o assistente somar a coisa errada ou escolher a coluna de data errada. O
exemplo da receita com frete multiplicado (\veja{sec:consulta-errada}) é exatamente o tipo
de erro que um grão bem documentado ajuda a evitar.

\section{Documentos e recuperação}
\label{sec:documentos-recuperacao}

Para responder com base em documentos (procedimentos, contratos, políticas, relatórios), os
documentos precisam ser preparados. Eles são divididos em trechos de tamanho adequado, cada
trecho recebe metadados (de que documento veio, de que data, quem é o dono, quem pode ler) e
é indexado de forma que a busca encontre os trechos pelo significado, e não só pelas
palavras exatas. Essa busca por significado usa representações numéricas dos textos
(\emph{embeddings}), em que textos parecidos ficam próximos.

O raciocínio da arquitetura medalhão serve aqui também (Figura~\ref{fig:medalhao}). O
documento original fica guardado como chegou, como no bronze. A versão limpa e dividida em
trechos, com os metadados, faz o papel da prata. O índice pronto para a consulta do
assistente é o ouro. Separar as camadas permite refazer a divisão ou trocar a forma de
indexar sem voltar a coletar os documentos.

Dois cuidados são decisivos. O primeiro é a permissão: o assistente não pode responder a uma
pessoa com base num documento que ela não poderia ler. As permissões precisam vir da fonte e
ser respeitadas na busca, e não aplicadas depois, na resposta. O segundo é a atualização:
contexto velho gera resposta velha com cara de nova. Se a política mudou e o índice não foi
atualizado, o assistente vai responder com a regra antiga, com toda a segurança. Tudo o que o
capítulo~\ref{cap:frequencia} discute sobre frequência de atualização vale para o contexto da
IA.

\section{Organizar para a IA}
\label{sec:organizar-para-ia}

Juntando o que foi dito, organizar os dados para servirem de contexto à IA se resume a uma
lista curta, e quase toda ela já apareceu neste livro com outros nomes:

\begin{itemize}
  \item um catálogo em que cada fonte tem dono e descrição;
  \item um glossário de negócio, com os termos e siglas internos;
  \item métricas definidas uma vez, com grão, filtros e dono;
  \item master data conciliado, para que \enquote{Agro Silva} e \enquote{AGRO SILVA LTDA}
    sejam reconhecidos como o mesmo cliente (capítulo~\ref{cap:master-data});
  \item exemplos de perguntas com a resposta ou a consulta correta;
  \item permissões herdadas da fonte;
  \item frequência de atualização definida para cada fonte de contexto.
\end{itemize}

E, para saber se tudo isso funciona, um conjunto de avaliação: perguntas com a resposta
certa conhecida, rodadas a cada mudança de dado, de contexto ou de modelo
(\veja{sec:consumidor-modelo}). Sem ele, cada ajuste é feito no escuro, e uma melhoria numa
pergunta pode piorar outras dez sem que ninguém perceba.

\begin{destaque}
Um assistente que responde perguntas sobre os dados da empresa é tão bom quanto a
documentação e a modelagem que estão por trás dele.
\end{destaque}

\rascunho{
  Relato: o que eu precisei organizar para uma IA responder bem perguntas sobre os nossos
  dados (descrições, métricas, exemplos), o que deu errado no começo e o que mais fez
  diferença. Sem nomes.
}

Com isso, a Parte~\ref{parte:modelagem} se fecha. Ela mostrou que o mesmo dado precisa ser
organizado de formas diferentes para usos diferentes, e que os cadastros, a governança e a
documentação são a base comum de todos eles. A próxima parte trata de como construir e operar
tudo isso com a disciplina da engenharia de software.

\begin{leituras}
  \leitura{lewis2020}{O artigo que propôs a geração aumentada por recuperação (RAG),
    base da maior parte dos assistentes que respondem com documentos da empresa.}
  \leitura{liu2024}{O estudo que mostra por que mais contexto nem sempre é melhor: modelos
    usam pior a informação no meio de contextos longos.}
  \leitura{huyen2022}{Para a visão de sistema: como dados, modelo e avaliação se
    encaixam em produção.}
\end{leituras}


\part{Dados com práticas de DevOps}
\label{parte:devops}
% Capitulo: Dados como código
\chapter{Dados como código}
\label{cap:dados-codigo}

Durante muito tempo, o trabalho de dados viveu fora das práticas que a engenharia de
software levou décadas para construir. Consultas ficavam salvas no computador de quem as
escreveu, a lógica de transformação estava escondida dentro da ferramenta de BI, e a versão
\enquote{oficial} de um painel era a que estava publicada, sem histórico de como chegou até
ali. Quando algo dava errado, ninguém sabia o que tinha mudado. Esta parte do livro trata de
trazer para os dados a disciplina da engenharia de software, e começa pela base: tratar
transformações, testes e, sempre que possível, painéis como código.

O movimento tem nome e uma literatura própria. A cultura DevOps nasceu para aproximar
desenvolvimento e operação de software, e se apoia em três ideias: fazer o trabalho fluir
em pequenos lotes do desenvolvimento à produção, ter retorno rápido sobre o que dá errado, e
aprender continuamente com isso \parencite{kim2016}. As três se aplicam aos dados quase sem
adaptação.

\section{Por que versionar}
\label{sec:versionar}

Versionar é guardar cada mudança no código com a informação de quem mudou, quando e por
quê, num sistema que permite ver o histórico, comparar versões e voltar atrás. A ferramenta
padrão para isso é o Git, e há um livro completo e gratuito sobre ela \parencite{chacon2014}.
Para quem vem da análise, o mínimo necessário cabe em poucos comandos:

\begin{terminal}
git switch -c receita-por-cliente         # cria um ramo para a mudança
# edita o modelo e o teste correspondente
git add modelos/ testes/
git commit -m "Receita por cliente com frete no grão do pedido"
git push -u origin receita-por-cliente    # envia e abre o pedido de revisão
\end{terminal}

Esse pequeno ritual muda muita coisa. Quando um número muda de um dia para o outro, o
histórico mostra qual alteração causou a mudança. Quando uma alteração quebra algo, voltar à
versão anterior leva minutos. Quando duas pessoas mexem no mesmo modelo, o sistema mostra o
conflito em vez de uma sobrescrever o trabalho da outra. E quando alguém novo chega ao time,
o histórico conta como cada regra de negócio chegou ao formato atual.

O que deve ser versionado vai além das consultas: os modelos de transformação, os testes, as
definições de métricas, a configuração dos pipelines, a documentação e, quando a ferramenta
permite, os próprios painéis. Muitas ferramentas de BI já conseguem exportar a definição dos
painéis como arquivos de texto, o que permite tratá-los com as mesmas práticas. Onde isso
não é possível, vale ao menos manter a lógica fora do painel, na camada de transformação,
onde ela pode ser versionada e testada.

Com a IA, o versionamento ganha uma importância nova. Quando parte do código é escrita por
um assistente ou por um agente, o histórico é o registro de cada alteração e de quem a
aprovou. É ele que permite responder, semanas depois, por que aquele filtro está ali.

\section{Revisão de código}
\label{sec:revisao-codigo}

Revisão de código é a prática de uma segunda pessoa ler cada alteração antes que ela chegue à
produção. Na forma moderna, acontece pelo pedido de revisão (\emph{pull request}): a mudança
fica num ramo separado, o revisor lê as diferenças, comenta, pede ajustes e só então aprova a
integração ao código principal.

O motivo mais citado para revisar código é encontrar defeitos. Um estudo com desenvolvedores
de uma grande empresa de software mostrou que, embora encontrar defeitos fosse a principal
expectativa, boa parte do valor real da revisão estava em outras coisas: transferência de
conhecimento, consciência do time sobre o que está mudando, e soluções alternativas que o
autor não tinha considerado \parencite{bacchelli2013}. Em times de dados, isso é especialmente
verdadeiro. A revisão é muitas vezes o momento em que alguém que conhece o negócio percebe que
a regra implementada não é a regra que a área usa.

Revisar SQL tem algumas perguntas próprias, que valem mais do que qualquer guia de estilo:

\begin{itemize}
  \item Qual é o grão de cada tabela que entra nesta consulta, e qual é o grão do resultado?
  \item Alguma junção pode multiplicar linhas (\veja{sec:consulta-errada})?
  \item Os filtros fazem sentido para o negócio? Algum registro legítimo está ficando de fora?
  \item Os nomes de tabelas e colunas são claros para quem não escreveu o código?
  \item Há testes que garantem as regras importantes, e eles passam?
\end{itemize}

Com a IA, a revisão vira o principal ponto de controle do processo. Quando o código é gerado
em segundos, a etapa que limita a qualidade deixa de ser a escrita e passa a ser a leitura.
Isso tem duas consequências práticas. A primeira é que o autor de uma alteração gerada por IA
continua sendo o responsável por ela, e deve revisá-la antes de pedir a revisão de outra
pessoa. A segunda é que alterações grandes, geradas de uma vez, são difíceis de revisar bem.
Vale pedir à IA alterações pequenas e focadas, do mesmo jeito que se pediria a um colega.

\antesdepois{A revisão era uma etapa entre muitas. O autor escrevia o código devagar,
entendia cada linha, e o revisor conferia principalmente estilo e regras de negócio.}
{A revisão é o principal controle de qualidade. O código chega rápido e em volume, e o revisor
precisa conferir também se o autor entendeu o que a IA escreveu.}

\section{Desenvolvimento e produção}
\label{sec:ambientes}

Testar em produção é a forma mais rápida de descobrir que algo estava errado, e a mais cara.
Um modelo alterado direto no ambiente que alimenta os painéis da diretoria pode publicar
números errados por horas antes de alguém perceber. A solução, conhecida há décadas na
engenharia de software, é separar ambientes.

O mínimo é ter dois: desenvolvimento, onde cada pessoa trabalha e testa as próprias
alterações, e produção, onde ficam os dados que as pessoas usam para decidir. Em armazéns de
dados na nuvem, isso costuma se traduzir em esquemas ou bancos separados, com a mesma
estrutura e permissões diferentes. Cada desenvolvedor pode ter o seu espaço de trabalho, com
uma cópia ou uma amostra dos dados de produção.

Um terceiro ambiente, de integração, completa o desenho. Nele, cada pedido de revisão é
construído e testado de forma isolada, antes de qualquer integração ao código principal. É o
que permite que o revisor veja não só o código, mas também o resultado: as tabelas geradas,
os testes que passaram e os que falharam. O capítulo~\ref{cap:ci-cd} trata de como
automatizar esse caminho.

Dois cuidados acompanham a separação de ambientes. O primeiro são as credenciais: senhas e
chaves de acesso nunca entram no código versionado, e cada ambiente usa as suas, com o
mínimo de permissão necessário. O segundo são os dados sensíveis: ambientes de
desenvolvimento costumam ter controles mais frouxos, e dados pessoais copiados de produção
para lá precisam ser mascarados (\veja{sec:seguranca}).

\rascunho{
  Relato: como foi introduzir versionamento, revisão e ambientes separados num time que
  não tinha essas práticas (resistências, o que convenceu as pessoas, o que mudou nos
  resultados).
}

\begin{leituras}
  \leitura{kim2016}{O manual da cultura DevOps. Os princípios de fluxo, retorno rápido e
    aprendizado contínuo se aplicam diretamente a times de dados.}
  \leitura{chacon2014}{O livro de referência sobre Git, completo e disponível gratuitamente
    na internet.}
  \leitura{bacchelli2013}{O estudo que mostra que o valor da revisão de código vai muito além
    de encontrar defeitos.}
\end{leituras}

% Capitulo: Testes e qualidade de dados
\chapter{Testes e qualidade de dados}
\label{cap:testes}

Todo time de dados já passou por isto: um número errado chega a um painel, alguém da área
percebe antes do time, e a confiança em todos os outros números cai junto. Testes de dados
existem para inverter essa ordem, fazendo o time descobrir o problema antes do usuário. Este
capítulo trata do que testar, de contratos entre quem produz e quem consome dados, e de uma
mudança que considero das mais importantes trazidas pela IA ao trabalho com dados: escrever
os testes antes dos modelos deixou de ser uma boa intenção e pode virar requisito.

\section{O que testar}
\label{sec:o-que-testar}

Um teste de dados é uma verificação automática que roda a cada carga e falha quando uma regra
é violada. Alguns tipos cobrem a maior parte dos problemas:

\begin{itemize}
  \item \textbf{Unicidade.} A chave de cada tabela é única. Se o modelo deveria ter uma
    linha por cliente e tem duas para alguns, toda soma por cliente está errada.
  \item \textbf{Não nulos.} Campos obrigatórios estão preenchidos. Uma venda sem data ou sem
    cliente não deveria chegar à camada de consumo.
  \item \textbf{Integridade referencial.} Toda chave que aponta para outra tabela encontra
    correspondência. Uma venda de um produto que não existe na dimensão de produtos some dos
    relatórios por categoria sem avisar ninguém.
  \item \textbf{Valores aceitos.} Campos com lista fechada de valores (status, tipo, região)
    só contêm valores conhecidos.
  \item \textbf{Regras de negócio.} As regras que a área usa para conferir os números: a
    receita líquida nunca é maior que a bruta, o saldo de estoque nunca é negativo, o total
    do modelo bate com o total da fonte.
  \item \textbf{Atualização e volume.} O dado chegou no prazo, e chegou na quantidade
    esperada. Uma carga que trouxe metade das linhas de sempre passa em todos os outros
    testes e está errada.
\end{itemize}

Os quatro primeiros tipos são genéricos e fáceis de aplicar a qualquer tabela; muitas
ferramentas de transformação já os trazem prontos. Os dois últimos exigem conhecer o negócio
e o comportamento dos dados, e são os que mais pegam problemas reais. Colunas financeiras
merecem atenção especial, com testes de faixa e de soma, porque é nelas que um erro custa
mais caro e é notado mais rápido.

\rascunho{
  Relato: um erro que chegou à produção e que um teste simples teria pego (qual era o erro,
  quanto tempo levou para ser percebido, quem percebeu, qual teste passou a existir depois).
}

\section{Testes antes do modelo}
\label{sec:tdd}

O desenvolvimento guiado por testes, ou TDD, inverte a ordem habitual de trabalho. Em vez de
escrever o código e depois testá-lo, primeiro se escreve um teste que descreve o
comportamento esperado, ele falha porque o código ainda não existe, e então se escreve o
código mínimo para fazer o teste passar \parencite{beck2003}. Na engenharia de software, a
prática tem mais de duas décadas. Em dados, sempre foi recomendada e quase nunca praticada. O
motivo era simples: escrever testes custava tempo que o prazo não dava, e ficava para depois,
que nunca chegava.

A IA mudou essa conta. A partir de uma regra de negócio descrita em texto, um assistente
escreve em segundos os testes correspondentes, sugere casos de borda que ninguém tinha
pensado e monta dados de exemplo para exercitá-los. O custo de escrever testes caiu tanto
que o TDD em dados pode deixar de ser um diferencial e passar a ser requisito do time.

O fluxo fica assim. Antes de escrever o modelo, escrevem-se os testes que descrevem as regras
de negócio: \enquote{toda venda cancelada sai da receita}, \enquote{o saldo nunca fica
negativo}, \enquote{o total do frete bate com o dos pedidos}. Esses testes falham, porque o
modelo ainda não existe. Depois se escreve o modelo até que todos passem.

A consequência mais interessante não é técnica. Os testes viram a especificação executável da
regra de negócio. Em vez de validar com o stakeholder cada versão do modelo, o que exige a
presença dele a cada iteração, a validação acontece uma vez, sobre os casos de teste:
\enquote{se o pedido for cancelado depois de faturado, ele sai da receita do mês do
cancelamento ou do mês do faturamento?} Respondida a pergunta, ela vira teste. Daí em diante,
o time itera sobre o modelo sem depender da agenda do stakeholder, e uma mudança de regra vira
uma mudança de teste, discutida e aprovada antes de qualquer código.

\begin{destaque}
Com a IA, escrever testes ficou barato. Os testes viram a especificação da regra de negócio,
validada uma vez com o stakeholder, e o time passa a iterar sem depender da presença dele.
\end{destaque}

Há dois papéis nessa divisão de trabalho. O da IA é gerar os testes a partir da regra
descrita, sugerir casos de borda e montar dados de exemplo. O do profissional é conferir se
os testes dizem mesmo o que o negócio quer, o que exige conhecer o negócio
(capítulo~\ref{cap:negocio}). E há um cuidado decisivo: testes gerados depois do modelo, a
partir do próprio modelo, tendem a confirmar o que ele já faz, inclusive os erros. É por isso
que a ordem importa. O teste escrito antes, a partir da regra, é uma verificação; o teste
escrito depois, a partir do código, é uma fotografia.

\rascunho{
  Relato: um projeto em que os testes foram escritos antes do modelo, com a ajuda da IA, e o
  efeito no tempo de validação com o stakeholder.
}

\section{Contratos de dados}
\label{sec:contratos}

Muitos problemas de dados não nascem no time de dados, e sim na fonte. Um sistema muda o
formato de um campo, renomeia uma coluna ou passa a registrar algo de outra forma, e o
pipeline quebra, ou pior, continua rodando com dados errados. Contratos de dados são uma
tentativa de prevenir isso formalizando o que quem produz o dado se compromete a entregar a
quem o consome \parencite{jones2023}.

Um contrato costuma definir três coisas. A estrutura: quais campos existem, de que tipo, quais
são obrigatórios. O significado: o que cada campo representa e quais valores são válidos. E o
nível de serviço: com que frequência o dado é entregue e com que prazo. Mudanças que quebram o
contrato precisam ser combinadas antes, com tempo para quem consome se adaptar.

Na prática, um contrato só funciona se for verificado automaticamente. Um documento que
descreve a estrutura esperada, sem um teste que confira a cada carga se ela foi respeitada,
envelhece como qualquer documentação. Por isso os contratos e os testes andam juntos: o
contrato é o acordo entre as partes, e os testes são a forma de saber se ele está sendo
cumprido.

\section{Um teste em SQL}
\label{sec:teste-sql}

A forma mais simples e portável de escrever um teste de dados é uma consulta que devolve as
linhas que violam a regra. Se a consulta não devolve nada, o teste passou. O primeiro exemplo
verifica a unicidade da chave de um modelo de receita por cliente:

\begin{sql}
-- Teste: cada cliente aparece uma única vez no modelo
SELECT id_cliente, COUNT(*) AS ocorrencias
FROM receita_por_cliente
GROUP BY id_cliente
HAVING COUNT(*) > 1;
\end{sql}

O segundo verifica uma regra de negócio que pegaria o erro do frete multiplicado
(\veja{sec:consulta-errada}): o frete somado no modelo precisa ser igual ao frete registrado
nos pedidos.

\begin{sql}
-- Teste: o frete do modelo bate com o frete dos pedidos
WITH origem AS (
    SELECT SUM(valor_frete) AS frete FROM pedidos
),
modelo AS (
    SELECT SUM(frete) AS frete FROM receita_por_cliente
)
SELECT o.frete AS frete_origem, m.frete AS frete_modelo
FROM origem AS o
CROSS JOIN modelo AS m
WHERE ABS(o.frete - m.frete) > 0.01;
\end{sql}

Escrito antes do modelo, esse teste falha até que o modelo calcule o frete no grão certo. É
exatamente o tipo de regra que um stakeholder da área financeira confirma em um minuto e que
protege o modelo para sempre. As ferramentas de transformação como código permitem declarar
testes assim junto aos modelos e rodá-los a cada alteração, o que leva ao assunto do próximo
capítulo.

\begin{leituras}
  \leitura{beck2003}{O livro que sistematizou o desenvolvimento guiado por testes. Escrito
    para software, mas a lógica de escrever o teste primeiro se aplica aos dados.}
  \leitura{jones2023}{Uma introdução prática aos contratos de dados, com atenção à relação
    entre quem produz e quem consome.}
  \leitura{moses2022}{Para levar os testes a uma visão mais ampla de observabilidade e
    confiabilidade dos dados.}
\end{leituras}

% Capitulo: Integração e entrega contínuas para dados
\chapter{Integração e entrega contínuas para dados}
\label{cap:ci-cd}

Versionar o código e escrever testes só resolve metade do problema se alguém precisar se
lembrar de rodar os testes e publicar as alterações à mão. Mais cedo ou mais tarde, alguém
esquece, ou publica com pressa numa sexta-feira à tarde. A integração e a entrega contínuas
automatizam o caminho do código até a produção, de forma que cada alteração seja testada e
publicada sempre do mesmo jeito. Este capítulo trata desse caminho, do monitoramento depois
que os dados estão em produção e da documentação que se gera sozinha a partir do código.

\section{Integração contínua}
\label{sec:ci}

Integração contínua é a prática de integrar as alterações de todo o time ao código principal
com frequência, cada uma delas verificada automaticamente. Em dados, isso significa que,
quando alguém abre um pedido de revisão, um processo automático constrói os modelos alterados
num ambiente isolado, roda os testes e mostra o resultado no próprio pedido. O revisor passa a
ver não só o código, mas também se ele funciona.

Alguns cuidados tornam essa verificação útil em vez de lenta. O primeiro é construir só o que
mudou e o que depende disso, em vez do projeto inteiro, o que reduz o tempo e o custo de cada
verificação. O segundo é usar um ambiente isolado por pedido de revisão, de forma que duas
alterações em paralelo não interfiram uma na outra (\veja{sec:ambientes}). O terceiro é
incluir na verificação, além dos testes de dados, as regras de estilo e de boas práticas do
time: nomes, formatação, proibições como selecionar todas as colunas de uma tabela em modelos
de consumo. Regras que dependem da boa vontade de cada revisor são esquecidas; regras
verificadas automaticamente, não.

Com alterações escritas cada vez mais com ajuda de IA, a integração contínua ganha um papel de
rede de segurança. Ela não substitui a revisão humana, mas garante que nenhuma alteração chegue
à revisão sem ter passado pelos testes, o que libera o revisor para olhar o que só uma pessoa
consegue avaliar: se a regra está certa para o negócio.

\section{Entrega contínua}
\label{sec:cd}

Entrega contínua é o passo seguinte: depois de revisada e integrada, a alteração vai para a
produção sem nenhum passo manual \parencite{humble2010}. O mesmo processo que construiu e
testou no ambiente de integração constrói e testa em produção, na mesma ordem e com as mesmas
verificações.

A principal vantagem não é a velocidade, e sim a segurança. Publicar à mão é uma sequência de
passos que alguém precisa lembrar, e cada passo esquecido é um risco. Publicar
automaticamente é sempre igual. E, quando publicar fica fácil, as alterações ficam menores e
mais frequentes, o que torna cada uma mais fácil de revisar e, se der errado, mais fácil de
desfazer. Uma grande publicação mensal com trinta alterações é muito mais arriscada do que
trinta publicações pequenas ao longo do mês.

A pesquisa sobre desempenho de times de software reforça essa intuição. Um estudo de vários
anos com milhares de profissionais identificou quatro indicadores que distinguem os times de
melhor desempenho: a frequência com que publicam, o tempo entre uma alteração pronta e a
produção, a proporção de publicações que causam falha e o tempo para se recuperar de uma
falha. Os melhores times eram, ao mesmo tempo, mais rápidos e mais estáveis
\parencite{forsgren2018}. Os quatro indicadores se aplicam a times de dados com poucas
adaptações, e são um bom ponto de partida para medir a maturidade do processo.

\section{Observabilidade}
\label{sec:observabilidade}

Testes verificam regras conhecidas no momento da carga. Observabilidade é a capacidade de
perceber, em produção, quando algo está diferente do normal, inclusive problemas que ninguém
pensou em testar. A ideia vem do monitoramento de sistemas de software e foi adaptada para
dados, com atenção a algumas dimensões: se o dado foi atualizado no prazo, se o volume está
dentro do esperado, se a estrutura mudou, se a distribuição dos valores mudou e de onde cada
dado veio \parencite{moses2022}.

Um conceito emprestado da engenharia de confiabilidade ajuda a transformar isso em
compromisso: o objetivo de nível de serviço \parencite{beyer2016}. Em vez de prometer que os
dados estarão sempre atualizados, o time combina, por exemplo, que o painel de vendas terá
os dados do dia anterior disponíveis até as 8 horas em pelo menos 99\% dos dias úteis. O
compromisso fica mensurável, e o monitoramento passa a avisar quando ele está em risco. O
capítulo~\ref{cap:frequencia} volta a essa ideia ao discutir com que frequência os dados
precisam atualizar.

O alerta precisa chegar a quem pode agir, e na hora certa. Um alerta que dispara para o time
inteiro a cada pequena variação vira ruído e passa a ser ignorado; um que só aparece num
painel que ninguém abre não serve para nada. O objetivo é que o time descubra o problema antes
do usuário, e que o usuário, quando o problema for inevitável, seja avisado pelo time, e não
descubra sozinho.

\section{Documentação viva}
\label{sec:documentacao}

Documentação escrita à parte do código envelhece no dia em que é publicada. A alternativa é
gerar a documentação a partir do próprio código: as descrições de tabelas e colunas ficam nos
mesmos arquivos que os modelos, os testes declarados viram parte da documentação, e a
linhagem, o mapa de quais tabelas alimentam quais, é desenhada automaticamente a partir das
dependências. Quando o código muda, a documentação muda junto, porque é gerada na mesma
publicação.

Essa documentação tem hoje três leitores. O analista que precisa entender um modelo antes de
usá-lo. O revisor que precisa entender o impacto de uma alteração. E, cada vez mais, o
assistente de IA que responde perguntas sobre os dados (\veja{sec:metadados-contexto}). Para
os três, uma documentação viva, versionada e gerada a partir do código é muito mais confiável
do que um documento à parte.

\antesdepois{Publicar era um evento: alguém rodava os scripts na ordem certa, conferia à mão e
avisava o time. A documentação ficava num documento à parte, quando ficava.}
{Publicar é rotina: cada alteração revisada passa pelos testes e vai para produção sozinha. A
documentação é gerada junto e serve de contexto também para a IA.}

\rascunho{
  Relato: a implantação de um processo de integração e entrega contínuas num time de dados
  (o que foi automatizado primeiro, que problema deixou de acontecer, como a frequência de
  publicação mudou).
}

\begin{leituras}
  \leitura{humble2010}{O livro que definiu a entrega contínua. Os princípios valem para
    pipelines de dados tanto quanto para software.}
  \leitura{forsgren2018}{A pesquisa que mostrou, com dados, que velocidade e estabilidade
    andam juntas, e os quatro indicadores para medir as duas.}
  \leitura{beyer2016}{Para o conceito de objetivo de nível de serviço e a cultura de
    confiabilidade que vem com ele.}
\end{leituras}

% Capitulo: Do pipeline ao dashboard
\chapter{Do pipeline ao dashboard}
\label{cap:pipeline-dashboard}

Em muitas empresas, há uma parede entre quem constrói os dados e quem constrói os painéis. O
time de engenharia entrega tabelas e considera o trabalho terminado; o time de visualização
recebe as tabelas, aplica as próprias regras dentro da ferramenta de BI e publica. Quando um
número sai errado, cada lado aponta para o outro. Este capítulo trata de como derrubar essa
parede: definir as métricas uma vez, deixar claro quem é dono de quê e tratar o caminho do
dado bruto até o painel como um fluxo só.

\section{Métricas definidas uma vez}
\label{sec:camada-semantica}

O problema dos \enquote{mil dashboards e nenhuma verdade} (\veja{sec:mil-dashboards}) quase
sempre tem a mesma origem: a mesma métrica definida em vários lugares. O faturamento é
calculado num painel com um filtro, em outro com outro, e numa planilha com um terceiro
critério. Cada definição faz sentido para quem a escreveu, e juntas elas destroem a confiança
nos números.

A solução é uma camada semântica: um lugar único onde cada métrica é definida, com a fórmula,
os filtros padrão, o grão em que faz sentido e o dono, e de onde todos os painéis e
ferramentas leem. Quando a definição muda, muda num lugar só, e todos os consumidores passam a
usar a definição nova ao mesmo tempo. Como o capítulo~\ref{cap:nuvem} observou, a ideia não é
nova (\veja{sec:nomes-novos}); o que mudou foi que hoje ela pode viver junto com o código de
transformação, versionada e testada como ele.

A camada semântica ganhou um consumidor que torna a sua existência ainda mais importante: a
IA. Um assistente que responde perguntas sobre os dados precisa de definições claras para não
inventar a própria fórmula de faturamento (\veja{sec:metadados-contexto}). A mesma definição
que garante que dois painéis mostrem o mesmo número garante que o assistente responda com
ele.

\section{Quem é dono de quê}
\label{sec:donos}

A parede entre engenharia e visualização é, em boa parte, um problema de responsabilidade mal
definida. Quando ninguém sabe quem é dono da métrica, todos acham que é o outro. A
Tabela~\ref{tab:donos} propõe uma divisão que costuma funcionar.

\begin{table}[htbp]
  \centering
  \caption{Uma divisão de responsabilidades do dado de origem ao painel.}
  \label{tab:donos}
  \small
  \begin{tabularx}{\textwidth}{@{}>{\raggedright\arraybackslash}p{0.2\textwidth}
      >{\raggedright\arraybackslash}p{0.28\textwidth}>{\raggedright\arraybackslash}X@{}}
    \toprule
    O quê & Dono & Responde por \\
    \midrule
    Dado na origem & Área dona do sistema de origem & Registrar corretamente, avisar mudanças no sistema \\
    \addlinespace
    Plataforma & Engenharia de dados & Ingestão, disponibilidade, desempenho, custo \\
    \addlinespace
    Modelos de dados & Analytics engineering & Transformação correta, testes, atualização no prazo \\
    \addlinespace
    Métrica & Dono de negócio da métrica & A definição: o que entra, o que sai, em que critério \\
    \addlinespace
    Painel & Quem o construiu, com a área usuária & Desenho, uso, manutenção e desligamento \\
    \bottomrule
  \end{tabularx}
  \fonte{elaborado pelo autor.}
\end{table}

Duas linhas da tabela costumam surpreender. A primeira é que o dono da métrica é alguém do
negócio, e não do time de dados. O time de dados implementa a definição e garante que ela seja
calculada corretamente, mas quem decide se a devolução desconta do faturamento é a área que
usa o número para decidir. A segunda é que todo painel tem um dono, inclusive para decidir
quando ele deixa de existir. Painéis sem dono são os que se acumulam, consomem atualização e
ninguém tem coragem de apagar (capítulo~\ref{cap:frequencia}).

\section{Um fluxo só}
\label{sec:fluxo-unico}

Com as métricas centralizadas e os donos definidos, o passo seguinte é tratar o caminho do dado
até o painel como um fluxo único, e não como duas entregas separadas.

Na prática, isso significa algumas coisas. A primeira é manter a lógica fora do painel: regras
de negócio, filtros padrão e cálculos ficam na camada de transformação ou na camada semântica,
onde são versionados e testados; o painel só apresenta. A segunda é conhecer a linhagem até o
fim: quando um modelo vai ser alterado, o time sabe quais painéis dependem dele e pode avaliar
o impacto antes, em vez de descobrir depois, por uma reclamação. A terceira, quando a
ferramenta permite, é versionar e publicar os painéis com o mesmo processo dos modelos
(capítulo~\ref{cap:ci-cd}). E a quarta é organizacional: um quadro de trabalho único para
engenharia e visualização, em que uma demanda de painel e a demanda de dados que ela depende
aparecem juntas, como partes da mesma entrega (Parte~\ref{parte:agil}).

O efeito mais visível dessa integração é na conversa com o negócio. Em vez de o pedido passar
por dois times, com duas filas e duas interpretações, ele é tratado do começo ao fim por um
fluxo só, com um responsável pela entrega. O efeito menos visível, e talvez mais importante, é
na qualidade: quando quem constrói o painel enxerga o modelo, e quem constrói o modelo
enxerga o painel, os erros aparecem antes de chegar à reunião.

\rascunho{
  Relato: a integração entre o time de engenharia e o de visualização num projeto real (como
  era antes, o que mudou na forma de trabalhar, o que melhorou e o que ainda é difícil).
}

\begin{leituras}
  \leitura{kimball2013}{Para a base de tudo isto: métricas e dimensões conformadas, definidas
    uma vez e compartilhadas por todos os processos.}
  \leitura{skelton2019}{Para pensar a divisão de responsabilidades entre times de plataforma
    e times que entregam valor ao negócio.}
\end{leituras}

% Capitulo: Com que frequência os dados precisam atualizar?
\chapter{Com que frequência os dados precisam atualizar?}
\label{cap:frequencia}

Poucas decisões num projeto de dados são tomadas tão no automático quanto a frequência de
atualização. O stakeholder pede \enquote{o mais atualizado possível}, o time configura o
pipeline para rodar de hora em hora porque a ferramenta permite, e ninguém volta ao assunto.
Meses depois, descobre-se que o painel é aberto uma vez por semana, na reunião de
segunda-feira, e que todas as outras atualizações consumiram recurso sem que ninguém olhasse
o resultado.

A regra que este capítulo defende é simples de enunciar: atualizar o mais rápido que as fontes
permitem, mas só o quanto a decisão de negócio precisa. Aplicá-la exige entender como o
negócio decide, conhecer os limites das fontes, dividir a responsabilidade entre os papéis
do time e, principalmente, medir o uso dos próprios produtos de dados, porque a primeira
definição quase sempre está errada.

\section{O ritmo da decisão}
\label{sec:ritmo-decisao}

A pergunta \enquote{com que frequência você quer que o dado atualize?} tem uma resposta
previsível: o mais rápido possível. Ninguém pede dado velho. A pergunta certa é outra:
\enquote{quando e como você toma a decisão com esse número?}. A resposta a essa pergunta define
a frequência necessária.

Decisões operacionais, como despachar uma entrega, repor um item que acabou ou bloquear uma
transação suspeita, acontecem em minutos ou horas, e pedem dados atualizados no mesmo ritmo.
Decisões táticas, como ajustar a escala da semana ou revisar o preço de uma linha de produtos,
convivem bem com dados do dia anterior. Decisões estratégicas, como definir metas, aprovar
orçamento ou planejar uma safra, são tomadas uma vez por mês ou por ano, com dados fechados.

A Figura~\ref{fig:frequencia-consumo} cruza a frequência de atualização com a frequência com
que alguém realmente olha os números. Nos dois quadrantes alinhados, a atualização acompanha o
uso. Nos outros dois, há um problema. No quadrante do desperdício, o dado atualiza muito mais do
que é consultado: recurso gasto sem que ninguém aproveite. No quadrante do risco, muita gente
decide todo dia com um número que só atualiza de vez em quando, e as decisões são tomadas com
dado velho sem que ninguém perceba.

\begin{figure}[htbp]
  \centering
  % Frequencia de atualizacao contra frequencia de consumo (cap. 16).
\begin{tikzpicture}[diagrama, x=1mm, y=1mm]
  % quadrantes (os da diagonal, alinhados, em branco; os desalinhados, em cinza)
  \node[caixa, fill=fundo, minimum width=46mm, minimum height=26mm, text width=40mm] at (25,43)
    {\textbf{Desperdício}\\[1mm]atualiza de hora em hora, alguém olha uma vez por semana};
  \node[caixa, minimum width=46mm, minimum height=26mm, text width=40mm] at (75,43)
    {\textbf{Operacional}\\[1mm]a decisão é de minutos ou horas; atualizar rápido vale o custo};
  \node[caixa, minimum width=46mm, minimum height=26mm, text width=40mm] at (25,15)
    {\textbf{Estratégico}\\[1mm]decisão mensal ou semanal; atualização periódica basta};
  \node[caixa, fill=fundo, minimum width=46mm, minimum height=26mm, text width=40mm] at (75,15)
    {\textbf{Risco}\\[1mm]muita gente decide todo dia com um número velho};

  % eixos
  \draw[seta] (-3,-1) -- (104,-1);
  \draw[seta] (-1,-3) -- (-1,61);
  \node[rotulo, anchor=north] at (50,-2) {com que frequência alguém olha os números};
  \node[rotulo, anchor=south, rotate=90] at (-2,29) {com que frequência os dados atualizam};
  \node[rotulo, anchor=north west] at (0,-2) {menos};
  \node[rotulo, anchor=north east] at (102,-2) {mais};
\end{tikzpicture}

  \caption{Frequência de atualização contra frequência de consumo. Nos quadrantes em
    branco as duas estão alinhadas; nos quadrantes em cinza, não.}
  \label{fig:frequencia-consumo}
  \fonte{elaborado pelo autor.}
\end{figure}

Seja qual for a frequência escolhida, o painel deve mostrar, de forma visível, a data e a hora
da última atualização. O usuário precisa saber a idade do número que está vendo, e um painel
que atualiza todo dia às 7 horas, mas que naquele dia falhou, deveria dizer isso em vez de
mostrar silenciosamente os números de ontem.

\rascunho{
  Relato de abertura, da minha experiência como gerente: um pedido de \enquote{tempo real}
  que, olhando o uso, era uma reunião semanal, e o que foi feito a partir disso.
}

\section{O que as fontes permitem}
\label{sec:limites-fontes}

Se o ritmo da decisão define a frequência necessária, as fontes definem a frequência possível.
Alguns sistemas só permitem extrações em janelas noturnas, para não pesar na operação. Algumas
APIs limitam o número de chamadas por hora. Alguns dados só ficam corretos depois de um
fechamento: o lançamento contábil do dia pode ser ajustado até o fim do mês. E nem todo sistema
oferece captura de mudanças, que permite receber só o que mudou em vez de reler tudo.

Atualizar mais rápido que a fonte muda não traz dado novo, só custo. Se o sistema de origem
consolida as vendas uma vez por dia, rodar o pipeline de hora em hora vai copiar o mesmo dado
vinte e quatro vezes.

Também vale lembrar que o custo e a complexidade não crescem de forma linear com a frequência.
Passar de diário para horário costuma ser simples; passar de horário para contínuo, em
fluxo, exige outra arquitetura, outras ferramentas e outra forma de monitorar
\parencite{reis2022}. E a forma de carregar pesa tanto quanto a frequência: uma carga
incremental, que processa só os registros novos ou alterados, pode tornar viável uma
frequência que seria cara demais com uma carga completa.

\section{O que cada profissional precisa olhar}
\label{sec:papeis-frequencia}

A frequência não é responsabilidade de um papel só. O gerente decide, mas cada papel enxerga
uma parte da conta, e a decisão é melhor quando todas as partes chegam até ele. A
Tabela~\ref{tab:papeis-frequencia} resume o que cada um precisa olhar.

\begin{table}[htbp]
  \centering
  \caption{A frequência de atualização vista por cada papel.}
  \label{tab:papeis-frequencia}
  \small
  \begin{tabularx}{\textwidth}{@{}>{\raggedright\arraybackslash}p{0.27\textwidth}
      >{\raggedright\arraybackslash}X@{}}
    \toprule
    Papel & Com o que precisa se preocupar \\
    \midrule
    Gerente ou líder de dados & Entender a necessidade real do negócio, negociar a
      frequência com o stakeholder, equilibrar custo e valor, acompanhar as métricas de
      uso e revisar a decisão de tempos em tempos. \\
    \addlinespace
    Stakeholder (dono do negócio) & Dizer quando e como decide com o número, e aceitar
      que atualizar mais rápido tem custo. \\
    \addlinespace
    Analista ou desenvolvedor de BI & Desenhar o produto para o ritmo de quem consome,
      mostrar a data da última atualização e não prometer uma frequência que a fonte não
      sustenta. \\
    \addlinespace
    Analytics engineer & Escolher a materialização certa (visão, tabela, incremental)
      para a frequência combinada e testar a atualização do dado. \\
    \addlinespace
    Engenheiro de dados & Conhecer os limites das fontes, o custo de cada execução e
      monitorar atrasos e falhas de carga. \\
    \bottomrule
  \end{tabularx}
  \fonte{elaborado pelo autor.}
\end{table}

O papel do gerente merece um destaque. É ele quem precisa traduzir o pedido do stakeholder em
necessidade real, e quem precisa sustentar a conversa quando a resposta for \enquote{diário é
suficiente}. Isso exige conhecer o negócio bem o bastante para propor a frequência certa, e não
só aceitar a pedida (capítulo~\ref{cap:negocio}). E exige números: a conversa fica muito mais
fácil quando o gerente pode mostrar quantas vezes o painel foi aberto no último mês e quanto
custou cada atualização.

Uma forma prática de formalizar o combinado é o objetivo de nível de serviço
(\veja{sec:observabilidade}): \enquote{os dados do dia anterior estarão disponíveis até as 8
horas em 99\% dos dias úteis}. Ele explicita a frequência e o prazo, e cria uma base objetiva
para monitorar e para rediscutir \parencite{beyer2016}.

\section{Métricas sobre os próprios produtos de dados}
\label{sec:metricas-uso}

Para saber se a frequência combinada está certa, é preciso tratar os painéis, tabelas e
modelos como produtos, com métricas próprias. Algumas das mais úteis:

\begin{itemize}
  \item número de acessos e de usuários distintos por período;
  \item data do último acesso;
  \item número de atualizações por período;
  \item custo de cada atualização e custo total por produto;
  \item custo por acesso, ou atualizações por acesso.
\end{itemize}

A última é a que mais revela desperdício. Um painel que atualiza vinte e quatro vezes por dia
e é aberto duas vezes por semana tem mais de oitenta atualizações por acesso. Quase toda
ferramenta de BI e quase todo armazém de dados guarda, de alguma forma, os registros de quem
acessou o quê e de quanto custou cada execução. A consulta abaixo mostra a ideia, sobre duas
tabelas de metadados hipotéticas.

\begin{sql}
-- Atualizações contra acessos de cada produto de dados nos últimos 30 dias
WITH atualizacoes AS (
    SELECT produto,
           COUNT(*)            AS qtd_atualizacoes,
           SUM(custo_creditos) AS creditos
    FROM meta_atualizacoes
    WHERE executado_em >= CURRENT_DATE - INTERVAL '30' DAY
    GROUP BY produto
),
acessos AS (
    SELECT produto,
           COUNT(*)                AS qtd_acessos,
           COUNT(DISTINCT usuario) AS usuarios
    FROM meta_acessos
    WHERE acessado_em >= CURRENT_DATE - INTERVAL '30' DAY
    GROUP BY produto
)
SELECT a.produto,
       a.qtd_atualizacoes,
       COALESCE(v.qtd_acessos, 0)            AS qtd_acessos,
       COALESCE(v.usuarios, 0)               AS usuarios,
       a.creditos,
       a.creditos / NULLIF(v.qtd_acessos, 0) AS creditos_por_acesso
FROM atualizacoes AS a
LEFT JOIN acessos AS v ON v.produto = a.produto
ORDER BY creditos_por_acesso DESC NULLS FIRST;  -- sem acesso nenhum aparece primeiro
\end{sql}

O resultado é uma lista priorizada. No topo ficam os produtos que consomem recurso e ninguém
abre, candidatos a desligar. Logo abaixo, os que atualizam muito mais do que são consultados,
candidatos a baixar a frequência. A mesma lógica vale para assistentes de IA: tokens gastos por
resposta efetivamente usada são o equivalente, na era dos modelos de linguagem, dos créditos
por acesso.

\section{Revisar a definição}
\label{sec:revisar-frequencia}

A frequência combinada no começo do projeto é uma hipótese. O stakeholder imaginava que
olharia o painel todo dia; a rotina dele mostrou que olha uma vez por semana. Ou o contrário:
um painel pensado para a reunião mensal virou ferramenta diária da operação, e agora atualiza
devagar demais. As métricas de uso dizem qual era a realidade, e a definição precisa ser
revisada com base nelas.

Um ritual simples resolve: uma vez por mês ou por trimestre, revisar os produtos mais caros e
os menos acessados, com a lista da seção anterior na mão. Essa revisão cabe naturalmente na
cadência do Scrum com Kanban (capítulo~\ref{cap:scrumban}), como um item fixo da revisão
periódica.

Baixar a frequência ou desligar um produto também é entrega, e precisa ser tratado como tal.
Isso significa conversar com o dono do produto, mostrar os números de uso, combinar a mudança
e comunicá-la a quem usa. Desligar sem avisar destrói confiança; não desligar nunca destrói
orçamento. Numa época em que quase tudo é cobrado por uso, em créditos de processamento ou em
tokens, o custo de errar nessa definição deixou de ser invisível, e é o assunto do próximo
capítulo.

\begin{destaque}
Atualize o mais rápido que as fontes permitem, mas só o quanto a decisão precisa. E meça o
uso, porque a primeira definição quase sempre está errada.
\end{destaque}

\rascunho{
  Relato: uma revisão de frequência ou de produtos sem uso, e o que ela economizou. Números
  só em ordem de grandeza.
}

\begin{leituras}
  \leitura{beyer2016}{De onde vem a ideia de objetivo de nível de serviço (SLO). Os
    capítulos sobre SLOs se aplicam quase sem adaptação a prazos de atualização de dados.}
  \leitura{moses2022}{Trata a atualização do dado como uma das dimensões da
    observabilidade e mostra como monitorar atrasos antes do usuário reclamar.}
  \leitura{reis2022}{Para a discussão de lote contra fluxo contínuo e do custo de cada
    escolha no ciclo de vida da engenharia de dados.}
\end{leituras}

% Capitulo: De usuários a tokens: custo e consumo
\chapter{De usuários a tokens: custo e consumo}
\label{cap:custo-consumo}

Durante muito tempo, o custo de uma ferramenta de dados era uma linha fixa no orçamento: um
valor por usuário, negociado na renovação do contrato. O uso, alto ou baixo, não mudava a
fatura. Essa lógica está mudando. Os armazéns de dados na nuvem passaram a cobrar pelo
processamento, as ferramentas de BI passaram a oferecer modelos por capacidade ou por
crédito, e os modelos de linguagem cobram por token, a unidade de texto que entra e sai de
cada pergunta. Com isso, gerir o consumo virou uma preocupação de todo gerente de dados. Este
capítulo trata dessa mudança, do que ela exige de quem lidera um time e de como os tokens
entram na conta.

\section{O modelo por usuário}
\label{sec:modelo-usuario}

No modelo tradicional, paga-se uma licença por pessoa com acesso à ferramenta, muitas vezes
em níveis diferentes: quem cria conteúdo paga mais, quem só visualiza paga menos. O custo é
previsível, e é negociado de tempos em tempos, quase sempre por compras e TI.

A alavanca de controle do gerente, nesse modelo, é o acesso: quem precisa de licença e de
qual nível. O desperdício típico também é de acesso: licenças pagas para pessoas que saíram da
área, que nunca usaram a ferramenta ou que precisariam só do nível mais barato. É um
desperdício real, mas limitado e fácil de enxergar numa planilha de usuários.

O que o modelo por usuário não faz é ligar o custo ao uso. Um painel que atualiza de hora em
hora custa o mesmo que um que atualiza uma vez por dia; uma consulta que lê a base inteira custa
o mesmo que uma bem escrita. Para o orçamento, a eficiência do trabalho técnico é invisível.

\section{Créditos, consumo e tokens}
\label{sec:modelo-consumo}

No modelo por consumo, paga-se pelo que foi processado: horas de computação, volume de dados
lido por consulta, atualizações executadas, tokens de entrada e de saída de um modelo de
linguagem. A Tabela~\ref{tab:usuario-consumo} compara os dois modelos.

\begin{table}[htbp]
  \centering
  \caption{Cobrança por usuário e por consumo.}
  \label{tab:usuario-consumo}
  \small
  \begin{tabularx}{\textwidth}{@{}>{\raggedright\arraybackslash}p{0.2\textwidth}
      >{\raggedright\arraybackslash}X>{\raggedright\arraybackslash}X@{}}
    \toprule
    & Por usuário & Por consumo (créditos, tokens) \\
    \midrule
    O que se paga & Uma licença por pessoa com acesso & O que foi processado: consultas, atualizações, tokens \\
    \addlinespace
    Previsibilidade & Alta, custo fixo por período & Baixa, varia com o uso \\
    \addlinespace
    Alavanca de controle & Quem tem acesso e com qual nível & O que roda, com que frequência e com qual modelo \\
    \addlinespace
    Desperdício típico & Licença de quem não usa & Atualização ou consulta que ninguém aproveita \\
    \addlinespace
    Quem acompanha & Compras e TI, na renovação & O gerente do time, de forma contínua \\
    \bottomrule
  \end{tabularx}
  \fonte{elaborado pelo autor.}
\end{table}

A mudança começou com os armazéns de dados na nuvem, que separaram armazenamento de
processamento e passaram a cobrar o processamento pelo uso (\veja{sec:dw-nuvem}). A IA
acelerou a tendência, porque cada pergunta feita a um modelo tem um custo em tokens, e cada
agente que executa uma tarefa faz dezenas ou centenas de chamadas ao modelo. E ferramentas de
BI e de dados que antes cobravam só por usuário passaram a oferecer, ou a impor, modelos por
capacidade ou por crédito. \pendente{conferir e citar exemplos atuais de ferramentas de BI
com cobrança por crédito ou capacidade}

O modelo por consumo tem vantagens reais. Paga-se pelo uso efetivo, um time pequeno começa
barato e cresce conforme precisa, e a eficiência técnica passa a aparecer na fatura, o que dá
ao time um incentivo concreto para fazer bem feito. O risco é a imprevisibilidade: um processo
mal configurado, um agente em laço ou uma consulta pesada agendada por engano podem consumir o
orçamento do mês em poucos dias, e só aparecer na fatura depois.

\antesdepois{O custo da ferramenta era uma linha fixa no orçamento. O gerente controlava
  quem tinha licença, e o uso, alto ou baixo, não mudava a fatura.}
{O custo acompanha o uso. O gerente controla o que roda, com que frequência e com qual
  modelo, e acompanha o consumo como acompanha qualquer outro indicador do time.}

\section{O gerente como gestor de consumo}
\label{sec:gestao-consumo}

No modelo por consumo, o gerente de dados passa a ser, também, gestor de um orçamento
variável. Isso não é exclusivo de dados: a disciplina de gestão de custos de nuvem, conhecida
como FinOps, propõe exatamente isso para qualquer área que consome recursos de nuvem, com três
ideias centrais: dar visibilidade ao custo, tornar cada time responsável pelo próprio consumo
e otimizar de forma contínua, e não em campanhas \parencite{storment2023,finops2026}.

Para um time de dados, algumas práticas traduzem essas ideias para o dia a dia:

\begin{itemize}
  \item \textbf{Orçamento por time e por produto}, com alertas quando o consumo se aproxima do
    limite, e não só depois que a fatura chega.
  \item \textbf{Custo por produto de dados}, a partir das métricas de metadados do capítulo
    anterior (\veja{sec:metricas-uso}): quanto custa manter cada painel, cada pipeline, cada
    assistente.
  \item \textbf{Regras técnicas que evitam desperdício}, verificadas automaticamente na
    integração contínua: não ler todas as colunas de uma tabela em modelos de consumo, não
    fazer junções sem condição, preferir cargas incrementais (\veja{sec:ci}).
  \item \textbf{Ambientes dimensionados para cada uso}, com o processamento de
    desenvolvimento separado e menor que o de produção, e desligado quando ninguém usa.
  \item \textbf{Custo visível para quem pede.} Quando o stakeholder vê quanto custa atualizar
    o painel de hora em hora, a conversa sobre frequência muda de natureza
    (\veja{sec:papeis-frequencia}).
\end{itemize}

A última prática é a que mais muda a relação com o negócio. No modelo por usuário, pedir mais
uma atualização ou mais um painel parecia gratuito. No modelo por consumo, cada pedido tem
um custo que pode ser mostrado, e a decisão passa a ser de valor: vale pagar isso por esse
benefício?

\rascunho{
  Relato: como eu acompanho o consumo hoje (que indicadores olho, com que frequência, que
  alertas existem) e um caso de otimização que reduziu o custo de forma relevante. Sem nomes
  e com valores só em ordem de grandeza.
}

\section{Tokens no trabalho com dados}
\label{sec:tokens-dados}

Os tokens aparecem em vários pontos do trabalho com dados: no assistente que ajuda a escrever
consultas, na pergunta em linguagem natural feita sobre um painel, no agente que executa uma
tarefa inteira, na documentação gerada automaticamente, no motor de conciliação que pede uma
sugestão à IA (\veja{sec:motor-conciliacao}). Em cada caso, o custo depende do tamanho do que
entra (o pedido e o contexto) e do que sai (a resposta), e do modelo usado.

Três fatores costumam fazer o custo crescer mais do que o necessário. O primeiro é o contexto
grande: mandar a tabela inteira, a documentação inteira ou o histórico inteiro a cada
pergunta. Além de caro, nem sempre melhora a resposta (\veja{sec:modelo-so-sabe}). O segundo
são as perguntas repetidas sem reaproveitamento: a mesma pergunta, feita por dez pessoas,
processada dez vezes. O terceiro é usar o modelo mais caro para tudo, inclusive para tarefas
simples que um modelo menor faria igualmente bem.

As práticas que ajudam seguem a mesma lógica do resto do capítulo. Pré-calcular o que é
perguntado sempre, em vez de responder do zero a cada vez. Escolher o modelo pelo tamanho da
tarefa. Limitar o contexto ao necessário, com recuperação de trechos relevantes em vez de
documentos inteiros. E medir tokens por resposta útil, do mesmo jeito que se mede créditos por
acesso num painel.

Vale, por fim, a mesma pergunta do capítulo anterior, agora aplicada à IA: não adianta
responder em segundos, com um modelo caro, uma pergunta que poderia ser atendida por um
relatório semanal. A frequência e a forma de entrega certas continuam sendo decisões de
negócio, e são elas que separam o uso da IA que gera valor do que só gera fatura.

\begin{leituras}
  \leitura{storment2023}{A referência sobre FinOps: como dar visibilidade ao custo de
    nuvem, distribuir a responsabilidade entre os times e otimizar sem travar o trabalho.}
  \leitura{finops2026}{O modelo mantido pela comunidade de FinOps, atualizado com frequência
    e com material específico sobre custos de IA.}
  \leitura{dageville2016}{Para entender, na fonte, por que o custo dos armazéns de dados na
    nuvem passou a acompanhar o processamento.}
\end{leituras}


\part{Ágil em times de dados}
\label{parte:agil}
% Capitulo: A estrutura dos times de dados
\chapter{A estrutura dos times de dados}
\label{cap:estrutura-times}

Antes de falar de como um time de dados trabalha, vale olhar para quem compõe esse time. A
estrutura mais comum nas empresas médias e grandes se formou ao longo das ondas descritas na
Parte~\ref{parte:origens}, e cada onda acrescentou um papel. O resultado é um time grande,
especializado e com muitas passagens de bastão. Este capítulo descreve essa estrutura, mostra
o que a IA comprime nela e o que não comprime, propõe uma versão enxuta possível e termina
com algumas apostas sobre o futuro.

\section{A estrutura tradicional}
\label{sec:estrutura-tradicional}

Numa estrutura típica, há um responsável pela área de dados, abaixo dele um gerente de
engenharia, com um time de engenheiros de dados, e um gerente de analytics, com um time de
analistas de dados e de BI. Muitas vezes há também um product manager, que traduz as
demandas do negócio em prioridades, e alguém dedicado a práticas ágeis, como um scrum master
ou um agile coach. E isso sem contar os times de negócio, que têm os seus próprios analistas e
gestores e são, ao mesmo tempo, clientes e fontes de conhecimento para o time de dados
(parte de cima da Figura~\ref{fig:estrutura-times}).

Essa estrutura surgiu por bons motivos. As ferramentas de cada etapa eram especializadas e
difíceis: extrair e carregar dados exigia uma formação, modelar exigia outra, montar painéis
exigia outra. Ter gente dedicada a cada etapa era a forma de ter qualidade em cada uma.

O custo aparece no caminho de um pedido. Uma pergunta de negócio chega ao product manager,
vira uma história no quadro, vai para a engenharia, que prepara os dados, volta para os
analistas, que montam o painel, e só então chega a quem perguntou. Cada passagem perde um
pouco de contexto e acrescenta um pouco de espera. E, quando o número sai errado, cada etapa
tende a achar que o problema está na outra.

Há também um custo menos visível, conhecido há meio século. Frederick Brooks observou que o
número de canais de comunicação num time cresce muito mais rápido do que o número de pessoas:
com $n$ pessoas, há $n(n-1)/2$ pares que podem precisar conversar \parencite{brooks1975}. Um
time de quatro pessoas tem seis canais; um de doze tem sessenta e seis. Parte do trabalho de
um time grande é coordenar a si mesmo.

E há a observação de Melvin Conway, de 1968: a estrutura dos sistemas que uma organização
constrói tende a copiar a estrutura de comunicação dessa organização \parencite{conway1968}.
Em dados, dois times separados, um de engenharia e um de visualização, costumam produzir duas
camadas que não conversam bem, exatamente como o capítulo~\ref{cap:pipeline-dashboard}
descreveu.

\begin{figure}[htbp]
  \centering
  % Estrutura tradicional de um time de dados e uma versao enxuta possivel (cap. 23).
\begin{tikzpicture}[diagrama, x=1mm, y=1mm,
    titulo/.style={anchor=west, font=\footnotesize\scshape, text=cinza},
    cargo/.style={caixa, minimum width=24mm, minimum height=8mm, font=\scriptsize},
    pilha/.style={cargo, double copy shadow={shadow xshift=1mm, shadow yshift=-1mm, draw=tinta, fill=white}},
    negocio/.style={caixa, draw=cinza, dashed, font=\scriptsize, text=cinza},
    hierarquia/.style={draw=tinta, line width=0.5pt}]

  % --- estrutura tradicional
  \node[titulo] at (0,58) {estrutura tradicional};
  \node[cargo] (head) at (58,50) {Head de dados};
  \node[cargo] (ge) at (14,37) {Gerente de\\engenharia};
  \node[cargo] (ga) at (44,37) {Gerente de\\analytics};
  \node[cargo] (pm) at (74,37) {Product\\manager};
  \node[cargo] (ag) at (103,37) {Agilista\\(scrum master)};
  \foreach \n in {ge,ga,pm,ag} \draw[hierarquia] (head.south) -- ++(0,-3) -| (\n.north);
  \node[pilha] (eng) at (14,23) {Engenheiros\\de dados};
  \node[pilha] (ana) at (44,23) {Analistas de\\dados e BI};
  \draw[hierarquia] (ge) -- (eng);
  \draw[hierarquia] (ga) -- (ana);
  \node[negocio, minimum width=50mm, minimum height=8mm] (neg1) at (88,23)
    {Times de negócio: analistas de\\negócio, gestores, usuários};
  \draw[ligacao, dashed] (pm.south) -- (pm.south |- neg1.north);

  % --- versao enxuta
  \draw[fio] (0,12) -- (118,12);
  \node[titulo] at (0,7) {uma versão enxuta possível};
  \node[destaquecaixa, minimum width=34mm, minimum height=8mm, font=\scriptsize] (lider) at (58,0)
    {Líder de dados\\(produto e fluxo)};
  \node[pilha] (prof) at (40,-14) {Profissionais de\\dados (ponta a ponta)};
  \node[caixa, draw=cinza, minimum width=24mm, minimum height=8mm, font=\scriptsize] (ia) at (76,-14)
    {Agentes de IA\\(revisados pelo time)};
  \draw[hierarquia] (lider.south) -- ++(0,-3) -| (prof.north);
  \draw[ligacao] (lider.south) -- ++(0,-3) -| (ia.north);
  \node[negocio, minimum width=26mm, minimum height=8mm] (neg2) at (105,-14)
    {Negócio, com\\analistas próprios};
  \draw[ligacao, dashed] (ia.east) -- (neg2.west);
\end{tikzpicture}

  \caption{A estrutura tradicional de um time de dados e uma versão enxuta possível com
    a IA. Em linha tracejada, a relação com o negócio.}
  \label{fig:estrutura-times}
  \fonte{elaborado pelo autor.}
\end{figure}

\rascunho{
  Relato: as estruturas de time que eu vi ou montei, quantas passagens de bastão havia entre
  o pedido e a entrega, e o que isso custava. Sem nomes.
}

\section{O que a IA comprime}
\label{sec:ia-comprime}

A IA não elimina etapas do trabalho com dados, mas encurta muitas delas e, principalmente,
diminui a distância entre elas. Com um assistente ajudando nas partes que não domina, a mesma
pessoa consegue ir da ingestão ao painel: um analista que sabe SQL e negócio escreve o modelo
com testes, um engenheiro monta um painel razoável, um analytics engineer cuida do caminho
inteiro (\veja{sec:elt}). A fronteira entre engenharia e análise, que justificava dois times,
fica mais fina.

Alguns papéis também tendem a se juntar. Em times menores, o líder de dados acumula a gestão
de produto e a gestão do fluxo, e as práticas ágeis passam a ser responsabilidade do time,
não de um cargo dedicado. Não porque esses papéis tenham perdido importância, mas porque um
time de quatro pessoas raramente justifica uma pessoa só para cada um deles.

O que a IA não comprime é igualmente importante. Conhecer o negócio continua exigindo tempo e
contato com as pessoas (capítulo~\ref{cap:negocio}). Revisar o que a IA produz continua
exigindo quem saiba ler código e conferir números (capítulo~\ref{cap:programar}). E decidir o
que construir, com que frequência e a que custo continua sendo trabalho de gente
(capítulos~\ref{cap:frequencia} e~\ref{cap:custo-consumo}). Um time enxuto que perde essas
capacidades não ficou eficiente; ficou frágil.

Há dois riscos a evitar. O primeiro é reduzir a estrutura sem reduzir a demanda: um time menor
com o mesmo volume de pedidos e a mesma forma de trabalhar vira um time sobrecarregado, e a
qualidade cai justamente na revisão, que é o que a IA mais exige. O segundo é cortar as
posições de entrada. Se o trabalho de júnior encolher e ninguém mais for contratado para
aprender, daqui a alguns anos faltarão os seniores que hoje revisam o que a IA faz.

\section{Uma estrutura enxuta}
\label{sec:estrutura-enxuta}

A parte de baixo da Figura~\ref{fig:estrutura-times} mostra uma possibilidade. Um líder de
dados, que cuida do produto e do fluxo de trabalho. Alguns profissionais de dados que
trabalham de ponta a ponta, da fonte ao painel, cada um com uma especialidade mais funda, mas
todos capazes de atravessar as camadas. Agentes de IA executando parte das tarefas, sempre
revisados pelo time. E o negócio mais perto dos dados, com analistas próprios que consomem as
bases certificadas e fazem as suas análises (\veja{sec:governanca}).

Duas referências ajudam a pensar essa organização com mais cuidado. \textcite{skelton2019}
propõem um vocabulário para tipos de time: times alinhados a um fluxo de valor, que entregam
ao negócio; times de plataforma, que oferecem serviços internos para os outros; e times de
apoio, que ajudam os demais a adquirir uma capacidade. E propõem formas de interação entre
eles, da colaboração próxima ao consumo de um serviço. \textcite{dehghani2022} vai além e
propõe que os próprios domínios de negócio sejam donos dos seus produtos de dados, com uma
plataforma comum de autoatendimento e uma governança federada.

Nenhuma das duas propostas precisa ser adotada por inteiro. O que vale extrair delas é a
pergunta: o que fica centralizado numa plataforma, com poucas pessoas e muita automação, e o
que vai para perto do negócio? A IA torna viável centralizar menos, porque cada pessoa
consegue fazer mais sozinha, mas não dispensa a plataforma, a governança e as definições
comuns que impedem cada área de criar a sua própria verdade.

O Scrum com Kanban descrito nos próximos capítulos foi pensado para funcionar também num time
assim, de três ou quatro pessoas, sem cargos dedicados a cada prática.

\section{Uma previsão para o futuro}
\label{sec:previsao-times}

Previsões sobre tecnologia envelhecem mal, e esta deve ser lida como aposta. Algumas
tendências, porém, já são visíveis o bastante para arriscar uma direção.

Os times de dados tendem a ficar menores e mais seniores, com a IA fazendo parte do trabalho
que antes era de pessoas em início de carreira: escrever a primeira versão de uma consulta,
documentar tabelas, montar painéis padronizados. O valor de cada profissional passa a estar
mais no julgamento do que no volume do que produz.

O analista de dados e o analista de negócio tendem a convergir. Se a técnica fica mais
acessível, a pessoa que conhece o negócio consegue fazer mais análise sozinha, e a que sabe
análise precisa conhecer mais o negócio para se diferenciar. O ponto de encontro é um
profissional que vive perto da área de negócio e domina o suficiente da técnica para
trabalhar com a IA com critério.

O gerente de dados tende a acumular três gestões: a do produto (o que construir), a do fluxo
(como o trabalho anda) e a do consumo (quanto custa, em créditos e tokens). E, com a IA
proativa do capítulo~\ref{cap:estagios-ia}, uma quarta: a das regras e alçadas que dizem o
que a IA pode decidir sozinha.

E há uma pergunta em aberto, que nenhuma empresa respondeu bem ainda: como se forma um sênior
se o trabalho de júnior encolher? A resposta provavelmente passa por usar a IA como tutor, e
não como substituto, na formação de quem entra (\veja{sec:ia-tutor}), mas ainda está para ser
construída.

\rascunho{
  A minha aposta pessoal, em primeira pessoa: qual destas tendências eu acho mais provável,
  qual eu acho exagerada, e o que eu estou fazendo no meu próprio time por causa disso.
  Alinhar com o epílogo para não repetir.
}

\begin{leituras}
  \leitura{skelton2019}{Um vocabulário claro para tipos de time (de fluxo, de plataforma,
    de apoio) e para as formas de interação entre eles.}
  \leitura{dehghani2022}{A proposta de organizar dados por domínio de negócio, com times
    responsáveis pelos próprios produtos de dados.}
  \leitura{brooks1975}{Um clássico sobre por que times maiores nem sempre entregam mais. A
    observação sobre os canais de comunicação continua atual.}
\end{leituras}

% Capitulo: Por que o ágil de prateleira não encaixa em dados
\chapter{Por que o ágil de prateleira não encaixa em dados}
\label{cap:agil-prateleira}

Muitos times de dados já tentaram adotar o Scrum exatamente como descrito no guia: iterações
de duração fixa, estimativas em pontos, compromisso com um conjunto de histórias no começo de
cada ciclo. E muitos descobriram que algo não encaixa. As histórias mudam de tamanho no meio do
caminho, metade do time é interrompida por demandas urgentes, e a revisão no fim do ciclo
mostra mais trabalho inacabado do que entregue. Este capítulo explica por que isso acontece e
o que vale aproveitar de cada método. Não é uma crítica ao ágil, e sim ao ágil aplicado sem
adaptação.

\section{Descoberta e incerteza}
\label{sec:descoberta}

O ágil nasceu no desenvolvimento de software, e o manifesto que lhe deu nome valoriza as
pessoas, o software funcionando, a colaboração com o cliente e a resposta à mudança
\parencite{agilemanifesto2001}. Nada disso é estranho ao trabalho com dados. O que é diferente
é a natureza da incerteza.

Num projeto de software típico, a maior parte da incerteza está em como construir. Num projeto
de dados, uma parte grande está nos próprios dados, e só aparece quando alguém olha para eles.
Uma história como \enquote{criar o indicador de clientes ativos por região} parece pequena até
que se descobre que a região não está preenchida para um terço dos clientes, que há dois
cadastros de região que não batem e que a área comercial e a financeira definem
\enquote{ativo} de formas diferentes. A tarefa não ficou maior por erro de estimativa; ela
revelou o seu tamanho real quando os dados foram abertos.

A literatura sobre estimativa de software descreve essa incerteza como um cone: no começo de
um projeto, as estimativas podem errar por um fator grande para mais ou para menos, e a margem
só se estreita à medida que o trabalho avança e as incógnitas são resolvidas
\parencite{mcconnell2006}. Em dados, o cone é mais largo e se estreita mais tarde, porque
parte das incógnitas só se resolve olhando os dados. Um método que exige compromisso com o
escopo no começo de cada ciclo, com base em estimativas feitas antes de olhar os dados, está
pedindo ao time o que ele não tem como dar.

A consequência prática é separar descoberta de entrega. Antes de comprometer uma história de
construção, uma investigação curta, com prazo fixo, olha os dados e responde às perguntas que
definem o tamanho do trabalho. Só então a história de construção é refinada e estimada. O
capítulo~\ref{cap:scrumban} volta a essa separação.

\section{Dependências e sustentação}
\label{sec:dependencias}

Duas outras características do trabalho com dados atrapalham o ágil de prateleira. A primeira
são as dependências externas. Um time de dados depende das áreas donas dos sistemas de origem
para liberar acessos, explicar campos, corrigir cadastros e avisar mudanças, e depende das
áreas de negócio para definir regras e validar resultados. Nenhuma dessas pessoas está no time,
e nenhuma segue o ritmo das iterações dele. Uma história pode ficar parada dias esperando uma
resposta, sem que o time possa fazer nada.

A segunda é a sustentação. Um time de dados mantém em produção tudo o que já entregou:
pipelines, modelos, painéis. Quando uma carga falha ou um número sai estranho, a correção não
pode esperar o fim da iteração, porque alguém está decidindo com aquele número agora. Esse
trabalho não planejado chega todos os dias, em quantidade variável, e consome uma parte da
capacidade do time que nenhum planejamento consegue prever com precisão.

Num método que mede o sucesso pela entrega do que foi comprometido no começo do ciclo, as
dependências e a sustentação fazem o time parecer pior do que é. O planejamento fica sempre
otimista, a revisão mostra sempre trabalho inacabado, e o time passa a desconfiar do próprio
processo.

\rascunho{
  Relato: uma experiência com Scrum aplicado ao pé da letra num time de dados, e o que deu
  errado (histórias que mudavam de tamanho, interrupções de sustentação, revisões frustrantes).
}

\section{O que aproveitar de cada método}
\label{sec:aproveitar}

A saída não é abandonar o ágil, e sim combinar o que cada método tem de melhor para o trabalho
com dados. O Scrum e o Kanban, os dois mais usados, têm forças complementares.

Do Scrum, vale aproveitar as cadências \parencite{scrumguide2020}. Um planejamento periódico
obriga o time e o negócio a discutir prioridades com regularidade. A revisão, em que o
trabalho pronto é mostrado a quem pediu, mantém o negócio perto e permite corrigir o rumo cedo.
A retrospectiva, em que o time discute o próprio processo, é a forma de melhorar
continuamente. E a definição de pronto, um critério explícito do que precisa estar feito para
um item ser considerado entregue, evita a entrega que parece terminada e não está.

Do Kanban, vale aproveitar a gestão do fluxo \parencite{anderson2010}. Visualizar todo o
trabalho num quadro, inclusive a sustentação e as dependências, mostra onde as coisas estão
paradas. Limitar o trabalho em andamento força o time a terminar antes de começar, o que reduz
o tempo de cada item no sistema. Tornar as regras explícitas (o que entra em cada coluna,
quando um item está pronto para a próxima etapa) reduz discussões e mal-entendidos. E medir o
fluxo, com indicadores como o tempo que cada item leva do começo ao fim, dá uma base objetiva
para previsões, em vez de estimativas feitas no escuro.

A combinação dos dois, às vezes chamada de Scrumban, foi proposta originalmente como uma forma
de times acostumados ao Scrum evoluírem em direção a um sistema de fluxo \parencite{ladas2009}.
Para dados, ela resolve boa parte dos problemas descritos neste capítulo: as cadências mantêm o
negócio por perto, o fluxo acomoda a descoberta e a sustentação, e os limites de trabalho em
andamento protegem o time de começar tudo e terminar nada. O próximo capítulo descreve como
esse modelo pode ser montado na prática.

\begin{leituras}
  \leitura{anderson2010}{A referência do método Kanban para trabalho intelectual, com as
    práticas de visualizar, limitar o trabalho em andamento e gerir o fluxo.}
  \leitura{scrumguide2020}{O guia oficial do Scrum, curto e direto. Vale reler para separar
    o que o Scrum exige do que foi acrescentado pela prática.}
  \leitura{mcconnell2006}{Para entender por que estimativas no começo de um projeto erram
    tanto, e como lidar com a incerteza em vez de ignorá-la.}
\end{leituras}

% Capitulo: Scrum com Kanban: o modelo que usei
\chapter{Scrum com Kanban: o modelo que usei}
\label{cap:scrumban}

\rascunho{
  Este capítulo foi montado a partir dos princípios do Scrum com Kanban e das ideias do
  livro. Revisar cada seção e substituir as escolhas genéricas pelas que eu de fato usei
  (duração das iterações, colunas do quadro, limites, ritos), marcadas como pendências.
}

O capítulo anterior mostrou por que o Scrum aplicado ao pé da letra costuma frustrar times de
dados, e o que vale aproveitar dele e do Kanban. Este capítulo descreve um modelo que combina
os dois, aplicado em projetos de visualização e de engenharia de dados. Ele foi pensado para
ser copiado e adaptado: nenhuma das escolhas abaixo é a única certa, mas cada uma resolve um
problema concreto, e vale entender o problema antes de mudar a escolha \parencite{ladas2009}.

\section{Cadências}
\label{sec:cadencias}

Do Scrum, o modelo mantém as cadências, porque são elas que mantêm o negócio por perto e o
time melhorando. O que muda é o papel delas: em vez de marcar um compromisso fechado de
escopo, cada cadência é um momento de decisão e de ajuste.

O \textbf{planejamento} acontece no início de cada iteração \pendente{duração das iterações
no meu modelo}. Nele, o time e quem prioriza o trabalho revisam o que está pronto para entrar
e escolhem o que puxar primeiro, respeitando os limites de trabalho em andamento. Não há
compromisso de terminar uma lista fechada; há uma ordem de prioridade e uma previsão baseada
no histórico do fluxo (\veja{sec:metricas-fluxo}).

O \textbf{alinhamento diário} é curto e olha para o quadro, da direita para a esquerda: o que
está perto de ficar pronto e o que está travado, antes do que vai começar. A pergunta central
não é \enquote{o que cada um fez ontem}, e sim \enquote{o que impede o trabalho de fluir}.

A \textbf{revisão} mostra a quem pediu o que ficou pronto, com dados reais, e é o momento de
validar e de corrigir o rumo. É também o lugar natural para a revisão periódica dos produtos
de dados em uso, com as métricas de acesso e de custo do capítulo~\ref{cap:frequencia}.

A \textbf{retrospectiva} discute o processo: o que travou, que regra do quadro precisa mudar,
que tipo de demanda está chegando mal definida. É o rito que faz o modelo evoluir, e o
primeiro a ser cortado quando o time está sobrecarregado, o que é justamente quando ele mais
faz falta.

\rascunho{
  Confirmar quais desses ritos existiam no meu modelo, com que frequência e duração, e se
  havia algum rito próprio (por exemplo, uma reunião de refinamento com o negócio).
}

\section{O quadro e o limite de WIP}
\label{sec:quadro-wip}

Do Kanban, o modelo traz o quadro e a gestão do fluxo \parencite{anderson2010}. A
Figura~\ref{fig:quadro-scrumban} mostra um exemplo de quadro para um time de dados.

\begin{figure}[htbp]
  \centering
  % Exemplo de quadro de Scrum com Kanban para um time de dados (cap. 28).
\begin{tikzpicture}[diagrama, x=1mm, y=1mm,
    cab/.style={font=\scriptsize\bfseries, align=center, text width=14mm},
    raia/.style={font=\tiny, text=cinza, align=left, text width=13mm, anchor=west},
    cartao/.style={draw=tinta, fill=white, minimum width=11mm, minimum height=3.2mm, inner sep=0pt},
    cartaourg/.style={cartao, fill=fundo, line width=0.7pt}]
  % colunas
  \foreach \i/\nome in {0/Entrada, 1/Descoberta, 2/{Pronto p/\\fazer}, 3/{Fazendo\\(até 3)},
                        4/{Revisão\\(até 2)}, 5/{Validação\\negócio}, 6/Pronto} {
    \node[cab] at (21.5+\i*15,-3) {\nome};
  }
  % grade
  \draw[fio] (14,-7.5) -- (119,-7.5);
  \foreach \i in {0,...,7} \draw[fio] (14+\i*15,0) -- (14+\i*15,-49);
  \foreach \y in {-21.5,-35.5,-49} \draw[fio] (0,\y) -- (119,\y);
  % raias
  \node[raia] at (0,-14.5) {Urgente};
  \node[raia] at (0,-28.5) {Projeto};
  \node[raia] at (0,-42.5) {Sustentação};
  % cartoes (exemplo)
  \node[cartaourg] at (66.5,-14.5) {};
  \foreach \c/\y in {0/-25.5, 0/-30.5, 1/-28, 2/-25.5, 2/-30.5, 3/-25.5, 3/-30.5, 4/-28, 5/-28, 6/-25.5, 6/-30.5}
    \node[cartao] at (21.5+\c*15,\y) {};
  \foreach \c/\y in {0/-40, 0/-45, 3/-42.5, 6/-40, 6/-45}
    \node[cartao] at (21.5+\c*15,\y) {};
\end{tikzpicture}

  \caption{Um exemplo de quadro de Scrum com Kanban para um time de dados, com colunas que
    separam descoberta de construção e raias por classe de serviço. O cartão destacado é
    urgente.}
  \label{fig:quadro-scrumban}
  \fonte{elaborado pelo autor.}
\end{figure}

Três escolhas do quadro merecem explicação. A primeira é a coluna de \textbf{descoberta},
separada da construção. É onde ficam as investigações curtas que olham os dados antes de o
trabalho ser refinado (\veja{sec:descoberta}). Um item só passa para \enquote{pronto para
fazer} quando as perguntas que definem o tamanho dele foram respondidas.

A segunda é a coluna de \textbf{validação com o negócio}, separada da revisão técnica. A
revisão confere se o código está certo; a validação confere se o resultado responde à pergunta
de quem pediu. São coisas diferentes, feitas por pessoas diferentes, e misturá-las numa coluna
só esconde onde o trabalho está parado.

A terceira são as \textbf{raias por classe de serviço}. Nem todo trabalho tem a mesma
urgência, e o Kanban propõe tratar cada tipo de forma explícita: itens urgentes, que passam na
frente e têm limite de quantidade; itens de projeto, que seguem a prioridade normal; e itens
de sustentação, que precisam de uma parte reservada da capacidade para não competir com os
projetos a cada dia \parencite{anderson2010}. Visualizar a sustentação no quadro, em vez de
tratá-la como interrupção invisível, é uma das mudanças que mais ajudam um time de dados.

Os números ao lado das colunas \enquote{fazendo} e \enquote{revisão} são os limites de
trabalho em andamento. Eles parecem uma restrição, e são. A razão de existirem vem da teoria
das filas: num sistema estável, o tempo médio que um item leva para atravessar o sistema é
igual à quantidade média de itens em andamento dividida pela vazão
\parencite{little1961,reinertsen2009}. Se o time termina, em média, cinco itens por semana, ter
vinte itens em andamento significa que cada um leva, em média, quatro semanas. Limitar o
trabalho em andamento é a forma mais direta de encurtar o tempo de cada entrega, e de fazer o
time terminar antes de começar.

\rascunho{
  Confirmar as colunas e os limites do meu quadro real, a ferramenta usada e se havia raias
  ou classes de serviço diferentes das do exemplo.
}

\section{Refinamento}
\label{sec:refinamento}

Boa parte dos problemas de um time de dados começa antes do trabalho: pedidos vagos que viram
histórias vagas. O refinamento é o momento de transformar um pedido em algo que o time
consegue construir e o negócio consegue validar, e vale ter critérios explícitos para que uma
história possa entrar na coluna \enquote{pronto para fazer}.

Para histórias de dados, alguns elementos fazem muita diferença. A pergunta de negócio, ou a
decisão que a entrega vai apoiar, escrita de forma que alguém de fora entenda. A forma
tradicional das histórias de usuário ajuda: \enquote{como gerente comercial, quero ver a
margem por vendedor para decidir a revisão das metas} diz mais do que \enquote{painel de
margem} \parencite{cohn2004}. As fontes envolvidas, já investigadas na descoberta. O dono de
negócio de cada métrica. A frequência de atualização necessária, definida a partir do ritmo
da decisão, e não do pedido (\veja{sec:ritmo-decisao}).

E, cada vez mais, os testes de aceite. Com a IA, escrever os testes ficou barato o bastante
para que eles sejam escritos já no refinamento, junto com o stakeholder, a partir das regras
de negócio que ele descreve (\veja{sec:tdd}). A história entra no quadro com a regra
transformada em teste, e o time passa a construir para fazer o teste passar, em vez de
construir e depois descobrir, na validação, que a regra era outra.

\section{Definição de pronto}
\label{sec:definicao-pronto}

A definição de pronto é a lista do que precisa estar feito para um item sair da última coluna.
Ela protege o time da entrega que parece terminada e não está, e protege o negócio do painel
que funciona hoje e quebra na semana que vem. Para dados, uma definição de pronto útil
costuma incluir:

\begin{itemize}
  \item o código revisado por outra pessoa (\veja{sec:revisao-codigo});
  \item os testes escritos antes, e passando (\veja{sec:tdd});
  \item as tabelas, colunas e métricas novas documentadas;
  \item a frequência de atualização combinada e monitorada;
  \item o painel mostrando a data da última atualização;
  \item a publicação feita pelo processo automático, e não à mão (capítulo~\ref{cap:ci-cd});
  \item a validação de quem pediu, com dados reais;
  \item um dono definido para o produto entregue.
\end{itemize}

A lista parece longa, mas a maior parte dela é automática quando as práticas da
Parte~\ref{parte:devops} estão no lugar. E é ela que garante que cada entrega deixe o
ambiente de dados um pouco melhor, e não um pouco mais frágil.

\section{Métricas do fluxo}
\label{sec:metricas-fluxo}

O modelo troca as estimativas detalhadas por medição do fluxo real. Quatro indicadores
bastam para a maior parte das conversas. O \textbf{tempo de ciclo}: quanto tempo um item leva
do momento em que o trabalho começa até ficar pronto. A \textbf{vazão}: quantos itens ficam
prontos por semana. O \textbf{trabalho em andamento}: quantos itens estão sendo trabalhados ao
mesmo tempo. E a \textbf{idade dos itens em andamento}: há quanto tempo cada item está no
quadro, o que revela cedo o que está travado.

Com o histórico desses indicadores, as previsões ficam mais honestas do que qualquer
estimativa em pontos. Em vez de \enquote{essa história vale oito pontos}, o time pode dizer
\enquote{itens desse tipo levam, em 85\% dos casos, até nove dias}. É uma resposta que o
negócio entende e que incorpora, sem esforço adicional, as dependências, a sustentação e a
descoberta que as estimativas costumam ignorar.

\rascunho{
  Relato: como esse modelo foi implantado (o que mudou em relação ao Scrum anterior, como o
  time e o negócio reagiram, o que melhorou nos indicadores). Números só em ordem de grandeza.
}

\begin{leituras}
  \leitura{ladas2009}{O livro que propôs o Scrumban como caminho de evolução do Scrum para um
    sistema de fluxo.}
  \leitura{reinertsen2009}{Para entender, com a teoria das filas, por que limitar o trabalho
    em andamento acelera as entregas.}
  \leitura{cohn2004}{A referência sobre histórias de usuário, útil para escrever histórias de
    dados que o negócio entende.}
\end{leituras}

% Capitulo: Na prática: projetos de visualização
\chapter{Na prática: projetos de visualização}
\label{cap:proj-visualizacao}

Projetos de visualização têm uma armadilha própria: o produto final é tão visível que todo
mundo tem opinião sobre ele, e tão fácil de começar que a tentação é abrir a ferramenta e
sair montando gráficos. Este capítulo mostra como o modelo de Scrum com Kanban funciona num
projeto desse tipo, do pedido vago ao painel publicado, passando pelo que mais costuma dar
errado: construir antes de entender e entregar tudo de uma vez.

\section{Do pedido à pergunta}
\label{sec:pedido-pergunta}

Quase todo projeto de visualização começa com um pedido como \enquote{quero um dashboard de
vendas}. O pedido é legítimo, mas não é uma especificação. Ele não diz quem vai usar, com que
frequência, para decidir o quê, e o que conta como sucesso. Construir a partir dele é
apostar que o resultado vai coincidir com o que a pessoa tinha em mente, e a aposta quase
sempre perde.

O primeiro trabalho, portanto, é transformar o pedido em perguntas. Quem vai abrir o painel?
Que decisão essa pessoa toma, e com que frequência? O que ela faz hoje para tomar essa
decisão, e o que falta? Que número, se mudasse, faria ela agir de forma diferente? As
respostas costumam mostrar que o \enquote{dashboard de vendas} é, na verdade, a necessidade de
acompanhar três indicadores por região toda segunda-feira, com um alerta quando algum deles
sai da faixa normal. É um produto bem menor e bem mais útil do que o painel com trinta gráficos
que seria construído a partir do pedido original.

Essa conversa é também o momento de definir a frequência de atualização a partir do ritmo da
decisão (\veja{sec:ritmo-decisao}), de identificar o dono de cada métrica
(\veja{sec:donos}) e de escrever as primeiras regras que vão virar testes (\veja{sec:tdd}).
No quadro, ela corresponde à coluna de descoberta: nada vai para construção antes que essas
perguntas estejam respondidas.

\section{Protótipo e validação}
\label{sec:prototipo}

Com as perguntas respondidas, o passo seguinte ainda não é construir o painel de verdade. É
fazer um protótipo, o mais barato possível, e mostrá-lo a quem vai usar. Um rascunho em papel,
uma imagem montada numa ferramenta de desenho ou uma versão rápida na própria ferramenta de
BI, com dados fictícios ou com uma amostra, bastam. O objetivo é descobrir o que está errado
quando mudar ainda custa pouco.

A validação do protótipo deve ser feita com a pessoa que vai usar o painel, e não só com quem
pediu, porque nem sempre são a mesma pessoa. E deve ser conduzida com perguntas concretas:
olhando para este painel, que decisão você tomaria hoje? O que você procuraria primeiro? O que
está faltando? Quase sempre aparecem surpresas: o indicador que parecia central é ignorado, e
um detalhe que ninguém tinha mencionado é o que a pessoa procura primeiro.

As referências de visualização do capítulo~\ref{cap:self-service} ajudam a fazer protótipos
melhores desde o começo: uma tela só, poucos elementos, comparações sempre presentes e a
mensagem principal no título (\veja{sec:visualizacao}). A IA também ajuda aqui, gerando
rapidamente variações de layout ou dados fictícios realistas para o protótipo. O que ela não
faz é a conversa com quem vai usar.

\section{Entregas pequenas}
\label{sec:fatiar-dashboard}

O último erro comum é entregar o painel inteiro de uma vez, depois de semanas de construção.
Mesmo com protótipo validado, um painel grande entregue no fim acumula riscos: a fonte que
tinha um problema que ninguém viu, a regra que foi entendida de outro jeito, a mudança de
prioridade do negócio no meio do caminho.

A alternativa é fatiar. A unidade de entrega não é o painel, e sim a pergunta: a primeira
entrega responde à pergunta mais importante, com dados reais, testados e atualizados, mesmo
que o painel tenha só um gráfico. As entregas seguintes acrescentam as outras perguntas, uma de
cada vez, em ordem de valor. Cada fatia passa por todas as colunas do quadro, inclusive a
validação com o negócio, e por todos os itens da definição de pronto
(\veja{sec:definicao-pronto}).

Fatiar tem efeitos que vão além do risco. O usuário começa a usar o painel cedo, e o uso real
mostra o que importa mais do que qualquer conversa. Algumas fatias planejadas deixam de ser
necessárias, porque a primeira já resolveu o problema. E as métricas de uso, desde a primeira
entrega, indicam se o painel está sendo aberto na frequência imaginada
(\veja{sec:metricas-uso}).

\begin{dica}
Se um painel não pode ser dividido em entregas menores que façam sentido sozinhas, é
provável que ele esteja tentando responder perguntas demais. Vale voltar à descoberta.
\end{dica}

\rascunho{
  Relato de um projeto real de visualização, do pedido à entrega, sem nomes de empresa ou
  cliente: o pedido original, o que a descoberta revelou, o protótipo, como foi fatiado e o
  que aconteceu depois da primeira entrega. É o relato central do capítulo.
}

\begin{leituras}
  \leitura{few2006}{Para desenhar painéis que respondem a perguntas de relance, base de bons
    protótipos.}
  \leitura{knaflic2015}{Para conduzir a conversa com quem vai usar e para apresentar as
    entregas de forma que a mensagem fique clara.}
  \leitura{cohn2004}{Para escrever cada fatia como uma história com valor próprio para o
    usuário.}
\end{leituras}

% Capitulo: Na prática: projetos de engenharia de dados
\chapter{Na prática: projetos de engenharia de dados}
\label{cap:proj-engenharia}

Projetos de engenharia de dados têm o problema oposto ao dos projetos de visualização: o
resultado é pouco visível. Um pipeline novo, uma integração com um sistema de origem ou uma
melhoria de desempenho na plataforma raramente aparecem numa reunião de diretoria, e por isso
são difíceis de priorizar, de explicar e de valorizar. Este capítulo mostra como o modelo de
Scrum com Kanban se aplica a esse tipo de trabalho: como escrever histórias quando o usuário
é outro sistema, onde colocar o débito técnico e como conciliar projetos novos com a
sustentação do que já existe.

\section{Histórias quando o usuário é um sistema}
\label{sec:historias-tecnicas}

O formato tradicional das histórias de usuário (\enquote{como tal pessoa, quero tal coisa,
para tal objetivo}) funciona mal quando o trabalho é \enquote{ingerir os dados do novo
sistema de pedidos}. Não há uma pessoa esperando na ponta, e a tentação é escrever histórias
puramente técnicas, como \enquote{criar tabela de pedidos na camada bronze}, que ninguém de
fora do time entende nem consegue priorizar.

Duas práticas ajudam. A primeira é amarrar o trabalho técnico ao valor que ele destrava, mesmo
que esse valor esteja algumas histórias à frente. \enquote{Ingerir os pedidos do novo sistema
para que o painel de vendas continue atualizado depois da migração} diz por que o trabalho
importa e quando ele precisa estar pronto. Quando o usuário é de fato outro sistema, a história
pode ser escrita do ponto de vista dele, com o contrato de dados como critério de aceite
(\veja{sec:contratos}).

A segunda é cuidar do tamanho. Uma boa história é pequena o bastante para atravessar o quadro
em poucos dias e independente o bastante para ser entregue sozinha \parencite{cohn2004}.
Projetos de engenharia tendem a gerar histórias enormes (\enquote{migrar a plataforma}) que
ficam semanas na coluna \enquote{fazendo}. Quebrá-las em fatias que entregam algo verificável,
como uma fonte migrada por vez, com os testes comparando o resultado antigo com o novo, torna
o progresso visível e o risco menor.

Os critérios de aceite de histórias técnicas também podem, e devem, virar testes escritos antes
(\veja{sec:tdd}): o volume carregado bate com a origem, a atualização acontece dentro do
prazo, nenhuma chave fica duplicada. Para engenharia, essa é a forma mais concreta de definir
o que é \enquote{pronto}.

\section{Débito técnico}
\label{sec:debito-tecnico}

A metáfora do débito técnico foi proposta por Ward Cunningham, em 1992, para explicar a
pessoas de negócio por que um código precisava ser revisto: entregar algo rápido, com uma
solução provisória, é como pegar um empréstimo, que acelera hoje e cobra juros depois, até que
seja pago com a reescrita \parencite{cunningham1992}. Em dados, os juros aparecem como
pipelines frágeis que quebram toda semana, modelos que ninguém entende e consultas lentas que
encarecem cada execução.

O débito técnico em si não é o problema. Às vezes faz sentido assumi-lo, de forma consciente,
para cumprir um prazo de negócio. O problema é o débito invisível, que ninguém registrou e que
só aparece na forma de incidentes. A prática mais útil é simples: todo débito assumido vira um
item no quadro, com a descrição do que foi feito de forma provisória e do risco que isso traz.
Assim, ele pode ser priorizado como qualquer outro trabalho.

Para que o débito seja pago de fato, muitos times reservam uma parte fixa da capacidade de cada
iteração para reduzi-lo \pendente{confirmar se eu usava uma porcentagem fixa e qual}. Outra
forma é aproveitar cada história que toca um modelo problemático para melhorá-lo um pouco, em
vez de deixar a melhoria para um projeto de reescrita que nunca é aprovado. E a IA ajuda
aqui: refatorar, documentar e acrescentar testes a código antigo é uma das tarefas em que os
agentes de programação mais economizam tempo, desde que alguém revise.

\section{Projeto e sustentação}
\label{sec:projeto-sustentacao}

Todo time de engenharia de dados divide a capacidade entre construir coisas novas e manter o
que já existe. A sustentação (falhas de carga, números estranhos, pedidos de ajuste, dúvidas
de usuários) chega todos os dias, em quantidade imprevisível, e tem prioridade, porque alguém
está decidindo com aqueles dados agora. Se ela compete diretamente com os projetos, os projetos
atrasam sempre, e o time vive a sensação de nunca terminar nada.

Há duas formas principais de organizar essa convivência. A primeira são duas raias no mesmo
quadro, com uma parte da capacidade reservada para a sustentação
(Figura~\ref{fig:quadro-scrumban}). A segunda é um rodízio: a cada período, uma pessoa do time
fica responsável pela sustentação, como um plantão, e o resto do time fica protegido das
interrupções. Os dois modelos funcionam; a escolha depende do tamanho do time e do volume de
sustentação. Separar em dois times diferentes também é possível em estruturas grandes, mas cria
o risco de que quem constrói nunca sinta o custo de manter o que construiu.

\rascunho{
  Confirmar qual dessas formas eu usava (raias, rodízio, outra) e por quê.
}

Uma situação recorrente merece atenção à parte: os projetos paralelos puxados por um projeto de
negócio, como o de master data do capítulo~\ref{cap:master-data}
(\veja{sec:master-data-puxado}). Eles entram no quadro como um épico próprio, com histórias
amarradas às do projeto principal, e não como tarefas escondidas dentro dele. Isso dá
visibilidade ao trabalho, permite medir quanto ele custou e, principalmente, deixa claro para o
negócio que o domínio entregue vai ser reaproveitado pelos próximos projetos.

Por fim, vale medir a sustentação. Quanto da capacidade do time ela consome, que tipo de
problema mais aparece, quais produtos geram mais chamados. Esses números mostram onde o débito
técnico está cobrando juros mais altos e ajudam a justificar o investimento em pagá-lo.

\rascunho{
  Relato de um projeto real de engenharia de dados, sem nomes de empresa ou cliente: como foi
  quebrado em histórias, que débito foi assumido e como a sustentação foi organizada durante o
  projeto.
}

\begin{leituras}
  \leitura{cohn2004}{Para escrever histórias pequenas e independentes, inclusive quando o
    trabalho é técnico.}
  \leitura{reis2022}{Para a visão completa do ciclo de vida da engenharia de dados, útil para
    planejar projetos de plataforma.}
  \leitura{forsgren2018}{Para medir a sustentação e a recuperação de falhas com os mesmos
    indicadores usados pelos melhores times de software.}
\end{leituras}

% Capitulo: Antes e depois da IA no meu trabalho
\chapter{Antes e depois da IA no meu trabalho}
\label{cap:antes-depois}

\rascunho{
  Capítulo escrito como síntese das mudanças discutidas no livro. Acrescentar, em cada seção,
  um exemplo concreto do meu trabalho (sem nomes e com números só em ordem de grandeza). Se
  preferir um título menos pessoal: \enquote{Antes e depois da IA no trabalho com dados}.
}

Ao longo do livro, a IA apareceu em quase todos os capítulos: mudando o que vale estudar, como
os dados precisam ser organizados, como o código é escrito e revisado, quanto as coisas custam
e como os times se estruturam. Este capítulo junta essas mudanças num lugar só e compara, tema a
tema, o mesmo fluxo de trabalho de um time de dados antes e depois da IA. Não é uma lista de
vantagens. Algumas coisas melhoraram muito, outras pioraram, e quase todas exigiram cuidados
novos.

\section{Estimativa e refinamento}
\label{sec:ad-estimativa}

A parte do trabalho que mais encolheu foi a construção: escrever a consulta, montar o modelo,
gerar a primeira versão do painel. Isso mudou o tamanho relativo das tarefas. O que antes
ocupava a maior parte de uma história agora ocupa uma fração, e o que antes era rápido, como
entender a pergunta de negócio e investigar os dados, passou a dominar o tempo.

\antesdepois{A estimativa era dominada pela construção. Uma história grande era uma história
com muito código para escrever, e o refinamento discutia principalmente como fazer.}
{A construção encolheu e a descoberta passou a dominar. O refinamento discute principalmente o
que fazer, quais são as regras e como saber se está certo.}

A consequência para o refinamento é que ele ficou mais importante, e não menos. Uma história mal
definida, que antes era descoberta no meio da construção, agora é construída rapidamente e da
forma errada. A separação entre descoberta e entrega (\veja{sec:descoberta}) e os critérios de
entrada no quadro (\veja{sec:refinamento}) passaram a ser a principal defesa contra retrabalho.

\rascunho{Exemplo meu: como mudou o tamanho ou a natureza das histórias no meu time.}

\section{Código e revisão}
\label{sec:ad-codigo}

Escrever código ficou muito mais rápido. Revisar ficou mais importante, e não ficou mais fácil.
O volume de código produzido cresceu, e boa parte dele foi escrita por alguém, a IA, que não
estará na reunião quando o número estiver errado (capítulo~\ref{cap:programar}).

\antesdepois{O autor escrevia devagar e entendia cada linha. A revisão conferia principalmente
estilo e regras de negócio.}
{O código chega rápido e em volume. A revisão passou a ser o principal controle de qualidade, e
precisa conferir também se o autor entendeu o que a IA escreveu.}

Na prática, isso levou a algumas regras que antes pareciam excesso de zelo: alterações pequenas,
mesmo quando a IA consegue gerar uma grande; o autor revisando o próprio código gerado antes de
pedir revisão; e as perguntas de grão e de junção feitas em toda revisão de SQL
(\veja{sec:revisao-codigo}).

\rascunho{Exemplo meu: um caso em que a revisão pegou um erro de código gerado por IA.}

\section{Testes}
\label{sec:ad-testes}

Talvez a mudança mais importante tenha sido nos testes. O desenvolvimento guiado por testes
sempre foi recomendado e quase nunca praticado em dados, porque escrever testes custava um tempo
que o prazo não dava. Com a IA, esse custo caiu, e os testes passaram a vir antes dos modelos
(\veja{sec:tdd}).

\antesdepois{Os testes ficavam para quando sobrasse tempo. As regras de negócio eram conferidas
com o stakeholder a cada versão do modelo, o que exigia a presença dele em cada iteração.}
{Os testes vêm antes do modelo e viram a especificação da regra. A validação com o stakeholder
acontece uma vez, sobre os casos de teste, e o time itera sem depender da agenda dele.}

\rascunho{Exemplo meu: um projeto em que escrever os testes antes mudou o tempo de validação.}

\section{Documentação}
\label{sec:ad-documentacao}

A documentação era a eterna pendência de todo time de dados: todos concordavam que era
importante, e ela sempre ficava para depois. Duas coisas mudaram. A IA passou a escrever a
primeira versão das descrições de tabelas, colunas e métricas, a partir do código e dos dados. E a
documentação ganhou um consumidor novo, os assistentes que respondem perguntas sobre os dados, para
os quais ela é insumo essencial (capítulo~\ref{cap:contexto-ia}).

\antesdepois{A documentação ficava num documento à parte, desatualizado, e servia
principalmente para quem chegava ao time.}
{A documentação vive junto com o código, a IA escreve a primeira versão e alguém que conhece o
negócio revisa. Ela serve ao time e à IA, que a usa a cada pergunta.}

O mesmo vale para a transferência de conhecimento. Explicar um modelo antigo a uma pessoa nova
ficou muito mais fácil quando um assistente consegue ler o código e responder às perguntas dela.
O que não ficou mais fácil é transferir o conhecimento de negócio, que continua dependendo de
pessoas.

\section{Custo}
\label{sec:ad-custo}

A forma de pagar pelas ferramentas mudou junto com elas. O que era uma licença fixa por usuário
passou a ser, cada vez mais, consumo: créditos de processamento no armazém de dados e tokens nos
modelos de linguagem (capítulo~\ref{cap:custo-consumo}).

\antesdepois{O custo das ferramentas era fixo e previsível, negociado na renovação. O gerente
controlava o acesso.}
{O custo acompanha o uso. O gerente acompanha o consumo de forma contínua, por time e por
produto, e a eficiência técnica aparece diretamente na fatura.}

\rascunho{Exemplo meu: como passei a acompanhar consumo, e uma otimização que fez diferença.}

\section{O que exigiu novos cuidados}
\label{sec:ad-cuidados}

Nem tudo melhorou. Algumas mudanças exigiram cuidados que antes não existiam, e vale
reuni-los aqui.

O primeiro é a privacidade. Ficou fácil colar num assistente um trecho de dados que nunca
deveria sair da empresa, e as regras sobre o que pode e o que não pode ir para um modelo
precisaram ser escritas e comunicadas (\veja{sec:seguranca}).

O segundo é a confiança excessiva. O código gerado parece correto, a resposta vem com
segurança, e a tendência humana é aceitar. Os estudos citados ao longo do livro mostram que
essa confiança nem sempre é justificada, e que a percepção de ganho de produtividade pode não
corresponder ao ganho real \parencite{becker2025,perry2023}.

O terceiro é a formação de quem está chegando. Quem aprende com a IA fazendo tudo aprende menos
\parencite{shen2026}, e um time que só contrata quem já sabe vai ficar sem seniores no futuro
(\veja{sec:ia-comprime}).

O quarto é o custo invisível do excesso. Ficou tão fácil criar painéis, modelos, aplicações e
agentes que a quantidade cresce mais rápido do que a capacidade de manter. As métricas de uso
dos próprios produtos de dados (\veja{sec:metricas-uso}) e a disciplina de desligar o que não é
usado passaram a ser parte do trabalho.

\begin{destaque}
A IA não diminuiu o trabalho de um time de dados. Ela mudou de lugar: menos construção, mais
definição, revisão, governança e conhecimento de negócio.
\end{destaque}

\rascunho{Exemplo meu: o cuidado novo que mais me custou aprender.}

\begin{leituras}
  \leitura{dellacqua2023}{Para entender por que a IA ajuda muito em algumas tarefas e
    atrapalha em outras, e como identificar a diferença.}
  \leitura{becker2025}{Para a distância entre o ganho percebido e o ganho medido, um bom
    antídoto contra o entusiasmo sem critério.}
  \leitura{forsgren2018}{Para medir, com indicadores objetivos, se a forma de trabalhar do
    time de fato melhorou.}
\end{leituras}


% Epilogo: fora das Partes, sem numero, mas no sumario.
\chapter*{Epílogo}
\addcontentsline{toc}{part}{Epílogo}
\markboth{Epílogo}{Epílogo}
\label{cap:epilogo}

Cada onda contada neste livro prometeu, de alguma forma, acabar com o trabalho de quem analisa
dados. A planilha ia dispensar o centro de processamento. O data warehouse ia acabar com as
extrações manuais. O self-service ia dispensar os analistas, porque cada gestor montaria o
próprio painel. O big data ia encontrar as respostas sozinho, sem que ninguém precisasse fazer
perguntas. E agora a IA vai escrever as consultas, montar os modelos e, no estágio proativo,
tomar as decisões.

Nenhuma dessas promessas se cumpriu do jeito que foi feita, e todas mudaram o trabalho de
verdade. A planilha levou a análise para perto de quem decide. O data warehouse trouxe a ideia de
uma base única e integrada. O self-service encurtou a distância entre a pergunta e o gráfico. A
nuvem tornou acessível o que antes exigia um projeto de anos. Cada onda tirou do caminho uma parte
do trabalho e deixou mais visível a parte que restava.

A IA está fazendo o mesmo, numa escala maior. Ela tira do caminho boa parte da construção, e o
que sobra fica mais evidente do que nunca: entender o negócio, definir as regras, desenhar os
dados, conferir os resultados, decidir o que vale a pena fazer e assumir a responsabilidade pelo
que foi entregue. Não é pouco. É, na verdade, a parte do trabalho que sempre foi a mais valiosa, e
que a técnica, por ser difícil, escondia.

Há algo que atravessa todas as ondas sem mudar, e que talvez seja a principal lição deste livro.
Luhn, em 1958, já falava em levar a informação certa à pessoa certa para orientar uma ação
(\veja{sec:termo-bi}). Ackoff, em 1967, já alertava que mais informação não é melhor informação
(\veja{sec:sad}). As ferramentas mudaram tantas vezes que é difícil lembrar os nomes de todas. A
pergunta de fundo continua a mesma: que decisão este dado ajuda a tomar, e como ter certeza de que
ele está certo?

\rascunho{
  A minha aposta para os próximos anos, em primeira pessoa e assumindo que é aposta. Não repetir
  o que está na previsão sobre os times (\veja{sec:previsao-times}); pode ser sobre o estágio
  proativo da IA, sobre o papel do gerente de dados ou sobre o que eu acho que vai parecer
  ingênuo, daqui a cinco anos, em tudo o que escrevi.
}

Para quem está começando, o conselho que fecha o livro é simples, e não depende de nenhuma
ferramenta. Aprenda a técnica a fundo o bastante para saber quando a máquina errou. Aprenda um
negócio a fundo o bastante para saber o que vale a pena perguntar. Use a IA para ir mais longe, e
não para deixar de aprender. E nunca publique um número que você não saberia explicar.

\rascunho{
  Fecho pessoal opcional: uma ou duas frases minhas, de agradecimento ao leitor ou sobre o que
  me levou a escrever o livro, se não estiver no prefácio.
}


% ---------------------------------------------------------------- fechamento
\appendix
\chapter{Glossário}
\label{cap:glossario}

\rascunho{
  Termos em ordem alfabética, com a definição em uma ou duas frases e, quando
  houver, o termo em inglês entre parênteses. Ir completando enquanto os
  capítulos são escritos.
}

\begin{description}[style=nextline, leftmargin=1.5em, font=\normalfont\scshape]
  \item[ELT] Extração, carga e transformação (\emph{extract, load, transform}). \pendente{definir}
  \item[ETL] Extração, transformação e carga (\emph{extract, transform, load}). \pendente{definir}
  \item[WIP] Trabalho em andamento (\emph{work in progress}). \pendente{definir}
\end{description}


\backmatter
\printbibliography[heading=bibintoc,title={Referências}]
% Ultima pagina do miolo.
\cleardoublepage
\chapter*{Sobre o autor}
\addcontentsline{toc}{chapter}{Sobre o autor}
\markboth{Sobre o autor}{Sobre o autor}
\thispagestyle{empty}

\pendente{escrever a biografia curta, em terceira pessoa, 1 ou 2 parágrafos}


\end{document}
